% Les fêtes religieuses ou nationales — Allemand 1ère A — Cours signé OnBuch+
% Programme officiel MINESEC. Compiler avec : tectonic ALA22-les-fetes-religieuses-ou-nationales.tex
\newcommand{\DOCMATIERE}{Allemand}
\newcommand{\DOCNIVEAU}{1ère A}
\newcommand{\DOCMODULE}{Chapitre 5 : Citoyenneté}
\newcommand{\DOCLECON}{ALA22}
\newcommand{\DOCTITRE}{Les fêtes religieuses ou nationales}
% OnBuch+ — préambule « cours signés » ENRICHI (compilé avec tectonic / XeLaTeX)
% Système visuel « Pop » d'OnBuch+ : fond crème (jamais blanc), bordures encre
% épaisses, ombres « dures » décalées, titres Archivo Black en capitales.
% Version MATHÉMATIQUES : graphes (pgfplots/tikz), boîtes pédagogiques
% (définition, propriété, méthode, attention, le savais-tu, exercice résolu,
% exercice, TP, bilan), page de garde et signature OnBuch+ ancrée en bas.
%
% Macros attendues dans main.tex AVANT \input{preamble} :
%   \DOCMATIERE  \DOCNIVEAU  \DOCMODULE  \DOCLECON (numéro)  \DOCTITRE
\documentclass[11pt]{article}

\usepackage{fontspec}
\usepackage[ngerman,french]{babel} % français = langue principale ; l'allemand via \all{..} / textebox
\usepackage[a4paper,margin=2cm,top=3.4cm,bottom=2.8cm,headheight=1.4cm,footskip=1.3cm]{geometry}
\usepackage{amsmath,amssymb,amsfonts}
\usepackage{mathtools} % \xmapsto, \coloneqq... souvent utilisés par les modèles
\usepackage{cancel} % simplifications barrées (bilans de forces)
% \abs{z} (module, valeur absolue) et \norm{u} (norme), utilisés par les modèles
\providecommand{\abs}{}\let\abs\relax\DeclarePairedDelimiter{\abs}{\lvert}{\rvert}
\providecommand{\norm}{}\let\norm\relax\DeclarePairedDelimiter{\norm}{\lVert}{\rVert}
\usepackage[table]{xcolor} % [table] : fournit aussi \rowcolors (alternance de lignes)
\usepackage{pagecolor}
\usepackage{tikz}
\usetikzlibrary{arrows.meta,positioning,calc,shapes.geometric,shapes.misc,shapes.arrows,decorations.pathmorphing,decorations.pathreplacing,decorations.markings,patterns,shadows,mindmap,trees,fit,backgrounds,matrix,intersections,angles,quotes}
\usepackage{pgfplots}
\usepackage[version=4]{mhchem} % équations nucléaires \ce{^{235}_{92}U}
\usepackage[european,siunitx]{circuitikz} % schémas électriques
\pgfplotsset{compat=1.18}
\usepgfplotslibrary{fillbetween,groupplots} % aires entre courbes (calcul intégral)
\usepackage{enumitem}
\usepackage{fancyhdr}
% [most] charge toutes les bibliothèques tcolorbox (listings, minted, raster,
% documentation...) et gonfle inutilement la mémoire TeX sur un document long
% (~300 boîtes) : on ne charge que ce qui est réellement utilisé.
\usepackage[skins,breakable]{tcolorbox}
\usepackage{graphicx}
\usepackage{colortbl}
\usepackage{booktabs}
\usepackage{tabularx}
\usepackage{diagbox} % case d'en-tête diagonale des tableaux à double entrée
% Colonnes extensibles centrée / gauche / droite (C, L, R), souvent utilisées par les modèles
\newcolumntype{C}{>{\centering\arraybackslash}X}
\newcolumntype{L}{>{\raggedright\arraybackslash}X}
\newcolumntype{R}{>{\raggedleft\arraybackslash}X}
\usepackage{array}
\usepackage{multicol}
\usepackage{multirow}
\usepackage{siunitx}
% parse-numbers=false : accepte aussi \SI{2,5 \times 10^{-3}}{...}
\sisetup{locale=FR,output-decimal-marker={,},group-minimum-digits=4,parse-numbers=false}
\DeclareSIUnit\ohm{\ensuremath{\symup{\Omega}}} % U+2126 absent de la police
\DeclareSIUnit\division{div} % réglages d'oscilloscope (ms/div, V/div)
\DeclareSIUnit\tr{tr} % tours (vitesse de rotation, ex. \SI{1500}{\tr\per\minute})
\DeclareSIUnit\molar{M} % concentration molaire (ex. \SI{3}{\molar}, \SI{1}{\micro\molar})
\DeclareSIUnit\an{an} % année (le modèle confond parfois avec \year, un compteur TeX) — ex. \SI{5}{\kilo\meter\per\an}
\DeclareSIUnit\FCFA{FCFA} % franc CFA (ex. \SI{500}{\FCFA\per\kilo\gram})
\DeclareSIUnit\degreeBrix{\textdegree Brix} % teneur en sucre d'un sirop/jus (réfractométrie)
\DeclareSIUnit\month{mois} % le modèle utilise parfois \month, un compteur TeX, comme unité de durée
\usepackage{qrcode}
% \degree : commande fréquemment utilisée par les modèles pour un angle ou une
% température en dehors de \SI{}{} (non fournie par les paquets chargés).
\providecommand{\degree}{^{\circ}}
% \qed (fin de démonstration), utilisé par les modèles sans amsthm
\providecommand{\qed}{\hfill$\square$}
% \barre : double barre des valeurs interdites dans un tableau de variations (tkz-tab)
\providecommand{\barre}{\ensuremath{\|}}
% environnement proof (amsthm non chargé) utilisé par les modèles
\ifdefined\proof\else\newenvironment{proof}[1][Démonstration]{\par\noindent\textit{#1.}\ }{\qed\par}\fi
\usepackage{lastpage}
\usepackage{hyperref}
\hypersetup{hidelinks}

% ============================================================
% Palette « Pop » OnBuch+ — mêmes hex que la classe P (plus_ui.dart)
% ============================================================
\definecolor{poporange}{HTML}{F59321}
\definecolor{popdark}{HTML}{E07A0C}
\definecolor{popcream}{HTML}{FFF6E8}
\definecolor{popink}{HTML}{1C1714}
\definecolor{poppurple}{HTML}{7A5AE0}
\definecolor{popgreen}{HTML}{2E9E6A}
\definecolor{poppink}{HTML}{E0457B}
\definecolor{popblue}{HTML}{2F7DE1}
\definecolor{popgold}{HTML}{C98A12}
\definecolor{popmuted}{HTML}{8A7F76}
\definecolor{popline}{HTML}{EADFD2}
\definecolor{popgray}{HTML}{8A7F76}
\colorlet{popgrey}{popgray}
% Alias tolérants (le modèle nomme parfois les couleurs différemment)
\colorlet{popwhite}{white}
\colorlet{popblack}{popink}
\colorlet{popred}{poppink}
\colorlet{popyellow}{popgold}
% Teintes claires (fonds de tableaux/encadrés) — le modèle nomme parfois
% les couleurs avec un suffixe L (ex. popblueL) sans les avoir définies.
\colorlet{poporangeL}{poporange!20}
\colorlet{popdarkL}{popdark!20}
\colorlet{poppurpleL}{poppurple!20}
\colorlet{popgreenL}{popgreen!20}
\colorlet{poppinkL}{poppink!20}
\colorlet{popblueL}{popblue!20}
\colorlet{popgoldL}{popgold!20}
\colorlet{popredL}{popred!20}
\colorlet{popgrayL}{popgray!20}
% Teintes claires pour fonds de tableaux / zones
\colorlet{popblueL}{popblue!10!white}
\colorlet{poporangeL}{poporange!14!white}
\colorlet{popgreenL}{popgreen!12!white}
\colorlet{poppurpleL}{poppurple!12!white}
\colorlet{poppinkL}{poppink!12!white}
\colorlet{popgoldL}{popgold!14!white}

\pagecolor{popcream}

% ============================================================
% Typographie
% ============================================================
\defaultfontfeatures{Ligatures=TeX}
\setmainfont{PlusJakartaSans-Medium.ttf}[Path=../../../fonts/,BoldFont=PlusJakartaSans-Bold.ttf]
% Maths en sans-serif (Fira Math) avec chiffres et lettres droites de Plus
% Jakarta, pour que les formules se fondent dans le texte.
\usepackage[math-style=ISO,bold-style=ISO]{unicode-math}
\setmathfont{FiraMath-Regular.otf}[Path=../../../fonts/]
\setmathfont{PlusJakartaSans-Medium.ttf}[Path=../../../fonts/,range={up/num,up/{Latin,latin},\mathup}]
\usepackage{icomma} % virgule décimale sans espace en mode math
% Caractères mathématiques Unicode absents des polices de texte (Plus Jakarta,
% Archivo Black) : rendus par leur équivalent mathématique, en texte comme en
% formule — sinon ils s'affichent en carré vide (ex. « Arithmétique dans ℤ »).
\usepackage{newunicodechar}
\newunicodechar{ℝ}{\ensuremath{\mathbb{R}}}
\newunicodechar{ℤ}{\ensuremath{\mathbb{Z}}}
\newunicodechar{ℂ}{\ensuremath{\mathbb{C}}}
\newunicodechar{ℕ}{\ensuremath{\mathbb{N}}}
\newunicodechar{ℚ}{\ensuremath{\mathbb{Q}}}
\newunicodechar{𝔻}{\ensuremath{\mathbb{D}}}
\newunicodechar{α}{\ensuremath{\alpha}}
\newunicodechar{β}{\ensuremath{\beta}}
\newunicodechar{γ}{\ensuremath{\gamma}}
\newunicodechar{δ}{\ensuremath{\delta}}
\newunicodechar{ε}{\ensuremath{\varepsilon}}
\newunicodechar{θ}{\ensuremath{\theta}}
\newunicodechar{λ}{\ensuremath{\lambda}}
\newunicodechar{μ}{\ensuremath{\mu}}
\newunicodechar{σ}{\ensuremath{\sigma}}
\newunicodechar{τ}{\ensuremath{\tau}}
\newunicodechar{φ}{\ensuremath{\varphi}}
\newunicodechar{ψ}{\ensuremath{\psi}}
\newunicodechar{ω}{\ensuremath{\omega}}
\newunicodechar{Σ}{\ensuremath{\Sigma}}
% majuscules produites par \MakeUppercase dans les titres (α-aminés -> Α)
\newunicodechar{Α}{\ensuremath{\alpha}}
\newunicodechar{Β}{\ensuremath{\beta}}
\newunicodechar{Γ}{\ensuremath{\Gamma}}
\newunicodechar{Δ}{\ensuremath{\Delta}}
\newunicodechar{Ε}{\ensuremath{\varepsilon}}
\newunicodechar{Θ}{\ensuremath{\Theta}}
\newunicodechar{Λ}{\ensuremath{\Lambda}}
\newunicodechar{Μ}{\ensuremath{\mu}}
\newunicodechar{Τ}{\ensuremath{\tau}}
\newunicodechar{Φ}{\ensuremath{\Phi}}
\newunicodechar{Ψ}{\ensuremath{\Psi}}
\newunicodechar{Ω}{\ensuremath{\Omega}}
\newunicodechar{↦}{\ensuremath{\mapsto}}
\newunicodechar{∈}{\ensuremath{\in}}
\newunicodechar{∉}{\ensuremath{\notin}}
\newunicodechar{∧}{\ensuremath{\wedge}}
\newunicodechar{⇔}{\ensuremath{\Leftrightarrow}}
\newunicodechar{⟺}{\ensuremath{\Longleftrightarrow}}
\newunicodechar{⇒}{\ensuremath{\Rightarrow}}
\newunicodechar{⟹}{\ensuremath{\Longrightarrow}}
\newunicodechar{∘}{\ensuremath{\circ}}
\newunicodechar{∪}{\ensuremath{\cup}}
\newunicodechar{∩}{\ensuremath{\cap}}
\newunicodechar{≡}{\ensuremath{\equiv}}
\newunicodechar{⊂}{\ensuremath{\subset}}
\newunicodechar{⊥}{\ensuremath{\perp}}
\newunicodechar{∃}{\ensuremath{\exists}}
\newunicodechar{∀}{\ensuremath{\forall}}
\newunicodechar{∛}{\ensuremath{\sqrt[3]{\,}}}
\newunicodechar{∜}{\ensuremath{\sqrt[4]{\,}}}
\newunicodechar{⃗}{\ensuremath{{}^{\to}}}
\newunicodechar{ⁿ}{\ifmmode^{n}\else\textsuperscript{\ensuremath{n}}\fi}
\newunicodechar{ˣ}{\ifmmode^{x}\else\textsuperscript{\ensuremath{x}}\fi}
\newunicodechar{ᵇ}{\ifmmode^{b}\else\textsuperscript{\ensuremath{b}}\fi}
\newunicodechar{ʳ}{\ifmmode^{r}\else\textsuperscript{\ensuremath{r}}\fi}
\newunicodechar{ᵘ}{\ifmmode^{u}\else\textsuperscript{\ensuremath{u}}\fi}
\newunicodechar{ʸ}{\ifmmode^{y}\else\textsuperscript{\ensuremath{y}}\fi}
\newunicodechar{ᵃ}{\ifmmode^{a}\else\textsuperscript{\ensuremath{a}}\fi}
\newunicodechar{ᵉ}{\ifmmode^{e}\else\textsuperscript{\ensuremath{e}}\fi}
\newunicodechar{ᵏ}{\ifmmode^{k}\else\textsuperscript{\ensuremath{k}}\fi}
\newunicodechar{ᵖ}{\ifmmode^{p}\else\textsuperscript{\ensuremath{p}}\fi}
\newunicodechar{ᴺ}{\ifmmode^{N}\else\textsuperscript{\ensuremath{N}}\fi}
\newunicodechar{ᶜ}{\ifmmode^{c}\else\textsuperscript{\ensuremath{c}}\fi}
\newunicodechar{ʰ}{\ifmmode^{h}\else\textsuperscript{\ensuremath{h}}\fi}
\newunicodechar{ᵠ}{\ifmmode^{q}\else\textsuperscript{\ensuremath{q}}\fi}
\newunicodechar{ₙ}{\ifmmode_{n}\else\textsubscript{\ensuremath{n}}\fi}
\newunicodechar{ₐ}{\ifmmode_{a}\else\textsubscript{\ensuremath{a}}\fi}
\newunicodechar{ᵢ}{\ifmmode_{i}\else\textsubscript{\ensuremath{i}}\fi}
\newunicodechar{ᵣ}{\ifmmode_{r}\else\textsubscript{\ensuremath{r}}\fi}
\newunicodechar{ₖ}{\ifmmode_{k}\else\textsubscript{\ensuremath{k}}\fi}
\newunicodechar{ᵦ}{\ifmmode_{\beta}\else\textsubscript{\ensuremath{\beta}}\fi}
\newunicodechar{ₓ}{\ifmmode_{x}\else\textsubscript{\ensuremath{x}}\fi}
\newfontfamily\popdisplay{ArchivoBlack-Regular.ttf}[Path=../../../fonts/]
\newfontfamily\poplogo{SpaceGrotesk-Bold.ttf}[Path=../../../fonts/]
\newfontfamily\popsemi{PlusJakartaSans-SemiBold.ttf}[Path=../../../fonts/]
\newfontfamily\popxbold{PlusJakartaSans-ExtraBold.ttf}[Path=../../../fonts/]
\newfontfamily\popxboldls{PlusJakartaSans-ExtraBold.ttf}[Path=../../../fonts/,LetterSpace=3.0]

\setlist[enumerate]{leftmargin=1.7em,itemsep=5pt,topsep=4pt}
\setlist[itemize]{leftmargin=1.7em,itemsep=4pt,topsep=4pt,label={\color{poporange}\textbullet}}
\setlength{\parindent}{0pt}
\setlength{\parskip}{5pt}
\renewcommand{\arraystretch}{1.25}

% pgfplots : style Pop par défaut
\pgfplotsset{
  every axis/.append style={axis line style={popink,line width=1pt},tick style={popink},
    grid style={popline,line width=0.5pt},label style={font=\small},tick label style={font=\footnotesize}},
  popcurve/.style={poporange,line width=1.8pt,no markers,smooth},
}
% Même style disponible dans un \draw TikZ simple (hors pgfplots)
\tikzset{popcurve/.style={poporange,line width=1.8pt}}
\tikzset{popgrid/.style={popline,very thin}}
\colorlet{popcurve}{poporange} % le nom du style est parfois utilisé comme couleur

% ============================================================
% Pill / squiggle
% ============================================================
\newcommand{\poppill}[3]{%
  \tikz[baseline=-0.55ex]\node[draw=popink,line width=1.2pt,rounded corners=2.6pt,%
    fill=#2,inner xsep=6.5pt,inner ysep=3pt]{%
    \popxboldls\fontsize{7}{7}\selectfont\color{#3}\MakeUppercase{#1}};%
}
\newcommand{\popsquiggle}[1]{%
  \tikz[baseline=-0.4ex]{
    \draw[#1,line width=2.6pt,line cap=round]
      plot [smooth] coordinates {(0,0.09) (0.34,0.01) (0.68,0.21) (1.02,0.01) (1.36,0.10)};
  }%
}
% Mot-clé surligné (marqueur orange sous le texte)
\newcommand{\cle}[1]{{\popxbold\color{popdark}#1}}

% ============================================================
% En-tête / pied
% ============================================================
\pagestyle{fancy}
\fancyhf{}
\renewcommand{\headrulewidth}{0pt}
\renewcommand{\footrulewidth}{0pt}
\lhead{\raisebox{-0.15ex}{\poplogo\fontsize{13}{13}\selectfont\color{popink}OnBuch\color{poporange}+}}
\rhead{\poppill{\DOCMATIERE\ \textbullet\ \DOCNIVEAU\ \textbullet\ Leçon \DOCLECON}{white}{popink}}
\lfoot{\popxbold\fontsize{7.6}{7.6}\selectfont\color{popmuted}\MakeUppercase{Cours signé OnBuch+}\ \textbullet\ {\popsemi\DOCTITRE}}
\rfoot{\tikz[baseline=-0.6ex]\node[rounded corners=8.5pt,fill=popink,minimum height=17pt,minimum width=17pt,inner xsep=5pt,inner sep=0]{\popxbold\fontsize{8}{8}\selectfont\color{popcream}\thepage\,/\,\pageref*{LastPage}};}

% ============================================================
% Titres
% ============================================================
\newcommand{\chaptitre}[2]{%
  \begin{tcolorbox}[enhanced,colback=poporange,colframe=popink,boxrule=2.6pt,arc=11pt,
      drop shadow=popink,left=15pt,right=15pt,top=12pt,bottom=11pt,before skip=3pt,after skip=18pt]
    \poppill{#2}{popink}{popcream}\par\vspace{9pt}
    {\raggedright\hyphenpenalty=10000\exhyphenpenalty=10000\popdisplay\fontsize{22}{24}\selectfont\color{popcream}\MakeUppercase{#1}\par}
  \end{tcolorbox}
}
% Section numérotée : pastille orange avec le numéro + titre Archivo
\newcounter{popsec}
\newcommand{\coursec}[1]{%
  \stepcounter{popsec}\setcounter{popsub}{0}%
  \par\vspace{16pt}\noindent
  \begin{minipage}{\linewidth}
  \tikz[baseline=-0.7ex]\node[draw=popink,line width=1.6pt,fill=poporange,rounded corners=4pt,minimum size=22pt,inner sep=0]{\popdisplay\fontsize{12}{12}\selectfont\color{popcream}\thepopsec};%
  \hspace{8pt}{\popdisplay\fontsize{15}{17}\selectfont\color{popink}\MakeUppercase{#1}}\par
  \vspace{4pt}\noindent{\color{poporange}\rule{36pt}{3.6pt}}
  \end{minipage}\par\nopagebreak\vspace{8pt}
}
\newcounter{popsub}
\newcommand{\courssub}[1]{%
  \stepcounter{popsub}%
  \par\vspace{9pt}\noindent{\popxbold\fontsize{11.5}{13}\selectfont\color{popink}{\color{poporange}\thepopsec.\thepopsub}\hspace{6pt}#1}\par\nopagebreak\vspace{4pt}
}

% ============================================================
% Boîtes pédagogiques — même recette : fond blanc, bordure épaisse colorée,
% ombre dure encre, badge plein. Titre optionnel [#1].
% ============================================================
\tcbset{popbase/.style={breakable,enhanced,colback=white,boxrule=2.3pt,arc=9pt,
  shadow={3pt}{-3pt}{0pt}{popink},left=14pt,right=14pt,top=11pt,bottom=11pt,before skip=11pt,after skip=15pt,
  fontupper=\popsemi\fontsize{10.6}{14.5}\selectfont\color{popink},
  fontlower=\popsemi\fontsize{10.6}{14.5}\selectfont\color{popink}}}
% Coche dessinée (le glyphe ✓ n'existe pas dans les polices du document)
\newcommand{\popcheck}[1]{\tikz[baseline=-0.5ex]\draw[#1,line width=1.6pt,line cap=round,line join=round](-0.13,0.0)--(-0.04,-0.09)--(0.13,0.1);}
\providecommand{\coche}{\popcheck{popgreen}}
\providecommand{\resultat}[1]{\par\smallskip\noindent\fcolorbox{popgreen}{popgreenL}{\parbox{\dimexpr\linewidth-2\fboxsep-2\fboxrule\relax}{\textbf{Résultat.} #1}}\par\smallskip}
\newcommand{\popboxhead}[3]{\poppill{#1}{#2}{white}\ifx\relax#3\relax\else\hspace{6pt}{\popxbold\fontsize{10.6}{12}\selectfont\color{popink}#3}\fi\par\vspace{7pt}}

\newtcolorbox{aretenir}[1][]{popbase,colframe=popgreen,before upper={\popboxhead{À retenir}{popgreen}{#1}}}
% Texte en allemand (document de lecture, dialogue, extrait) : cadre
% d'étude. Un titre optionnel (source, auteur) s'affiche dans l'en-tête.
\newtcolorbox{textebox}[1][]{popbase,colframe=popblue,before upper={\popboxhead{Text}{popblue}{#1}\begin{otherlanguage}{ngerman}},after upper={\end{otherlanguage}}}
% Marques ✓ / ✗ (forme correcte / fautive) souvent utilisées par les modèles
\providecommand{\cross}{\textcolor{poppink}{\ensuremath{\times}}}
\providecommand{\xmark}{\cross}\providecommand{\crossmark}{\cross}\providecommand{\cmark}{\ensuremath{\checkmark}}
% Ligne de réponse / justification à compléter (inventée par les modèles)
\providecommand{\justification}[1]{\par\noindent\textit{Justification~:} \dotfill\par}
% \textipa (package tipa non chargé) : la police ne couvre pas l'API -> simple passage
\providecommand{\textipa}[1]{#1}
% Soulignement (ulem non chargé) : \uline, \ul
\providecommand{\uline}[1]{\underline{#1}}\providecommand{\ul}[1]{\underline{#1}}
% \glossaire{\item ...} inventée par les modèles : liste de glossaire
\providecommand{\glossaire}[1]{\begin{itemize}#1\end{itemize}}
% \alert (beamer) inventée par les modèles : texte mis en évidence
\providecommand{\alert}[1]{\textbf{\textcolor{poppink}{#1}}}
% variantes ASCII d'ordinaux écrites en macro
\providecommand{\iere}{\textsuperscript{re}}\providecommand{\ieme}{\textsuperscript{e}}\providecommand{\ere}{\textsuperscript{re}}\providecommand{\eme}{\textsuperscript{e}}\providecommand{\er}{\textsuperscript{er}}\providecommand{\ie}{\textsuperscript{e}}
% Ordinaux français écrits en macro par les modèles : 1\ière, 3\ième
\newcommand{\ière}{\textsuperscript{re}}\newcommand{\ième}{\textsuperscript{e}}
% \exemple (inventée par les modèles) : étiquette « Exemple : »
\providecommand{\exemple}{\textbf{Exemple~:}\ }
% \square utilisable aussi hors mode math (cases à cocher)
\AtBeginDocument{\let\squareMath\square\renewcommand{\square}{\ensuremath{\squareMath}}}
% Mot ou phrase en allemand à l'intérieur d'un paragraphe français
\newcommand{\all}[1]{\ifmmode\text{\textit{#1}}\else\begingroup\selectlanguage{ngerman}\itshape #1\endgroup\fi}
\let\esp\all % alias toléré
\newtcolorbox{exemplebox}[1][]{popbase,colframe=poporange,before upper={\popboxhead{Exemple}{poporange}{#1}}}
\newtcolorbox{definition}[1][]{popbase,colframe=popblue,before upper={\popboxhead{Définition}{popblue}{#1}}}
\newtcolorbox{propriete}[1][]{popbase,colframe=poppurple,before upper={\popboxhead{Propriété}{poppurple}{#1}}}
\newtcolorbox{methode}[1][]{popbase,colframe=popgold,before upper={\popboxhead{Méthode}{popgold}{#1}}}
\newtcolorbox{attention}[1][]{popbase,colframe=poppink,before upper={\popboxhead{Attention}{poppink}{#1}}}
\newtcolorbox{experience}[1][]{popbase,colframe=popblue,colback=popblueL,before upper={\popboxhead{Expérience}{popblue}{#1}}}
% « Le savais-tu ? » : carte violette pleine (contexte, culture, Cameroun)
\newtcolorbox{savaistu}[1][]{popbase,colback=poppurple,colframe=popink,
  fontupper=\popsemi\fontsize{10.4}{14.2}\selectfont\color{white},
  before upper={\poppill{Le savais-tu\,?}{white}{poppurple}\ifx\relax#1\relax\else\hspace{6pt}{\popxbold\color{white}#1}\fi\par\vspace{7pt}}}
% Exercice résolu : énoncé en haut, \tcblower puis la solution sur fond crème
\newcounter{popexo}\newcounter{popexr}
\newtcolorbox{exoresolu}[1][]{popbase,colframe=poporange,colbacklower=popcream,
  segmentation style={popink,line width=1.2pt,dashed},
  before upper={\stepcounter{popexr}\popboxhead{Exercice résolu \thepopexr}{poporange}{#1}},
  before lower={\poppill{Solution}{popink}{popcream}\par\vspace{6pt}}}
% Exercice (non résolu en place) — la correction est dans la partie Corrigés
\newtcolorbox{exercice}[1][]{popbase,colframe=popink,
  before upper={\stepcounter{popexo}\popboxhead{Exercice \thepopexo}{popink}{#1}}}
% Corrigé
\newtcolorbox{corrige}[1][]{popbase,colframe=popgreen,colback=popgreenL,
  before upper={\popboxhead{Corrigé}{popgreen}{#1}}}
% Objectifs / prérequis / bilan
\newtcolorbox{objectifs}{popbase,colframe=popink,colback=poporangeL,
  before upper={\popboxhead{Objectifs de la leçon}{popink}{}}}
\newtcolorbox{prerequis}{popbase,colframe=popink,colback=popblueL,
  before upper={\popboxhead{Prérequis}{popblue}{}}}
\newtcolorbox{bilan}{popbase,colframe=popink,colback=white,boxrule=2.8pt,
  before upper={\popboxhead{Fiche bilan}{poporange}{L'essentiel à connaître pour l'examen}}}
% Figure encadrée avec légende
\newtcolorbox{popfigure}[1][]{enhanced,breakable=false,colback=white,colframe=popline,boxrule=1.4pt,arc=8pt,
  left=10pt,right=10pt,top=10pt,bottom=8pt,before skip=10pt,after skip=12pt,halign=center,
  lower separated=false,halign lower=center,
  fontlower=\popsemi\fontsize{9}{11.5}\selectfont\color{popmuted},
  #1}
\newcommand{\legende}[1]{\par\vspace{6pt}{\popsemi\fontsize{9}{11.5}\selectfont\color{popmuted}#1}}

% ============================================================
% Signature OnBuch+
% ============================================================
\newcommand{\poplogoinline}[1]{{\poplogo\fontsize{#1}{#1}\selectfont\color{popink}OnBuch\color{poporange}+}}
\newcommand{\signatureonbuch}{%
  \begin{tcolorbox}[enhanced,colback=popink,colframe=popink,boxrule=2.2pt,arc=10pt,
      drop shadow=poporange,left=14pt,right=14pt,top=10pt,bottom=10pt,before skip=0pt,after skip=0pt]
    \tikz[baseline=-0.7ex]\node[circle,fill=popgreen,draw=popcream,line width=1.3pt,minimum size=22pt,inner sep=0]{\popcheck{white}};%
    \hspace{9pt}\begin{minipage}[c]{0.62\linewidth}
      {\popxbold\fontsize{12}{13}\selectfont\color{popcream}Signé\ \textbullet\ {\poplogo OnBuch\color{poporange}+}}\par\vspace{2pt}
      {\raggedright\popsemi\fontsize{8}{10.5}\selectfont\color{popline}\MakeUppercase{Équipe pédagogique OnBuch+}\\\MakeUppercase{Conforme au programme officiel MINESEC}\par}
    \end{minipage}\hfill
    {\poplogo\fontsize{22}{22}\selectfont\color{popcream}OnBuch\color{poporange}+}
  \end{tcolorbox}%
}

% ============================================================
% Page de garde
% #1 titre · #2 matière · #3 niveau · #4 module · #5 n° leçon · #6 durée ·
% #7 description / objectifs (texte ou itemize)
% ============================================================
\newcommand{\pagegarde}[7]{%
  \thispagestyle{empty}
  \newgeometry{a4paper,margin=1.7cm,top=1.6cm,bottom=1.4cm}%
  \begin{tikzpicture}[remember picture,overlay]
    \fill[poporange!18] ([xshift=-3.2cm,yshift=-3.6cm]current page.north east) circle (4.6cm);
    \fill[poppurple!14] ([xshift=2.4cm,yshift=7.6cm]current page.south west) circle (3.8cm);
  \end{tikzpicture}%
  {\poplogo\fontsize{30}{30}\selectfont\color{popink}OnBuch\color{poporange}+}\hfill
  \raisebox{0.6ex}{\poppill{Cours signé \textbullet\ Premium}{popink}{popcream}}\par
  \vspace{1pt}\popsquiggle{poppurple}\par
  \vspace{8pt}
  \begin{tcolorbox}[enhanced,colback=poporange,colframe=popink,boxrule=2.8pt,arc=13pt,
      drop shadow=popink,left=18pt,right=18pt,top=12pt,bottom=13pt,after skip=10pt]
    \poppill{#2 \textbullet\ #3}{popink}{popcream}\hspace{5pt}\poppill{#4}{white}{popink}\par\vspace{8pt}
    {\popdisplay\fontsize{14}{14}\selectfont\color{popink}LEÇON #5}\par\vspace{4pt}
    {\raggedright\hyphenpenalty=10000\exhyphenpenalty=10000\popdisplay\fontsize{25}{27}\selectfont\color{popcream}\MakeUppercase{#1}\par}
    \vspace{8pt}
    {\popxbold\fontsize{9.5}{9.5}\selectfont\color{popink}\MakeUppercase{Durée conseillée : #6}}
  \end{tcolorbox}
  \begin{tcolorbox}[enhanced,colback=white,colframe=popink,boxrule=2.2pt,arc=10pt,
      drop shadow=popink,left=16pt,right=16pt,top=10pt,bottom=10pt,after skip=8pt]
    \poppill{Dans cette leçon}{poppurple}{white}\par\vspace{5pt}
    \popsemi\fontsize{9.3}{12.6}\selectfont\color{popink}#7
  \end{tcolorbox}
  \noindent\begin{minipage}[t]{0.487\linewidth}
    \begin{tcolorbox}[enhanced,colback=popgreen,colframe=popink,boxrule=2.2pt,arc=10pt,
        drop shadow=popink,left=11pt,right=11pt,top=10pt,bottom=10pt,height=3.1cm,valign=center]
      \tcbox[colback=white,colframe=popink,boxrule=1.3pt,arc=5pt,boxsep=0pt,nobeforeafter,
        left=3pt,right=3pt,top=3pt,bottom=3pt,valign=center]{\qrcode[height=1.9cm]{https://play.google.com/store/apps/details?id=com.luvvix.onbuch&hl=fr}}%
      \hspace{8pt}\begin{minipage}[c]{0.5\linewidth}
        {\popdisplay\fontsize{11}{12.5}\selectfont\color{white}\MakeUppercase{Télécharge OnBuch}}\par\vspace{4pt}
        {\raggedright\popsemi\fontsize{8.4}{11}\selectfont\color{white}Gratuit \textbullet\ Google Play\\Tuteur IA, annales, concours, résultats\par}
      \end{minipage}
    \end{tcolorbox}
  \end{minipage}\hfill
  \begin{minipage}[t]{0.487\linewidth}
    \begin{tcolorbox}[enhanced,colback=poppurple,colframe=popink,boxrule=2.2pt,arc=10pt,
        drop shadow=popink,left=11pt,right=11pt,top=10pt,bottom=10pt,height=3.1cm,valign=center]
      \tikz[baseline=-0.7ex]\node[circle,fill=white,draw=popink,line width=1.3pt,minimum size=1.6cm,inner sep=0]{\popdisplay\fontsize{24}{24}\selectfont\color{poppurple}+};%
      \hspace{8pt}\begin{minipage}[c]{0.6\linewidth}
        {\popdisplay\fontsize{11}{12.5}\selectfont\color{white}\MakeUppercase{Passe à OnBuch+}}\par\vspace{4pt}
        {\raggedright\popsemi\fontsize{8.4}{11}\selectfont\color{white}Tous les cours signés, Tuteur IA illimité, fiches PDF — onglet OnBuch+ de l'app\par}
      \end{minipage}
    \end{tcolorbox}
  \end{minipage}
  \vspace{8pt}
  \popcontenu
  \vfill
  \signatureonbuch
  \restoregeometry
  \newpage
}

% Rangée « Ce cours contient » (page de garde)
\newcommand{\poptile}[3]{% #1 couleur #2 gros texte #3 légende
  \begin{tcolorbox}[enhanced,colback=#1,colframe=popink,boxrule=2pt,arc=9pt,shadow={2.5pt}{-2.5pt}{0pt}{popink},
      left=6pt,right=6pt,top=9pt,bottom=9pt,halign=center,valign=center,height=1.75cm,width=\linewidth,nobeforeafter]
    {\popdisplay\fontsize{12.5}{13}\selectfont\color{white}\MakeUppercase{#2}}\par\vspace{3pt}
    {\centering\popsemi\fontsize{7.8}{9.5}\selectfont\color{white}#3\par}
  \end{tcolorbox}}
\newcommand{\popcontenu}{%
  \par\noindent{\popxboldls\fontsize{7.5}{7.5}\selectfont\color{popmuted}CE COURS SIGNÉ CONTIENT}\par\vspace{7pt}
  \noindent\begin{minipage}[t]{0.235\linewidth}\poptile{popblue}{Cours}{complet, illustré, pas à pas}\end{minipage}\hfill
  \begin{minipage}[t]{0.235\linewidth}\poptile{poporange}{Méthodes}{et exercices résolus}\end{minipage}\hfill
  \begin{minipage}[t]{0.235\linewidth}\poptile{poppink}{Exercices}{type examen + corrigés}\end{minipage}\hfill
  \begin{minipage}[t]{0.235\linewidth}\poptile{popgreen}{Fiche bilan}{l'essentiel à retenir}\end{minipage}\par}

% Sommaire
\newtcolorbox{sommaire}{enhanced,colback=white,colframe=popink,boxrule=2.2pt,arc=10pt,
  drop shadow=popink,left=18pt,right=18pt,top=13pt,bottom=13pt,before skip=6pt,after skip=20pt,
  before upper={\poppill{Sommaire}{poppurple}{white}\par\vspace{9pt}\popsemi\fontsize{10.6}{14.5}\selectfont\color{popink}}}

% Signature de fin de cours — poussée tout en bas de la dernière page
\newcommand{\findecours}{%
  \par\vfill\nopagebreak
  \begin{center}\popsquiggle{poporange}\end{center}\vspace{4pt}
  \signatureonbuch
}
\AtBeginDocument{\def\llbracket{\mathopen{[\mkern-3.5mu[}}\def\rrbracket{\mathclose{]\mkern-3.5mu]}}} % intervalles d'entiers (glyphe ⟦ absent de la police)
% Glyphes absents de Fira Math : redéfinis à partir de glyphes disponibles
\AtBeginDocument{%
  \def\setminus{\mathbin{\backslash}}\let\smallsetminus\setminus
  \def\vdots{\mathord{\vbox{\baselineskip3.2pt\lineskiplimit0pt\kern4pt\hbox{.}\hbox{.}\hbox{.}}}}%
  \def\ddots{\mathinner{\mkern1mu\raise6.5pt\hbox{.}\mkern2mu\raise3.6pt\hbox{.}\mkern2mu\raise0.7pt\hbox{.}\mkern1mu}}%
  \def\triangle{\mathord{\scalebox{0.9}{$\symup{\Delta}$}}}%
  \def\nabla{\mathord{\raisebox{\depth}{\scalebox{1}[-1]{$\symup{\Delta}$}}}}%
  \def\checkmark{\popcheck{popdark}}%
  \def\star{\mathbin{\ast}}\def\bigstar{\mathord{\ast}}%
  \def\diamond{\mathbin{\scalebox{0.6}{$\Diamond$}}}%
  \def\bigcup{\mathop{\mathchoice{\raisebox{-0.3em}{\scalebox{1.7}{$\cup$}}}{\scalebox{1.25}{$\cup$}}{\cup}{\cup}}\displaylimits}%
  \def\bigcap{\mathop{\mathchoice{\raisebox{-0.3em}{\scalebox{1.7}{$\cap$}}}{\scalebox{1.25}{$\cap$}}{\cap}{\cap}}\displaylimits}%
  \def\lesseqqgtr{\mathrel{\lessgtr}}\def\bigoplus{\mathop{\oplus}\displaylimits}%
}
\providecommand{\ssi}{si et seulement si} % abréviation fréquente
\AtBeginDocument{\providecommand{\wideparen}{\overparen}} % arc de cercle

\begin{document}

\pagegarde{Les fêtes religieuses ou nationales}{Allemand}{1ère A}{Chapitre 5 : Citoyenneté}{ALA22}{2 h}{Cette leçon propose la lecture et l'exploitation d'un texte sur les fêtes religieuses et nationales dans les pays germanophones (Weihnachten, Ostern, Tag der Deutschen Einheit). Elle permet d'acquérir le vocabulaire thématique des fêtes et deux outils grammaticaux essentiels : les dates avec les nombres ordinaux et les verbes à particule séparable.
\vspace{4pt}
\begin{multicols}{2}\small
\begin{itemize}
\item À la fin de cette leçon, je sais comprendre globalement puis en détail un texte sur les fêtes religieuses ou nationales.
\item Je sais employer correctement le vocabulaire des fêtes (Weihnachten, Ostern, der Nationalfeiertag, die Parade, die Fahne, der Gottesdienst) avec articles et pluriels.
\item Je sais former et employer les nombres ordinaux (der erste, der zweite, der zwanzigste...).
\item Je sais dire et écrire une date en allemand (am 20. Mai, der 20. Mai 2024).
\item Je sais reconnaître un verbe séparable et le conjuguer correctement au présent et au parfait.
\item Je sais répondre par écrit et oralement à des questions de compréhension sur un texte.
\end{itemize}\end{multicols}}

\begin{sommaire}
\begin{multicols}{2}
\begin{enumerate}
\item Texte support : Feste in Deutschland
\item Compréhension du texte
\item Lexique thématique : les fêtes
\item Grammaire 1 : les dates et les nombres ordinaux
\item Grammaire 2 : les verbes séparables
\item Commentaire et production guidée
\item Activité d'intégration
\item Méthodes et exercices résolus
\item Exercices
\item Corrigés des exercices
\item Fiche bilan
\end{enumerate}
\end{multicols}
\end{sommaire}

\begin{objectifs}
\begin{itemize}
\item À la fin de cette leçon, je sais comprendre globalement puis en détail un texte sur les fêtes religieuses ou nationales.
\item Je sais employer correctement le vocabulaire des fêtes (Weihnachten, Ostern, der Nationalfeiertag, die Parade, die Fahne, der Gottesdienst) avec articles et pluriels.
\item Je sais former et employer les nombres ordinaux (der erste, der zweite, der zwanzigste...).
\item Je sais dire et écrire une date en allemand (am 20. Mai, der 20. Mai 2024).
\item Je sais reconnaître un verbe séparable et le conjuguer correctement au présent et au parfait.
\item Je sais répondre par écrit et oralement à des questions de compréhension sur un texte.
\item Je sais rédiger un court paragraphe guidé présentant une fête camerounaise à un correspondant allemand.
\item Je sais comparer les fêtes du Cameroun et des pays germanophones en quelques phrases simples.
\end{itemize}
\end{objectifs}

\begin{prerequis}
\begin{itemize}
\item Les nombres cardinaux de 1 à 100 (Seconde).
\item Les mois, les jours de la semaine et les saisons (Seconde).
\item La conjugaison des verbes réguliers au présent et la place du verbe en position 2.
\item Le parfait avec haben et sein (début de 1ère).
\item Les articles définis et indéfinis au nominatif et à l'accusatif.
\end{itemize}
\end{prerequis}

\newpage

\coursec{Texte support : Feste in Deutschland}

\courssub{Situation d'accroche et problématique}

Le 20 mai, à Yaoundé comme à Douala, le Cameroun célèbre sa fête nationale : défilés militaires sur le boulevard, fanfares, discours officiels, drapeaux hissés partout. Le 25 décembre, les familles chrétiennes se retrouvent à l'église, vêtues de leurs plus beaux habits, avant le grand repas en famille. Ces moments font partie de ta vie. Mais que fête-t-on en Allemagne, en Autriche et en Suisse ? Comment dit-on « le 3 octobre » ou « le premier mai » en allemand ? Et pourquoi le verbe allemand se coupe-t-il parfois en deux dans la phrase, comme dans \all{Im Frühling findet Ostern statt} ?

\textbf{Problématique :} comment comprendre et raconter, en allemand correct, les fêtes religieuses et nationales des pays germanophones, en employant les dates, le vocabulaire des fêtes et les verbes à particule séparable ?

Dans cette première partie de la leçon, tu vas lire un texte original : André, un jeune Camerounais de Yaoundé, écrit à son correspondant allemand Lukas. André a passé une année scolaire en Allemagne et il raconte les fêtes qu'il a découvertes là-bas. Tu vas d'abord lire le texte globalement, puis le décortiquer phrase par phrase, avec son glossaire.

\courssub{Présentation du texte}

Avant de lire, voici la situation de communication :

\begin{itemize}
\item \textbf{Qui parle ?} André, 16 ans, élève camerounais en classe de Première, en échange scolaire à Munich.
\item \textbf{À qui ?} À Lukas, son correspondant allemand, par courriel.
\item \textbf{De quoi parle-t-il ?} Des fêtes religieuses (Noël, Pâques) et de la fête nationale allemande (le 3 octobre).
\item \textbf{Pourquoi ?} Pour comparer avec les fêtes du Cameroun et inviter Lukas à découvrir le 20 mai camerounais.
\end{itemize}

\begin{methode}[Bien lire un texte allemand en quatre étapes]
\begin{enumerate}
\item \textbf{Lire le titre et la première phrase} pour deviner le thème. Ici : \all{Feste in Deutschland} annonce clairement le sujet : les fêtes.
\item \textbf{Faire une première lecture globale, sans dictionnaire.} Ne t'arrête pas sur chaque mot inconnu : cherche d'abord le sens général.
\item \textbf{Repérer les mots transparents}, c'est-à-dire les mots qui ressemblent au français : \all{die Parade} (la parade), \all{der Nationalfeiertag} (la fête nationale), \all{die Familie}, \all{die Kirche}, \all{das Konzert}.
\item \textbf{Faire une deuxième lecture détaillée} avec le glossaire : souligne les dates, les noms de fêtes et les verbes.
\end{enumerate}
\end{methode}

\courssub{Lecture globale du texte}

\begin{textebox}[Feste in Deutschland — Brief von André an Lukas]
Lieber Lukas,\\
Deutschland hat viele schöne Feste, und ich möchte dir heute von meinen Erlebnissen erzählen. Am 25. Dezember feiert man hier Weihnachten: die Familien schmücken einen Tannenbaum mit Lichtern und Kugeln, und die Kinder stehen sehr früh auf, um die Geschenke zu öffnen. Am Abend gibt es ein großes Essen mit der ganzen Familie.\\
Im Frühling findet Ostern statt: die Kinder suchen bunte Eier im Garten, denn der Osterhase hat sie versteckt. Das finde ich sehr lustig!\\
Am 3. Oktober ist der Nationalfeiertag, der Tag der Deutschen Einheit: in den Städten gibt es Paraden und Konzerte, überall hängen schwarz-rot-goldene Fahnen, und am Abend sieht man ein großes Feuerwerk. Viele Menschen hören auch die Rede des Bundespräsidenten im Fernsehen.\\
Am Sonntag gehen viele Familien zum Gottesdienst in die Kirche, besonders an Weihnachten und Ostern.\\
Bei uns in Kamerun ist der 20. Mai der Nationalfeiertag, mit Paraden und Reden, und am 25. Dezember gehen wir auch in die Kirche. Unsere Feste sind also ähnlich und doch anders!\\
Ich mache gern bei diesen Festen mit, denn man lernt das Land und die Leute besser kennen. Und du, Lukas, welches Fest magst du am liebsten?\\
Viele Grüße\\
André
\end{textebox}

\textbf{Glossaire (mots difficiles du texte) :}

\begin{center}
\begin{tabularx}{\linewidth}{|l|X|}
\hline
\rowcolor{popblueL} \textbf{Deutsch} & \textbf{Français} \\
\hline
die Feier, -n & la fête, la célébration \\
\hline
der Feiertag, -e & le jour férié \\
\hline
feiern (hat gefeiert) & fêter, célébrer \\
\hline
schmücken (hat geschmückt) & décorer, orner \\
\hline
stattfinden (fand statt, hat stattgefunden) & avoir lieu \\
\hline
die Kerze, -n & la bougie \\
\hline
das Feuerwerk, -e & le feu d'artifice \\
\hline
das Erlebnis, -se & l'expérience vécue, l'aventure \\
\hline
der Tannenbaum, -bäume & le sapin (de Noël) \\
\hline
die Kugel, -n & la boule (de décoration) \\
\hline
das Geschenk, -e & le cadeau \\
\hline
der Osterhase, -n & le lièvre de Pâques \\
\hline
verstecken (hat versteckt) & cacher \\
\hline
die Rede, -n & le discours \\
\hline
der Bundespräsident, -en & le président fédéral \\
\hline
das Fernsehen & la télévision \\
\hline
der Gottesdienst, -e & l'office religieux, la messe \\
\hline
ähnlich & semblable \\
\hline
\end{tabularx}
\end{center}

\begin{definition}[Feier, Feiertag, feiern : trois mots de la même famille]
\begin{itemize}
\item \all{der Feiertag, -e} = le jour férié (jour officiel de fête, où l'on ne travaille pas) : \all{Der 3. Oktober ist ein Feiertag.} « Le 3 octobre est un jour férié. »
\item \all{die Feier, -n} = la fête, la célébration (l'événement) : \all{Die Feier beginnt um 18 Uhr.} « La fête commence à 18 heures. »
\item \all{feiern} = fêter, célébrer (le verbe) : \all{Wir feiern Weihnachten mit der Familie.} « Nous fêtons Noël en famille. »
\end{itemize}
Ces trois mots appartiennent à la même famille : retiens-les ensemble, avec leur article et leur pluriel.
\end{definition}

\courssub{Première compréhension : repérage global}

Après ta première lecture, tu dois pouvoir répondre à ces questions simples :

\begin{enumerate}
\item De quelles fêtes parle André ? (Trois fêtes allemandes et une fête camerounaise.)
\item Quelles dates sont citées dans le texte ?
\item Que font les enfants à Noël et à Pâques ?
\end{enumerate}

\begin{exoresolu}[Repérage global du texte]
\textbf{Énoncé :} Réponds en allemand, par une phrase complète : \emph{Welche Feste beschreibt André in seinem Brief?}
\tcblower
\textbf{Solution détaillée :} On relit le texte et on relève les noms de fêtes : \all{Weihnachten} (25. Dezember), \all{Ostern} (im Frühling), \all{der Tag der Deutschen Einheit / der Nationalfeiertag} (3. Oktober) et, pour le Cameroun, \all{der 20. Mai}. On construit une phrase complète avec le verbe conjugué en deuxième position :\\
\all{André beschreibt Weihnachten, Ostern und den Nationalfeiertag in Deutschland; er nennt auch den 20. Mai in Kamerun.}\\
« André décrit Noël, Pâques et la fête nationale en Allemagne ; il mentionne aussi le 20 mai au Cameroun. »
\end{exoresolu}

\courssub{Lecture détaillée : dates, verbes séparables, informations clés}

\textbf{1. Repérage des dates.} Souligne dans le texte toutes les dates. Tu remarques la structure allemande : \all{am 25. Dezember}, \all{am 3. Oktober}, \all{am 20. Mai}. La préposition \all{am} (= an + dem) introduit la date, et le nombre ordinal se termine par un point : \all{der 3. Oktober} se lit \all{der dritte Oktober} « le trois octobre ». Cette structure sera étudiée en détail dans la partie grammaire de la leçon.

\textbf{2. Repérage des verbes séparables.} Le texte contient trois verbes à particule séparable. Observe leur comportement dans la phrase :

\begin{center}
\begin{tabularx}{\linewidth}{|l|X|X|}
\hline
\rowcolor{popblueL} \textbf{Verbe} & \textbf{Phrase du texte} & \textbf{Traduction} \\
\hline
stattfinden & \all{Im Frühling findet Ostern statt.} & « Au printemps, Pâques a lieu. » \\
\hline
aufstehen & \all{Die Kinder stehen sehr früh auf.} & « Les enfants se lèvent très tôt. » \\
\hline
mitmachen & \all{Ich mache gern bei diesen Festen mit.} & « Je participe volontiers à ces fêtes. » \\
\hline
\end{tabularx}
\end{center}

Tu observes la règle essentielle : à la forme conjuguée, le verbe se \textbf{coupe en deux} : la partie verbale (\all{findet}, \all{stehen}, \all{mache}) prend la deuxième position dans la phrase, et la particule (\all{statt}, \all{auf}, \all{mit}) part \textbf{à la fin de la phrase}. C'est une différence majeure avec le français, où le verbe reste toujours en un seul bloc.

\textbf{3. Informations sur chaque fête.} Complète mentalement ce tableau en relisant le texte :

\begin{center}
\begin{tabularx}{\linewidth}{|l|X|X|}
\hline
\rowcolor{popgreenL} \textbf{Fête} & \textbf{Date / période} & \textbf{Activités} \\
\hline
Weihnachten & am 25. Dezember & Tannenbaum schmücken, Geschenke öffnen, großes Essen, Gottesdienst \\
\hline
Ostern & im Frühling & bunte Eier suchen, der Osterhase \\
\hline
Tag der Deutschen Einheit & am 3. Oktober & Paraden, Fahnen, Konzerte, Rede, Feuerwerk \\
\hline
Nationalfeiertag in Kamerun & am 20. Mai & Paraden und Reden \\
\hline
\end{tabularx}
\end{center}

\begin{attention}[Faux amis et confusions fréquentes]
\begin{itemize}
\item Ne confonds pas \all{feiern} (fêter, célébrer) et \all{feuern} (tirer un coup de feu ; licencier quelqu'un) : \all{Wir feiern Ostern.} « Nous fêtons Pâques », mais \all{Der Chef feuert den Arbeiter.} « Le patron licencie l'ouvrier. » Une seule lettre, un sens totalement différent !
\item \all{die Fahne} signifie « le drapeau » : \all{Die Fahnen hängen in den Straßen.} « Les drapeaux sont accrochés dans les rues. » Mais dans le langage familier, \all{eine Fahne haben} signifie « avoir mauvaise haleine » : le contexte décide !
\item Attention à la majuscule : tous les noms de fêtes prennent une majuscule en allemand : \all{Weihnachten}, \all{Ostern}, \all{der Nationalfeiertag}.
\item Ne traduis pas « avoir lieu » par \all{haben} : il faut le verbe séparable \all{stattfinden} : \all{Das Fest findet im Mai statt.} « La fête a lieu en mai. »
\item Attention à \all{die Kirche} : c'est « l'église » (le bâtiment et l'institution), et non « la chaire ». \all{Wir gehen in die Kirche.} « Nous allons à l'église. »
\end{itemize}
\end{attention}

\courssub{Carte mentale du vocabulaire des fêtes}

\begin{popfigure}
\begin{tikzpicture}[mindmap, grow cyclic, every node/.style={concept, align=center}, concept color=popblue!60, root concept/.append style={concept color=popblue, font=\large\bfseries}, level 1/.append style={level distance=3.2cm, sibling angle=120}, level 2/.append style={level distance=2.4cm, sibling angle=45, concept color=popblueL, font=\small}]
\node[root concept]{Feste}
  child { node[align=center]{Weihnachten\\ 25. Dezember}
    child { node {der Tannenbaum} }
    child { node {das Geschenk} }
    child { node {die Kerze} }
  }
  child { node[align=center]{Ostern\\ im Frühling}
    child { node {das Ei} }
    child { node {der Osterhase} }
    child { node {der Gottesdienst} }
  }
  child { node[align=center]{Nationalfeiertag\\ 3. Oktober}
    child { node {die Parade} }
    child { node {die Fahne} }
    child { node {das Feuerwerk} }
  };
\end{tikzpicture}
\legende{Carte mentale : les trois grandes fêtes du texte et leur vocabulaire-clé.}
\end{popfigure}

\courssub{Écoute et relecture à voix haute : la prononciation}

Relis le texte à voix haute après le professeur. Travaille en particulier :

\begin{itemize}
\item \textbf{Les deux sons du « ch » allemand.} Après \emph{a, o, u, au}, le « ch » se prononce au fond de la gorge, comme un « r » grasseyé doux : \all{Weihnachten} « vé-nakh-tèn », \all{acht} « akht », \all{suchen} « zou-khèn ». Après \emph{i, e} et en début de mot devant voyelle claire, le « ch » est doux, proche d'un « ch » français sifflé entre les dents : \all{ich} « ihh », \all{Kirche} « kir-h'e », \all{nicht} « niht ».
\item \textbf{Le « w » allemand se prononce [v]} comme le « v » français : \all{Weihnachten} « vé-nakh-tèn », \all{das Feuerwerk} « foïeur-vèrque », \all{der Weihnachtsmarkt} « vé-nakhts-markt ».
\item \textbf{Le « v » allemand se prononce souvent [f]} : \all{viele} « fi-lé », \all{viel} « fil », \all{der Vater} « fa-teur ».
\item \textbf{Le « s » en début de mot devant voyelle se prononce [z]} : \all{der Sonntag} « zon-tak », \all{suchen} « zou-khèn », \all{sieben} « zi-bèn ».
\item \textbf{Le « st » et « sp » en début de mot ou de syllabe} se prononcent « cht » et « chp » : \all{stattfinden} « chtat-fin-dèn », \all{stehen} « chté-èn », \all{sprechen} « chpré-h'è ».
\end{itemize}

\begin{savaistu}[Les fêtes nationales des pays germanophones]
Le 3 octobre, \all{der Tag der Deutschen Einheit} (« le Jour de l'Unité allemande »), commémore la réunification de l'Allemagne de l'Est et de l'Allemagne de l'Ouest en 1990, après la chute du Mur de Berlin en 1989. C'est un jour férié dans tout le pays. L'Autriche, elle, fête sa fête nationale le 26 octobre, et la Suisse le 1\textsuperscript{er} août, avec des feux d'artifice et des lanternes. Et le Cameroun ? Son 20 mai, jour de l'unité nationale, ressemble beaucoup au 3 octobre allemand : défilés, drapeaux et discours officiels dans les deux pays !
\end{savaistu}

\begin{exercice}[À toi de jouer : compréhension détaillée]
Réponds en allemand, par des phrases complètes :
\begin{enumerate}
\item Wann feiert man in Deutschland Weihnachten?
\item Was machen die Kinder am Weihnachtsmorgen?
\item Wann findet Ostern statt, und was suchen die Kinder?
\item Was ist der 3. Oktober, und was sieht man am Abend?
\item Welches Fest gibt es in Kamerun am 20. Mai?
\end{enumerate}
\end{exercice}

\begin{corrige}[Exercice : compréhension détaillée]
\begin{enumerate}
\item \all{Man feiert Weihnachten am 25. Dezember.} « On fête Noël le 25 décembre. »
\item \all{Die Kinder stehen früh auf und öffnen die Geschenke.} « Les enfants se lèvent tôt et ouvrent les cadeaux. » (Attention : la particule \all{auf} va en fin de phrase.)
\item \all{Ostern findet im Frühling statt; die Kinder suchen bunte Eier im Garten.} « Pâques a lieu au printemps ; les enfants cherchent des œufs colorés dans le jardin. »
\item \all{Der 3. Oktober ist der Nationalfeiertag, der Tag der Deutschen Einheit; am Abend sieht man ein Feuerwerk.} « Le 3 octobre est la fête nationale, le Jour de l'Unité allemande ; le soir, on voit un feu d'artifice. »
\item \all{Am 20. Mai ist in Kamerun der Nationalfeiertag, mit Paraden und Reden.} « Le 20 mai, c'est la fête nationale au Cameroun, avec des défilés et des discours. »
\end{enumerate}
\end{corrige}

\begin{aretenir}[L'essentiel de cette section]
\begin{itemize}
\item Le texte d'André présente trois fêtes allemandes : \all{Weihnachten} (25. Dezember), \all{Ostern} (im Frühling) et \all{der Tag der Deutschen Einheit} (3. Oktober), qu'il compare au 20 mai camerounais.
\item Vocabulaire-clé : \all{die Feier, der Feiertag, feiern, schmücken, stattfinden, die Kerze, das Feuerwerk, die Fahne, die Parade, der Gottesdienst}.
\item Les dates s'expriment avec \all{am} + nombre ordinal + point : \all{am 3. Oktober} « le 3 octobre ».
\item Les verbes séparables se coupent en deux : \all{Ostern findet im Frühling statt.} — la particule va à la fin de la phrase.
\item Méthode de lecture : titre et première phrase, lecture globale sans dictionnaire, mots transparents, puis lecture détaillée avec le glossaire.
\end{itemize}
\end{aretenir}

\coursec{Compréhension du texte}

Dans la section précédente, nous avons lu et écouté le texte \emph{Feste in Deutschland}, qui présente les grandes fêtes religieuses et nationales célébrées dans les pays germanophones. Nous allons maintenant vérifier notre compréhension, d'abord de manière globale (le sens général), puis de manière détaillée (les informations précises : dates, lieux, actions). C'est une compétence essentielle : à l'examen, on vous demandera toujours de répondre à des questions sur un texte \textbf{par des phrases complètes en allemand}.

\courssub{La méthode pour répondre à une question de compréhension}

Avant de répondre aux questions, apprenons la bonne démarche. Une bonne réponse reprend les mots de la question et forme une phrase complète, avec le verbe conjugué en position 2.

\begin{methode}[Répondre à une question de compréhension en quatre étapes]
\begin{enumerate}
\item \textbf{Lire la question} et identifier le mot interrogatif : \all{wann} (quand → une date), \all{wo} (où → un lieu), \all{was} (quoi → une action ou une chose), \all{wer} (qui → une personne), \all{wie} (comment → une manière).
\item \textbf{Repérer le mot-clé} de la question (par exemple \all{Weihnachten}, \all{die Kinder}, \all{die Ostereier}).
\item \textbf{Retrouver la phrase du texte} qui contient ce mot-clé et l'information demandée.
\item \textbf{Reformuler la réponse} en phrase complète, en reprenant les mots de la question : \all{Wann feiert man Weihnachten?} → \all{Man feiert Weihnachten am 25. Dezember.}
\end{enumerate}
\end{methode}

\begin{popfigure}
\begin{tikzpicture}[
  node distance=0.55cm,
  etape/.style={rectangle, rounded corners=3pt, draw=popblue, fill=popblueL, text=popdark, align=center, minimum width=4.6cm, minimum height=0.9cm, font=\small},
  fleche/.style={-{Stealth[length=2.5mm]}, thick, poporange}
]
\node[etape, align=center] (e1){\textbf{1. Lire la question}\\ (mot interrogatif : \all{wann}, \all{wo}, \all{was})};
\node[etape, below=of e1, align=center] (e2){\textbf{2. Repérer le mot-clé}\\ (\all{Weihnachten}, \all{die Kinder}, \dots)};
\node[etape, below=of e2, align=center] (e3){\textbf{3. Retrouver la phrase}\\ dans le texte};
\node[etape, below=of e3, draw=poppurple, fill=poppurpleL, align=center] (e4){\textbf{4. Reformuler la réponse}\\ en phrase complète (verbe en position 2)};
\draw[fleche] (e1) -- (e2);
\draw[fleche] (e2) -- (e3);
\draw[fleche] (e3) -- (e4);
\end{tikzpicture}
\legende{Organigramme de la méthode de compréhension : de la question à la réponse complète.}
\end{popfigure}

\begin{exemplebox}[Question et réponse modèles]
\all{Was machen die Kinder am Weihnachtsmorgen?} « Que font les enfants le matin de Noël ? »

Réponse complète : \all{Die Kinder stehen früh auf und öffnen die Geschenke.} « Les enfants se lèvent tôt et ouvrent les cadeaux. »

On observe : le mot interrogatif \all{was} demande une \textbf{action} ; la réponse reprend le sujet \all{die Kinder} et donne les deux actions du texte (\all{aufstehen}, \all{öffnen}).
\end{exemplebox}

\begin{attention}[L'inversion du sujet dans la réponse]
En allemand, on commence souvent la réponse par l'élément demandé (la date, le lieu). Dans ce cas, \textbf{le verbe reste en position 2} et le sujet passe après le verbe : c'est l'inversion.

\all{Am 25. Dezember feiert man Weihnachten.} « Le 25 décembre, on fête Noël. »

Ne dites jamais : \emph{Am 25. Dezember man feiert Weihnachten.} Le verbe conjugué doit toujours être le deuxième élément de la phrase. En français, l'ordre ne change pas (« on fête Noël le 25 décembre » / « le 25 décembre, on fête Noël ») ; en allemand, si, et c'est une erreur classique des francophones.
\end{attention}

\courssub{Compréhension globale du texte}

Les questions globales portent sur le sens général du texte, sans entrer dans les détails. Elles commencent souvent par \all{wovon} (de quoi), \all{wer} (qui) ou \all{welche} (quels).

\begin{exercice}[Questions de compréhension globale]
Répondez par des phrases complètes en allemand.
\begin{enumerate}
\item \all{Wovon handelt der Text?} (De quoi parle le texte ?)
\item \all{Wer spricht im Text?} (Qui parle dans le texte ?)
\item \all{Welche Feste werden erwähnt?} (Quelles fêtes sont mentionnées ?)
\item \all{Sind diese Feste religiös oder national?} (Ces fêtes sont-elles religieuses ou nationales ?)
\end{enumerate}
\end{exercice}

\begin{corrige}[Compréhension globale]
\begin{enumerate}
\item \all{Der Text handelt von den Festen in Deutschland.} « Le texte parle des fêtes en Allemagne. »
\item \all{Ein deutscher Schüler spricht und beschreibt die Feste seines Landes.} « Un élève allemand parle et décrit les fêtes de son pays. »
\item \all{Der Text erwähnt Weihnachten, Ostern und den Nationalfeiertag.} « Le texte mentionne Noël, Pâques et la fête nationale. »
\item \all{Weihnachten und Ostern sind religiöse Feste, der Nationalfeiertag ist ein nationales Fest.} « Noël et Pâques sont des fêtes religieuses, la fête nationale est une fête nationale. »
\end{enumerate}
Remarquez la question 3 : \all{erwähnen} (mentionner) est suivi de l'accusatif, d'où \all{den Nationalfeiertag} (masculin à l'accusatif).
\end{corrige}

\courssub{Compréhension détaillée du texte}

Les questions détaillées portent sur des informations précises. Le mot interrogatif vous indique le type d'information à chercher.

\begin{center}
\begin{tabularx}{\linewidth}{|l|X|X|}
\hline
\rowcolor{popblueL} \textbf{Mot interrogatif} & \textbf{Information demandée} & \textbf{Exemple de réponse} \\
\hline
\all{wann} (quand) & une date, un moment & \all{Man feiert Weihnachten am 25. Dezember.} \\
\hline
\all{wo} (où) & un lieu & \all{Man sucht die Ostereier im Garten.} \\
\hline
\all{was} (quoi) & une action, une chose & \all{Die Kinder öffnen die Geschenke.} \\
\hline
\all{wer} (qui) & une personne & \all{Die Kinder stehen früh auf.} \\
\hline
\all{wie} (comment) & une manière & \all{Die Familie feiert zusammen.} \\
\hline
\end{tabularx}
\end{center}

\begin{exoresolu}[Questions de compréhension détaillée]
Énoncé : répondez par des phrases complètes.
\begin{enumerate}
\item \all{Wann feiert man Weihnachten?}
\item \all{Was machen die Kinder am Weihnachtsmorgen?}
\item \all{Wo sucht man die Ostereier?}
\item \all{Was gibt es am Nationalfeiertag?}
\end{enumerate}
\tcblower
Solution détaillée :
\begin{enumerate}
\item \all{Wann} → une date. \all{Man feiert Weihnachten am 25. Dezember.} « On fête Noël le 25 décembre. » (On peut aussi dire, avec inversion : \all{Am 25. Dezember feiert man Weihnachten.})
\item \all{Was} → des actions. \all{Die Kinder stehen früh auf und öffnen die Geschenke.} « Les enfants se lèvent tôt et ouvrent les cadeaux. »
\item \all{Wo} → un lieu. \all{Man sucht die Ostereier im Garten.} « On cherche les œufs de Pâques dans le jardin. »
\item \all{Was gibt es} → ce qu'il y a. \all{Am Nationalfeiertag gibt es eine Parade und viele Fahnen.} « À la fête nationale, il y a un défilé et beaucoup de drapeaux. »
\end{enumerate}
\end{exoresolu}

\begin{attention}[Répondre avec \all{es gibt}]
La question \all{Was gibt es am Nationalfeiertag?} utilise l'expression \all{es gibt} (« il y a »), toujours suivie de l'\textbf{accusatif}. La réponse doit reprendre la même structure : \all{Es gibt eine Parade und viele Fahnen.} Ne traduisez jamais « il y a » par \emph{es ist} ou \emph{es hat} : c'est un faux ami de structure avec le français.
\end{attention}

\courssub{Vrai ou faux : justifier par une citation du texte}

Pour les questions vrai/faux, il ne suffit pas de répondre \all{richtig} ou \all{falsch} : il faut \textbf{justifier} en citant la phrase du texte qui prouve votre réponse. C'est la méthode attendue à l'examen.

\begin{exercice}[Richtig oder falsch ?]
Répondez par \all{richtig} ou \all{falsch} et justifiez avec une phrase du texte.
\begin{enumerate}
\item \all{Man feiert Weihnachten im Sommer.}
\item \all{Die Kinder öffnen am Weihnachtsmorgen die Geschenke.}
\item \all{Die Ostereier sucht man in der Schule.}
\item \all{Am Nationalfeiertag gibt es eine Parade.}
\end{enumerate}
\end{exercice}

\begin{corrige}[Richtig oder falsch ?]
\begin{enumerate}
\item \all{Falsch. Im Text steht: „Man feiert Weihnachten am 25. Dezember.“} — Le 25 décembre est en hiver, pas en été.
\item \all{Richtig. Im Text steht: „Die Kinder stehen früh auf und öffnen die Geschenke.“}
\item \all{Falsch. Im Text steht: „Man sucht die Ostereier im Garten.“} — C'est dans le jardin, pas à l'école.
\item \all{Richtig. Im Text steht: „Am Nationalfeiertag gibt es eine Parade und viele Fahnen.“}
\end{enumerate}
Remarquez la formule de justification : \all{Im Text steht: „…“} (« Dans le texte, il est écrit : … »). Apprenez-la par cœur.
\end{corrige}

\courssub{Recherche dans le texte : dates et verbes séparables}

Cette activité entraîne l'œil à repérer deux catégories grammaticales du programme : les \textbf{dates} (nombres ordinaux) et les \textbf{verbes à particule séparable}. Nous les étudierons en détail dans les sections de grammaire ; ici, il s'agit seulement de les \emph{repérer}.

\begin{exemplebox}[Ce qu'on relève dans le texte]
\textbf{Les dates} (avec un ordinal + la préposition \all{am}) :
\begin{itemize}
\item \all{am 25. Dezember} — le 25 décembre (on lit : \all{am fünfundzwanzigsten Dezember})
\item \all{am 3. Oktober} — le 3 octobre, fête nationale allemande (\all{am dritten Oktober})
\item \all{am 20. Mai} — le 20 mai, fête nationale camerounaise (\all{am zwanzigsten Mai})
\end{itemize}
\textbf{Les verbes séparables} (la particule va en fin de phrase au Präsens) :
\begin{itemize}
\item \all{aufstehen} : \all{Die Kinder stehen früh auf.} — se lever
\item \all{stattfinden} : \all{Die Parade findet am Morgen statt.} — avoir lieu
\item \all{anfangen} : \all{Das Fest fängt um 10 Uhr an.} — commencer
\end{itemize}
\end{exemplebox}

\begin{propriete}[Comment reconnaître un verbe séparable dans un texte]
Un verbe séparable se reconnaît à deux signes :
\begin{enumerate}
\item À l'infinitif, il commence par une petite particule accentuée : \all{\textbf{auf}stehen}, \all{\textbf{statt}finden}, \all{\textbf{an}fangen}, \all{\textbf{ein}laden}.
\item Dans la phrase conjuguée au Präsens, la particule est \textbf{séparée} et se place \textbf{à la fin} : \all{Die Kinder stehen früh \textbf{auf}.}
\end{enumerate}
Astuce de lecture : quand vous voyez une petite particule (\all{auf}, \all{an}, \all{ein}, \all{statt}, \all{mit}) isolée en fin de phrase, cherchez le verbe conjugué en position 2 : ils forment un seul verbe.
\end{propriete}

\begin{exoresolu}[Relever et classer]
Énoncé : dans les phrases suivantes tirées du texte, soulignez la date et encadrez le verbe séparable (verbe + particule).
\begin{enumerate}
\item \all{Am 25. Dezember stehen die Kinder früh auf.}
\item \all{Am 3. Oktober findet in Berlin eine große Parade statt.}
\item \all{Am 20. Mai fängt das Fest in Yaoundé um 10 Uhr an.}
\end{enumerate}
\tcblower
Solution détaillée :
\begin{enumerate}
\item Date : \all{am 25. Dezember} ; verbe séparable : \all{aufstehen} → \all{stehen … auf}.
\item Date : \all{am 3. Oktober} ; verbe séparable : \all{stattfinden} → \all{findet … statt}.
\item Date : \all{am 20. Mai} ; verbe séparable : \all{anfangen} → \all{fängt … an}.
\end{enumerate}
On remarque que la date en tête de phrase provoque l'inversion du sujet (\all{stehen die Kinder}), et que la particule reste bien à la fin.
\end{exoresolu}

\courssub{Entraînement oral : mini-dialogue sur les fêtes}

Lisons maintenant ce mini-dialogue à deux voix, en classe. Il réutilise le vocabulaire du texte et prépare l'expression orale.

\begin{textebox}[Mini-dialogue : Feierst du Weihnachten?]
Amina: Feierst du Weihnachten, Lukas?

Lukas: Ja, ich feiere Weihnachten mit meiner Familie. Wir stehen früh auf und öffnen die Geschenke. Und du?

Amina: Ich feiere Ostern und den Nationalfeiertag am 20. Mai.

Lukas: Was gibt es am Nationalfeiertag in Kamerun?

Amina: Es gibt eine große Parade und viele Fahnen. Die Parade findet am Morgen statt.

Lukas: Das ist schön! In Deutschland ist der Nationalfeiertag am 3. Oktober.
\end{textebox}

\textbf{Glossaire :} \all{feiern} (fêter) ; \all{das Geschenk, -e} (le cadeau) ; \all{die Parade, -n} (le défilé) ; \all{die Fahne, -n} (le drapeau) ; \all{stattfinden} (a lieu, \all{findet statt}, \all{hat stattgefunden}) ; \all{am Morgen} (le matin).

Questions orales sur le dialogue :
\begin{enumerate}
\item \all{Wann feiert Amina den Nationalfeiertag?} → \all{Sie feiert den Nationalfeiertag am 20. Mai.}
\item \all{Was gibt es am Nationalfeiertag in Kamerun?} → \all{Es gibt eine große Parade und viele Fahnen.}
\item \all{Wann ist der Nationalfeiertag in Deutschland?} → \all{Er ist am 3. Oktober.}
\end{enumerate}

\courssub{Question d'ouverture : les fêtes en Allemagne et au Cameroun}

La question d'ouverture invite à comparer les cultures : c'est le cœur de l'approche par les compétences. Elle prépare aussi le commentaire de la section 6.

\textbf{Question :} \all{Welches deutsche Fest ist wie ein Fest in Kamerun?} « Quelle fête allemande ressemble à une fête camerounaise ? »

\begin{exemplebox}[Modèles de réponses]
\begin{itemize}
\item \all{Weihnachten in Deutschland ist wie Weihnachten in Kamerun: Die Familie feiert zusammen und geht zum Gottesdienst.} « Noël en Allemagne ressemble à Noël au Cameroun : la famille fête ensemble et va au service religieux. »
\item \all{Der Nationalfeiertag am 3. Oktober ist wie der 20. Mai in Kamerun: Es gibt eine Parade und viele Fahnen.} « La fête nationale du 3 octobre ressemble au 20 mai au Cameroun : il y a un défilé et beaucoup de drapeaux. »
\end{itemize}
\end{exemplebox}

\begin{center}
\begin{tabularx}{\linewidth}{|X|X|X|}
\hline
\rowcolor{popblueL} \textbf{Deutschland} & \textbf{Kamerun} & \textbf{Point commun} \\
\hline
\all{Weihnachten (25. Dezember)} & \all{Weihnachten (25. Dezember)} & \all{der Gottesdienst}, la famille réunie \\
\hline
\all{Ostern} & \all{Ostern} & fête religieuse, repas en famille \\
\hline
\all{der Nationalfeiertag (3. Oktober)} & \all{der Nationalfeiertag (20. Mai)} & \all{die Parade}, \all{die Fahnen} \\
\hline
\end{tabularx}
\end{center}

\begin{experience}[Débat en binômes]
En binômes, posez-vous la question d'ouverture et répondez en trois phrases minimum, avec les structures du cours : \all{Ich feiere …}, \all{Es gibt …}, \all{Das Fest findet … statt}. Puis deux binômes présentent leur réponse devant la classe. Critères : phrase complète, verbe en position 2, particule séparable en fin de phrase.
\end{experience}

\begin{savaistu}[Le 3 octobre et le 20 mai]
La fête nationale allemande, le \all{Tag der Deutschen Einheit} (« Jour de l'unité allemande »), célèbre la réunification de l'Allemagne le 3 octobre 1990, après la chute du Mur de Berlin. Au Cameroun, le 20 mai commémore le référendum de 1972 qui a fait du pays une République unie : ce jour-là, les élèves des lycées de Yaoundé et de Douala défilent devant les autorités, comme les soldats lors de la parade allemande. Deux pays, deux histoires, mais le même symbole : \all{die Parade} et \all{die Fahnen}.
\end{savaistu}

\begin{aretenir}[L'essentiel de la compréhension de texte]
\begin{itemize}
\item Identifier le mot interrogatif : \all{wann} → date, \all{wo} → lieu, \all{was} → action, \all{wer} → personne.
\item Répondre toujours par une \textbf{phrase complète} qui reprend les mots de la question.
\item Si la réponse commence par la date ou le lieu, le verbe reste en position 2 : \all{Am 25. Dezember feiert man Weihnachten.}
\item Pour un vrai/faux, justifier avec la formule : \all{Im Text steht: „…“}
\item Repérer les verbes séparables : particule en fin de phrase (\all{stehen … auf}, \all{findet … statt}).
\end{itemize}
\end{aretenir}

\coursec{Lexique thématique : les fêtes}

Dans la section précédente, nous avons lu et compris le texte \all{Feste in Deutschland} : nous savons maintenant \emph{de quoi} il parle et \emph{ce qu'il dit}. Mais pour pouvoir parler nous-mêmes des fêtes — à l'oral comme à l'écrit — il faut posséder les mots exacts, avec leur \cle{article} et leur \cle{pluriel}. C'est l'objectif de cette section : construire, champ par champ, un vocabulaire solide et actif sur les fêtes religieuses et nationales. En allemand, un nom appris sans son article est un nom à moitié appris : retenez donc toujours le trio \cle{article — nom — pluriel}, par exemple \all{die Fahne, die Fahnen}.

\courssub{Observer d'abord : un petit texte pour découvrir les mots}

Avant tout tableau, lisons ce court texte original. Soulignez mentalement tous les mots qui se rapportent à une fête.

\begin{textebox}[Festtage in Deutschland und in Kamerun]
In Deutschland gibt es viele Feste. Im Dezember feiern die Familien Weihnachten. Die Kinder schmücken den Baum, und am Abend gehen viele Menschen zum Gottesdienst in die Kirche. Dort brennen Kerzen, und die Familie betet zusammen. Am Heiligabend schenken die Eltern den Kindern Geschenke. Im Frühjahr feiern die Christen Ostern.

Aber es gibt auch nationale Feste. Am Nationalfeiertag hängen überall Fahnen an den Häusern. Am Morgen findet eine große Parade statt, und der Präsident hält eine Rede. Viele Bürger nehmen an der Parade teil. Am Abend sehen die Leute ein Feuerwerk.

In Kamerun ist es ähnlich. Die Schüler in Yaoundé stehen am Feiertag früh auf und machen bei der Parade mit. Die Lehrer laden die Eltern ein, und alle feiern zusammen. Feste bringen die Menschen zusammen: man singt, man tanzt, man isst und man ist froh.
\end{textebox}

\textbf{Glossaire} : \all{der Gottesdienst} = l'office religieux ; \all{brennen} = brûler ; \all{beten} = prier ; \all{der Heiligabend} = la veille de Noël (24 décembre) ; \all{das Frühjahr} = le printemps ; \all{hängen} = être accroché ; \all{halten (eine Rede)} = prononcer (un discours) ; \all{der Bürger} = le citoyen ; \all{teilnehmen an} = participer à ; \all{ähnlich} = semblable ; \all{zusammenbringen} = rassembler.

\textbf{Questions de compréhension rapide} :
\begin{enumerate}
\item \all{Was machen die Kinder im Dezember?}
\item \all{Wohin gehen viele Menschen am Abend?}
\item \all{Was findet am Nationalfeiertag am Morgen statt?}
\item \all{Wer hält eine Rede?}
\end{enumerate}

\textbf{Réponses attendues} : 1. \all{Die Kinder schmücken den Baum.} 2. \all{Sie gehen zum Gottesdienst in die Kirche.} 3. \all{Am Morgen findet eine große Parade statt.} 4. \all{Der Präsident hält eine Rede.}

Vous avez repéré les mots-clés ? Classons-les maintenant en trois champs lexicaux.

\courssub{Carte mentale du vocabulaire}

\begin{popfigure}
\begin{tikzpicture}[mindmap, grow cyclic, every node/.style={concept, align=center, font=\small},
root concept/.append style={concept color=popblue, font=\bfseries},
level 1/.append style={level distance=3.4cm, sibling angle=120},
level 2/.append style={level distance=2.4cm, sibling angle=45}]
\node[root concept]{die Feste}
  child[concept color=popgreen]{ node {religiöse Feste}
    child { node[concept color=popgreenL, text=popdark] {Weihnachten Ostern} }
    child { node[concept color=popgreenL, text=popdark] {der Gottesdienst die Kirche} }
    child { node[concept color=popgreenL, text=popdark] {die Kerze beten} }
  }
  child[concept color=poporange]{ node {nationale Feste}
    child { node[concept color=poporangeL, text=popdark] {der Nationalfeiertag der Feiertag} }
    child { node[concept color=poporangeL, text=popdark] {die Parade die Fahne} }
    child { node[concept color=poporangeL, text=popdark] {das Feuerwerk die Rede} }
  }
  child[concept color=poppurple]{ node {Aktivitäten}
    child { node[concept color=poppurpleL, text=popdark] {feiern schmücken} }
    child { node[concept color=poppurpleL, text=popdark] {schenken einladen} }
    child { node[concept color=poppurpleL, text=popdark] {stattfinden mitmachen} }
  };
\end{tikzpicture}
\legende{Carte mentale : les trois champs lexicaux des fêtes}
\end{popfigure}

\courssub{Champ 1 — Les fêtes religieuses}

\begin{definition}[Les mots du culte]
\begin{itemize}
\item \all{Weihnachten} (pluriel : --, fête sans article) = Noël. On dit \all{zu Weihnachten} ou \all{an Weihnachten} = à Noël.
\item \all{Ostern} (pluriel : --, fête sans article) = Pâques. On dit \all{zu Ostern} = à Pâques.
\item \all{der Gottesdienst, die Gottesdienste} = l'office religieux, la messe.
\item \all{die Kirche, die Kirchen} = l'église.
\item \all{die Kerze, die Kerzen} = la bougie.
\item \all{beten} (verbe régulier) = prier : \all{ich bete, du betest, er betet}.
\end{itemize}
\end{definition}

Observez ces phrases tirées de la vie quotidienne :

\begin{exemplebox}[Les fêtes religieuses en contexte]
\all{Wir gehen zum Gottesdienst.} « Nous allons à l'office. » (\all{zum} = \all{zu dem})\par
\all{Die Familie betet vor dem Essen.} « La famille prie avant le repas. »\par
\all{In der Kirche brennen viele Kerzen.} « Dans l'église, beaucoup de bougies brûlent. »\par
\all{An Weihnachten sind die Kirchen voll.} « À Noël, les églises sont pleines. »
\end{exemplebox}

\begin{propriete}[Weihnachten et Ostern : des fêtes sans article]
Les noms de fêtes \all{Weihnachten} et \all{Ostern} s'emploient en général \cle{sans article} : \all{Weihnachten ist im Dezember.} « Noël est en décembre. »\par
Avec les prépositions, on dit : \all{zu Weihnachten} ou \all{an Weihnachten} (à Noël), \all{zu Ostern} (à Pâques).\par
Exception : si l'on ajoute un adjectif ou une détermination, le nom devient \cle{neutre} et prend l'article : \all{das letzte Weihnachten} « le Noël dernier », \all{das nächste Ostern} « le prochain Pâques ».
\end{propriete}

\courssub{Champ 2 — Les fêtes nationales}

\begin{definition}[Les mots de la fête nationale]
\begin{itemize}
\item \all{der Nationalfeiertag, die Nationalfeiertage} = la fête nationale.
\item \all{die Parade, die Paraden} = le défilé.
\item \all{die Fahne, die Fahnen} = le drapeau.
\item \all{das Feuerwerk, die Feuerwerke} = le feu d'artifice.
\item \all{die Rede, die Reden} = le discours.
\item \all{der Feiertag, die Feiertage} = le jour férié (jour où l'on ne travaille pas).
\end{itemize}
\end{definition}

\begin{exemplebox}[La fête nationale en contexte]
\all{Die Parade findet am Morgen statt.} « Le défilé a lieu le matin. »\par
\all{Die Fahnen hängen an den Häusern.} « Les drapeaux sont accrochés aux maisons. »\par
\all{Der Präsident hält eine Rede.} « Le président prononce un discours. »\par
\all{Am Abend gibt es ein Feuerwerk.} « Le soir, il y a un feu d'artifice. »
\end{exemplebox}

\begin{attention}[Faux amis et nuances]
\all{die Parade} existe en français, mais en allemand le \all{-e} final se prononce : « pa-ra-deu » (trois syllabes).\par
Ne confondez pas \all{der Feiertag} (le jour \emph{férié}, officiel, où les magasins sont fermés) et \all{der Festtag} (le jour de \emph{fête}, que l'on célèbre en famille). Cette nuance est souvent demandée au Bac : \all{Der 1. Januar ist ein Feiertag.} « Le 1er janvier est un jour férié. »\par
Autre piège : \all{die Rede} = le discours, mais \all{reden} = parler ; \all{die Kirche} = l'église (le bâtiment et l'institution), jamais « la messe » : la messe, c'est \all{der Gottesdienst}.
\end{attention}

\courssub{Champ 3 — Les actions de la fête}

\begin{definition}[Les verbes et noms de l'activité festive]
\begin{itemize}
\item \all{feiern} = fêter, célébrer : \all{Wir feiern ein Fest.} « Nous fêtons une fête. »
\item \all{schmücken} = décorer : \all{Die Kinder schmücken den Baum.} « Les enfants décorent le sapin. »
\item \all{schenken} = offrir (faire un cadeau) : \all{Die Eltern schenken den Kindern Spielzeug.} « Les parents offrent des jouets aux enfants. »
\item \all{das Geschenk, die Geschenke} = le cadeau.
\item \all{einladen} (verbe séparable) = inviter : \all{Wir laden die Nachbarn ein.} « Nous invitons les voisins. »
\item \all{stattfinden} (verbe séparable) = avoir lieu : \all{Die Parade findet statt.} « Le défilé a lieu. »
\item \all{mitmachen} (verbe séparable) = participer, être de la partie : \all{Alle machen mit.} « Tout le monde participe. »
\item \all{aufstehen} (verbe séparable) = se lever : \all{Am Feiertag stehen wir früh auf.} « Le jour férié, nous nous levons tôt. »
\end{itemize}
\end{definition}

Remarquez dans les exemples ci-dessus la place de la particule (\all{ein}, \all{statt}, \all{mit}, \all{auf}) : elle va \cle{à la fin de la phrase}. Ce point sera étudié en détail dans la section Grammaire 2 ; observez-le dès maintenant dans le vocabulaire.

\courssub{Les familles de mots : un mot en cache trois autres}

L'allemand construit ses mots comme des briques : à partir d'une racine, on forme des noms, des adjectifs, des verbes. Apprendre une famille, c'est apprendre quatre mots pour le prix d'un.

\begin{center}
\begin{tabularx}{\linewidth}{|X|X|X|}
\hline
\rowcolor{popblueL} Racine & Famille & Exemple traduit \\
\hline
\all{feiern} (fêter) & \all{die Feier} (la fête, la cérémonie) ; \all{der Feiertag} (le jour férié) & \all{Die Feier beginnt um 10 Uhr.} « La cérémonie commence à 10 heures. » \\
\hline
\all{die Fahne} (le drapeau) & \all{der Fahnenmast} (le mât à drapeau) & \all{Die Fahne hängt am Fahnenmast.} « Le drapeau est accroché au mât. » \\
\hline
\all{das Fest} (la fête) & \all{festlich} (festif, de fête) & \all{Die Stadt ist festlich geschmückt.} « La ville est décorée de façon festive. » \\
\hline
\all{schenken} (offrir) & \all{das Geschenk} (le cadeau) & \all{Das Geschenk ist für dich.} « Le cadeau est pour toi. » \\
\hline
\end{tabularx}
\end{center}

\begin{aretenir}[Stratégie de vocabulaire]
Quand vous rencontrez un mot nouveau, demandez-vous : quelle est sa racine ? Quels autres mots de la famille connaissez-vous déjà ? \all{feiern — die Feier — der Feiertag — das Fest — festlich} forment un seul réseau dans la mémoire.
\end{aretenir}

\courssub{Tableau récapitulatif : articles et pluriels}

\begin{center}
\begin{tabularx}{\linewidth}{|l|X|l|}
\hline
\rowcolor{popblueL} Nom allemand & Pluriel & Français \\
\hline
\all{der Gottesdienst} & \all{die Gottesdienste} & l'office religieux \\
\hline
\all{die Kirche} & \all{die Kirchen} & l'église \\
\hline
\all{die Kerze} & \all{die Kerzen} & la bougie \\
\hline
\all{der Nationalfeiertag} & \all{die Nationalfeiertage} & la fête nationale \\
\hline
\all{der Feiertag} & \all{die Feiertage} & le jour férié \\
\hline
\all{die Parade} & \all{die Paraden} & le défilé \\
\hline
\all{die Fahne} & \all{die Fahnen} & le drapeau \\
\hline
\all{das Feuerwerk} & \all{die Feuerwerke} & le feu d'artifice \\
\hline
\all{die Rede} & \all{die Reden} & le discours \\
\hline
\all{das Geschenk} & \all{die Geschenke} & le cadeau \\
\hline
\all{das Fest} & \all{die Feste} & la fête \\
\hline
\all{die Feier} & \all{die Feiern} & la cérémonie, la fête \\
\hline
\all{Weihnachten} / \all{Ostern} & -- (sans article) & Noël / Pâques \\
\hline
\end{tabularx}
\end{center}

\begin{propriete}[Lire un pluriel allemand]
Dans ce tableau, observez les marques de pluriel : beaucoup de noms féminins en \all{-e} prennent \all{-n} (\all{die Kerze → die Kerzen}, \all{die Fahne → die Fahnen}, \all{die Parade → die Paraden}) ; les noms masculins composés en \all{-tag} prennent \all{-e} (\all{der Feiertag → die Feiertage}). Apprenez toujours le pluriel en même temps que le mot.
\end{propriete}

\courssub{Expressions utiles pour parler des fêtes}

\begin{exemplebox}[À savoir par cœur]
\all{Frohe Weihnachten!} « Joyeux Noël ! »\par
\all{Frohe Ostern!} « Joyeuses Pâques ! »\par
\all{ein Fest feiern} « fêter une fête » : \all{Wir feiern ein Fest mit der Familie.} « Nous fêtons une fête en famille. »\par
\all{an einer Parade teilnehmen} « participer à un défilé » : \all{Die Soldaten nehmen an der Parade teil.} « Les soldats participent au défilé. »\par
\all{jemanden einladen} « inviter quelqu'un » : \all{Ich lade meine Freunde ein.} « J'invite mes amis. »
\end{exemplebox}

\begin{attention}[Erreurs fréquentes des francophones]
\begin{itemize}
\item Ne dites pas \all{die Weihnachten} : la fête s'emploie sans article (\all{Weihnachten kommt bald.}).
\item \all{schenken} se construit avec l'accusatif (la chose) et le datif (la personne) : \all{Ich schenke meiner Mutter Blumen.} « J'offre des fleurs à ma mère. »
\item N'oubliez pas la majuscule à tous les noms : \all{die Parade}, \all{das Feuerwerk}, \all{der Feiertag}.
\item \all{stattfinden} est séparable : on écrit \all{Die Parade findet am Morgen statt.}, jamais \all{stattfindet} en tête de phrase.
\end{itemize}
\end{attention}

\begin{savaistu}[Noël en Allemagne : le 24 au soir]
En Allemagne, Noël se fête surtout le 24 décembre au soir, le \all{Heiligabend} (littéralement « le saint soir »). Après le repas, on distribue les cadeaux : c'est la \all{Bescherung}. Dans le sud du pays (Bavière, Autriche), ce n'est pas le Père Noël mais le \all{Christkind} (l'enfant Jésus) qui apporte les cadeaux aux enfants. Et le 6 décembre, \all{der Nikolaus} (saint Nicolas) dépose des sucreries dans les chaussures des enfants sages !
\end{savaistu}

\begin{savaistu}[Deux fêtes nationales, deux histoires]
La fête nationale de l'Allemagne est le \all{Tag der Deutschen Einheit}, le 3 octobre : elle rappelle la réunification du pays en 1990. Au Cameroun, la fête nationale est le 20 mai, avec son célèbre défilé du boulevard du 20 Mai à Yaoundé : élèves, élèves-officiers et forces armées y défilent devant le président. Comparer les deux fêtes est un excellent sujet de production écrite au Bac !
\end{savaistu}

\begin{exoresolu}[Compléter avec le bon mot et le bon article]
Énoncé — Complétez avec : \all{die Fahne}, \all{der Gottesdienst}, \all{das Geschenk}, \all{die Parade}, \all{die Kerzen}.\par
1. \all{Am Nationalfeiertag findet \dots\ statt.}\par
2. \all{Am Sonntag gehen wir zu \dots\ .}\par
3. \all{In der Kirche brennen \dots\ .}\par
4. \all{Vor dem Haus hängt \dots\ .}\par
5. \all{Für meine Schwester kaufe ich \dots\ .}
\tcblower
Solution détaillée :\par
1. \all{Am Nationalfeiertag findet die Parade statt.} (accusatif féminin : \all{die} ne change pas ; \all{stattfinden} = avoir lieu, particule \all{statt} en fin de phrase.)\par
2. \all{Am Sonntag gehen wir zum Gottesdienst.} (\all{zu dem} se contracte en \all{zum} ; datif masculin après \all{zu}.)\par
3. \all{In der Kirche brennen die Kerzen.} (sujet au pluriel : \all{die Kerze → die Kerzen}.)\par
4. \all{Vor dem Haus hängt die Fahne.} (sujet féminin singulier au nominatif.)\par
5. \all{Für meine Schwester kaufe ich ein Geschenk.} (\all{das Geschenk} est neutre : accusatif \all{ein}.)
\end{exoresolu}

\begin{exercice}[À vous de jouer]
1. Traduisez en allemand : a) « Joyeux Noël ! » b) « Le défilé a lieu le matin. » c) « Nous invitons nos voisins. » d) « Les enfants décorent le sapin. »\par
2. Donnez l'article et le pluriel de : \all{Fahne}, \all{Feuerwerk}, \all{Geschenk}, \all{Rede}, \all{Feiertag}.\par
3. Rédigez trois phrases sur la fête nationale au Cameroun avec : \all{aufstehen}, \all{mitmachen}, \all{die Fahne}.
\end{exercice}

\begin{corrige}[Exercice « À vous de jouer »]
1. a) \all{Frohe Weihnachten!} b) \all{Die Parade findet am Morgen statt.} c) \all{Wir laden unsere Nachbarn ein.} d) \all{Die Kinder schmücken den Baum.}\par
2. \all{die Fahne, die Fahnen} ; \all{das Feuerwerk, die Feuerwerke} ; \all{das Geschenk, die Geschenke} ; \all{die Rede, die Reden} ; \all{der Feiertag, die Feiertage}.\par
3. Exemple : \all{Am Nationalfeiertag stehen wir früh auf. Die Schüler machen bei der Parade mit. Überall hängen Fahnen.}
\end{corrige}

\begin{aretenir}[L'essentiel du lexique]
Trois champs à maîtriser : les fêtes religieuses (\all{Weihnachten, Ostern, der Gottesdienst, die Kirche, die Kerze, beten}), les fêtes nationales (\all{der Nationalfeiertag, die Parade, die Fahne, das Feuerwerk, die Rede, der Feiertag}) et les actions (\all{feiern, schmücken, schenken, einladen, stattfinden, mitmachen, aufstehen}). Chaque nom s'apprend avec son article et son pluriel ; \all{Weihnachten} et \all{Ostern} s'emploient sans article. Ce vocabulaire sera réutilisé dans les sections de grammaire (dates, verbes séparables) et dans votre production écrite finale.
\end{aretenir}

\coursec{Grammaire 1 : les dates et les nombres ordinaux}

Dans la section précédente, nous avons découvert le vocabulaire des fêtes : \all{Weihnachten}, \all{Ostern}, \all{der Nationalfeiertag}, \all{die Parade}, \all{die Fahne}, \all{der Gottesdienst}. Mais une fête, c'est toujours une \cle{date} ! Pour dire que la fête nationale allemande a lieu « le 3 octobre » ou que le Cameroun célèbre sa fête nationale « le 20 mai », il faut savoir former les \cle{nombres ordinaux} et construire la date en allemand. C'est l'objet de cette section de grammaire.

\courssub{Observation : que remarque-t-on dans le texte ?}

Reprenons deux phrases du texte support \emph{Feste in Deutschland} :

\begin{exemplebox}[Deux dates du texte]
\all{Der Nationalfeiertag ist am 3. Oktober.} « La fête nationale est le 3 octobre. »\par
\all{Weihnachten ist am 25. Dezember.} « Noël est le 25 décembre. »
\end{exemplebox}

Observons attentivement :
\begin{itemize}
\item après le chiffre, il y a un \cle{point} : \all{3.}, \all{25.} ;
\item la date est introduite par la préposition \all{am} (= an + dem) pour répondre à la question \all{wann ?} « quand ? » ;
\item si on lit à voix haute, on ne dit pas \all{drei} ni \all{fünfundzwanzig}, mais \all{am dritten Oktober}, \all{am fünfundzwanzigsten Dezember} : le chiffre se prononce comme un \emph{adjectif ordinal} qui se termine en \all{-en}.
\end{itemize}

\begin{aretenir}[Le point = l'ordinal]
En allemand, le \cle{point} placé après un chiffre indique qu'il s'agit d'un \cle{nombre ordinal} : \all{3.} se lit \all{dritte} (« troisième »), \all{25.} se lit \all{fünfundzwanzigste} (« vingt-cinquième »). C'est l'équivalent de notre « 3\textsuperscript{e} » ou « 25\textsuperscript{e} » en français.
\end{aretenir}

\courssub{Formation des nombres ordinaux}

Un \cle{nombre ordinal} exprime un rang, un ordre : premier, deuxième, troisième\ldots{} En allemand, on le forme à partir du nombre cardinal (eins, zwei, drei\ldots) en ajoutant une terminaison.

\begin{propriete}[Formation des ordinaux]
\begin{itemize}
\item \textbf{De 1 à 19} : cardinal + \cle{-te} : \all{vier} $\rightarrow$ \all{der vierte} (4\textsuperscript{e}), \all{zehn} $\rightarrow$ \all{der zehnte} (10\textsuperscript{e}), \all{neunzehn} $\rightarrow$ \all{der neunzehnte} (19\textsuperscript{e}).
\item \textbf{À partir de 20} : cardinal + \cle{-ste} : \all{zwanzig} $\rightarrow$ \all{der zwanzigste} (20\textsuperscript{e}), \all{einundzwanzig} $\rightarrow$ \all{der einundzwanzigste} (21\textsuperscript{e}), \all{hundert} $\rightarrow$ \all{der hundertste} (100\textsuperscript{e}).
\item \textbf{Trois irréguliers à mémoriser} : \all{der erste} (1\textsuperscript{er}), \all{der dritte} (3\textsuperscript{e}), \all{der siebte} (7\textsuperscript{e}).
\end{itemize}
\end{propriete}

\begin{center}
\begin{tabular}{|l|l|l|}
\hline
\rowcolor{popblueL} Cardinal & Ordinal & Français \\
\hline
eins & der erste & le premier \\
\hline
zwei & der zweite & le deuxième \\
\hline
drei & der dritte & le troisième \\
\hline
vier & der vierte & le quatrième \\
\hline
sieben & der siebte & le septième \\
\hline
acht & der achte & le huitième \\
\hline
zwölf & der zwölfte & le douzième \\
\hline
neunzehn & der neunzehnte & le dix-neuvième \\
\hline
zwanzig & der zwanzigste & le vingtième \\
\hline
einundzwanzig & der einundzwanzigste & le vingt et unième \\
\hline
dreißig & der dreißigste & le trentième \\
\hline
einunddreißig & der einunddreißigste & le trente et unième \\
\hline
\end{tabular}
\end{center}

\begin{attention}[Petites irrégularités de forme]
Attention à deux formes légèrement modifiées : \all{drei} $\rightarrow$ \all{dritte} (et non « dreite ») et \all{sieben} $\rightarrow$ \all{siebte} (le \all{-en} tombe). Notez aussi \all{acht} $\rightarrow$ \all{achte} : le \all{t} du cardinal suffit, on n'écrit pas « achtte ».
\end{attention}

\begin{popfigure}
\begin{tikzpicture}[
  jour/.style={rectangle, rounded corners=2pt, draw=popblue, fill=popblueL, minimum width=1.7cm, minimum height=1.1cm, align=center, font=\small},
  fleche/.style={-{Stealth[length=2mm]}, thick, poporange}
]
\node[jour, align=center] (j1) at (0,0){\textbf{1.}\\ erste};
\node[jour, align=center] (j2) at (2.3,0){\textbf{2.}\\ zweite};
\node[jour, align=center] (j3) at (4.6,0){\textbf{3.}\\ dritte};
\node[jour, align=center] (j4) at (6.9,0){\textbf{7.}\\ siebte};
\node[jour, align=center] (j5) at (9.2,0){\textbf{19.}\\ neunzehnte};
\node[jour, draw=poporange, fill=poporangeL, align=center] (j6) at (11.5,0){\textbf{20.}\\ zwanzigste};
\draw[fleche] (j1) -- (j2);
\draw[fleche] (j2) -- (j3);
\draw[fleche] (j3) -- (j4);
\draw[fleche] (j4) -- (j5);
\draw[fleche] (j5) -- (j6);
\node[align=center, font=\small\itshape, popdark] at (4.6,-1.35) {cardinal + \textbf{-te} (de 1 à 19)};
\node[align=center, font=\small\itshape, popdark] at (11.5,-1.35) {cardinal + \textbf{-ste}\\ (à partir de 20)};
\draw[decorate, decoration={brace, amplitude=4pt, mirror}, popblue, thick] (-0.85,-0.7) -- (10.05,-0.7);
\draw[decorate, decoration={brace, amplitude=4pt, mirror}, poporange, thick] (10.65,-0.7) -- (12.35,-0.7);
\end{tikzpicture}
\legende{Frise du calendrier : la terminaison \all{-te} jusqu'au 19, la terminaison \all{-ste} à partir du 20.}
\end{popfigure}

\courssub{Dire et écrire la date}

\begin{propriete}[Poser la question et répondre]
Deux questions sont possibles :
\begin{itemize}
\item \all{Der wievielte ist heute?} « Le combien sommes-nous aujourd'hui ? » $\rightarrow$ réponse au \textbf{nominatif} : \all{Heute ist der zwanzigste Mai.} « Aujourd'hui, nous sommes le 20 mai. »
\item \all{Wann ist die Parade?} « Quand est le défilé ? » $\rightarrow$ réponse avec \cle{am} + ordinal au \textbf{datif} : \all{Die Parade ist am dritten Oktober.} « Le défilé est le 3 octobre. »
\end{itemize}
\end{propriete}

L'ordinal se comporte exactement comme un \cle{adjectif} : il prend la terminaison de l'adjectif après l'article défini. Après \all{der} (nominatif), la terminaison est \all{-e} : \all{der erste Mai}. Après \all{am} (= an + dem, datif), la terminaison est \all{-en} : \all{am ersten Mai}.

\begin{center}
\begin{tabularx}{\linewidth}{|X|X|X|}
\hline
\rowcolor{popblueL} Français & Deutsch & Remarque \\
\hline
le premier mai & der erste Mai & nominatif : -e \\
\hline
le 20 mai (quand ?) & am zwanzigsten Mai & datif après am : -en \\
\hline
le 3 octobre & am 3. Oktober & point obligatoire \\
\hline
la première parade & die erste Parade & accord avec un nom féminin \\
\hline
nous sommes le 25 décembre & heute ist der 25. Dezember & der + ordinal \\
\hline
\end{tabularx}
\end{center}

\begin{exemplebox}[Exemples traduits]
\all{Heute ist der erste Mai.} « Aujourd'hui c'est le premier mai. »\par
\all{Die Parade ist am dritten Oktober.} « Le défilé est le 3 octobre. »\par
\all{Am zwanzigsten Mai feiert Kamerun seinen Nationalfeiertag.} « Le 20 mai, le Cameroun fête sa fête nationale. »\par
\all{Der Gottesdienst beginnt am fünfundzwanzigsten Dezember um neun Uhr.} « L'office religieux commence le 25 décembre à neuf heures. »\par
\all{Ostern ist am ersten Sonntag nach dem Vollmond.} « Pâques est le premier dimanche après la pleine lune. »
\end{exemplebox}

\begin{methode}[Comment dire une date en quatre étapes]
\begin{enumerate}
\item Identifier la question : \all{der wievielte?} (quelle date sommes-nous ?) ou \all{wann?} (quand ?).
\item Former l'ordinal : cardinal + \all{-te} (1--19) ou \all{-ste} (20 et plus) ; vérifier les irréguliers \all{erste}, \all{dritte}, \all{siebte}.
\item Choisir la construction : \all{der wievielte?} $\rightarrow$ \all{der + ordinal + mois} ; \all{wann?} $\rightarrow$ \all{am + ordinal en -en + mois}.
\item À l'écrit, remplacer l'ordinal par le chiffre suivi d'un point : \all{am 3. Oktober}, \all{der 20. Mai}.
\end{enumerate}
\end{methode}

\courssub{L'écriture des dates et la lecture des années}

À l'écrit, les Allemands abrègent souvent la date en chiffres, dans l'ordre \textbf{jour.mois.année} :

\begin{exemplebox}[Dates abrégées]
\all{der 20.5.2025} se lit \all{der zwanzigste Mai zweitausendfünfundzwanzig} « le 20 mai 2025 ».\par
\all{am 1.1.2026} se lit \all{am ersten Januar zweitausendsechsundzwanzig} « le 1\textsuperscript{er} janvier 2026 ».
\end{exemplebox}

Pour les \cle{années}, deux lectures sont possibles :
\begin{itemize}
\item en un bloc : \all{2025} = \all{zweitausendfünfundzwanzig} ;
\item en deux blocs pour les années récentes : \all{1990} = \all{neunzehnhundertneunzig} (littéralement « dix-neuf cent quatre-vingt-dix »), exactement comme en français.
\end{itemize}

\begin{savaistu}[Comment les Allemands disent les années]
Comme les Français, les Allemands découpent les années récentes en deux blocs : \all{1990} se dit \all{neunzehnhundertneunzig}, \all{2000} se dit simplement \all{zweitausend}. La réunification allemande date du \all{3. Oktober 1990} (\all{am dritten Oktober neunzehnhundertneunzig}) : c'est pourquoi le 3 octobre est devenu la fête nationale allemande, le \all{Tag der Deutschen Einheit} (« Jour de l'Unité allemande »). Au Cameroun, la date clé est le \all{20. Mai 1972}, jour du référendum qui a donné naissance à la République unie du Cameroun.
\end{savaistu}

\courssub{Texte d'application : un calendrier de fêtes}

\begin{textebox}[Unsere Feste im Kalender]
Das Jahr hat viele Feste. Am ersten Januar feiert man Neujahr. Viele Familien stehen spät auf, denn die Nacht war lang. Im Frühling kommt Ostern: die Kinder suchen Ostereier im Garten, und viele Familien gehen am Sonntag zum Gottesdienst.\par
Am ersten Mai ist der Tag der Arbeit. In Deutschland ist das ein Feiertag, die Geschäfte sind geschlossen. Am zwanzigsten Mai feiert Kamerun seinen Nationalfeiertag: in Yaoundé und in Douala gibt es eine große Parade, die Schüler marschieren, und überall weht die Fahne. „Heute ist der zwanzigste Mai, unser Nationalfeiertag“, sagt die Lehrerin. „Wann ist der Nationalfeiertag in Deutschland?“, fragt ein Schüler. „Am dritten Oktober“, antwortet die Lehrerin. „An diesem Tag feiert Deutschland die Deutsche Einheit.“\par
Am fünfundzwanzigsten Dezember ist Weihnachten. Die Familien essen zusammen, die Kinder bekommen Geschenke, und am Abend gehen viele Menschen in die Kirche. Am sechsundzwanzigsten Dezember ist der zweite Weihnachtsfeiertag — auch das ist in Deutschland ein Feiertag.\par
Und du? Wann ist dein Geburtstag? Mein Geburtstag ist am siebten August. An diesem Tag laden meine Freunde zu mir ein, wir essen Kuchen und hören Musik. Jeder Tag im Kalender kann ein Fest sein!
\end{textebox}

\textbf{Glossaire :} \all{das Neujahr} (-) : le Nouvel An ; \all{aufstehen} (steht auf, stand auf, ist aufgestanden) : se lever ; \all{das Osterei} (-er) : l'œuf de Pâques ; \all{suchen} : chercher ; \all{der Feiertag} (-e) : le jour férié ; \all{das Geschäft} (-e) : le magasin ; \all{geschlossen} : fermé ; \all{marschieren} : défiler ; \all{wehen} : flotter (au vent) ; \all{antworten} : répondre ; \all{das Geschenk} (-e) : le cadeau ; \all{die Kirche} (-n) : l'église ; \all{einladen} (lädt ein) : inviter ; \all{der Kuchen} (-) : le gâteau.

\textbf{Questions de compréhension :}
\begin{enumerate}
\item \all{Wann ist der Tag der Arbeit?}
\item \all{Was gibt es am zwanzigsten Mai in Yaoundé?}
\item \all{Wann feiert Deutschland die Deutsche Einheit?}
\item \all{Was machen die Kinder am fünfundzwanzigsten Dezember?}
\item \all{Wann hat der Erzähler Geburtstag?}
\end{enumerate}

\textbf{Réponses attendues (modèles) :}
\begin{enumerate}
\item \all{Der Tag der Arbeit ist am ersten Mai.}
\item \all{Am zwanzigsten Mai gibt es in Yaoundé eine große Parade.}
\item \all{Deutschland feiert die Deutsche Einheit am dritten Oktober.}
\item \all{Am fünfundzwanzigsten Dezember bekommen die Kinder Geschenke.}
\item \all{Der Erzähler hat am siebten August Geburtstag.}
\end{enumerate}

\courssub{Exercice résolu et pièges à éviter}

\begin{exoresolu}[Écris les dates en toutes lettres]
Énoncé : écris en toutes lettres les dates suivantes, comme dans une phrase complète : 1) am 1. Januar ; 2) der 20. Mai ; 3) am 3. Oktober ; 4) der 7. August ; 5) am 25. Dezember.
\tcblower
Solution détaillée :
\begin{enumerate}
\item \all{am ersten Januar} — ordinal irrégulier \all{erste}, datif en \all{-en} après \all{am}.
\item \all{der zwanzigste Mai} — à partir de 20, terminaison \all{-ste} ; nominatif en \all{-e} après \all{der}.
\item \all{am dritten Oktober} — ordinal irrégulier \all{dritte}, datif en \all{-en}.
\item \all{der siebte August} — ordinal irrégulier \all{siebte} (le \all{-en} de \all{sieben} tombe).
\item \all{am fünfundzwanzigsten Dezember} — 25 $>$ 20, donc \all{-ste} + datif \all{-en}.
\end{enumerate}
\end{exoresolu}

\begin{attention}[Erreurs fréquentes des francophones]
\begin{itemize}
\item \textbf{Oublier le point} après le chiffre : on écrit \all{am 3. Oktober}, jamais « am 3 Oktober ». Sans point, le chiffre se lirait \all{drei} et la phrase serait incorrecte.
\item \textbf{Oublier l'accord} : après \all{am}, l'ordinal est au datif en \all{-en} : \all{am ersten Mai}, et non « am erste Mai ».
\item \textbf{Calquer le français} : en français on dit « le 20 mai » avec un simple cardinal ; en allemand, l'ordinal est obligatoire : \all{der zwanzigste Mai}.
\item \textbf{Mélanger les terminaisons} : \all{-te} seulement de 1 à 19 ; à partir de 20, c'est toujours \all{-ste} : \all{der zwanzigste}, et non « der zwanzigte ».
\item \textbf{Ordre des mots} : quand la date ouvre la phrase, le verbe reste en deuxième position : \all{Am zwanzigsten Mai \textbf{feiert} Kamerun seinen Nationalfeiertag.}
\end{itemize}
\end{attention}

\begin{exercice}[À toi de jouer]
1) Traduis en allemand : « Aujourd'hui, c'est le premier janvier. » — « Le défilé est le 20 mai. » — « Noël est le 25 décembre. »\par
2) Réponds en phrases complètes : \all{Der wievielte ist heute? Wann ist dein Geburtstag? Wann ist der Nationalfeiertag in Kamerun?}\par
3) Écris en toutes lettres : der 2. Februar ; am 12. April ; der 31. Oktober.
\end{exercice}

\begin{corrige}[Exercice « À toi de jouer »]
1) \all{Heute ist der erste Januar.} — \all{Die Parade ist am zwanzigsten Mai.} — \all{Weihnachten ist am fünfundzwanzigsten Dezember.}\par
2) Réponses possibles : \all{Heute ist der fünfte Juni. Mein Geburtstag ist am siebten August. Der Nationalfeiertag in Kamerun ist am zwanzigsten Mai.}\par
3) \all{der zweite Februar} ; \all{am zwölften April} ; \all{der einunddreißigste Oktober}.
\end{corrige}

\begin{aretenir}[L'essentiel de la section]
\begin{itemize}
\item Ordinaux : cardinal + \all{-te} (1--19), cardinal + \all{-ste} (à partir de 20) ; irréguliers : \all{erste}, \all{dritte}, \all{siebte}.
\item Le point après le chiffre marque l'ordinal : \all{3.} = \all{dritte}.
\item \all{Der wievielte ist heute?} $\rightarrow$ \all{Heute ist der zwanzigste Mai.} ; \all{Wann?} $\rightarrow$ \all{am zwanzigsten Mai} (datif en \all{-en}).
\item Écriture abrégée : \all{der 20.5.2025} = \all{der zwanzigste Mai zweitausendfünfundzwanzig}.
\end{itemize}
\end{aretenir}

\coursec{Grammaire 2 : les verbes séparables}

Dans la section précédente, nous avons appris à dire les dates avec les nombres ordinaux : \all{Weihnachten ist am fünfundzwanzigsten Dezember}. Mais pour raconter ce qui se passe ce jour-là — se lever tôt, participer à la fête, assister au défilé —, il faut maîtriser une structure typiquement allemande : les \cle{verbes séparables} (en allemand \all{trennbare Verben}). C'est l'une des grandes difficultés pour les francophones, car le français ne connaît rien d'équivalent. Observons d'abord, puis réglons.

\courssub{Observation : la particule part en fin de phrase}

Reprenons trois phrases tirées de notre texte support \emph{Feste in Deutschland} :

\begin{exemplebox}[Phrases observées dans le texte]
\all{Die Kinder stehen früh auf.} « Les enfants se lèvent tôt. »

\all{Ostern findet im Frühling statt.} « Pâques a lieu au printemps. »

\all{Ich mache gern mit.} « Je participe volontiers. »
\end{exemplebox}

Observons attentivement. Dans le dictionnaire, ces verbes s'écrivent en \textbf{un seul mot} : \all{aufstehen} (se lever), \all{stattfinden} (avoir lieu), \all{mitmachen} (participer). Pourtant, dans la phrase conjuguée, le verbe est \textbf{coupé en deux morceaux} : le radical conjugué (\all{stehen} → \all{stehen}, \all{finden} → \all{findet}, \all{machen} → \all{mache}) reste à sa place habituelle, et la petite particule (\all{auf}, \all{statt}, \all{mit}) est \textbf{rejetée tout à la fin de la phrase}.

\begin{definition}[Verbe séparable]
Un \cle{verbe séparable} (\all{trennbares Verb}) est un verbe composé d'un \textbf{préfixe} (\all{auf-}, \all{ein-}, \all{statt-}, \all{mit-}, \all{an-}, \all{aus-}...) et d'un \textbf{verbe de base}. À l'infinitif et au dictionnaire, il s'écrit en un seul mot : \all{aufstehen}, \all{stattfinden}, \all{mitmachen}. Mais quand le verbe est conjugué, le préfixe se détache et va en fin de phrase.
\end{definition}

\begin{popfigure}
\begin{tikzpicture}[font=\small]
\node[draw=popblue, fill=popblueL, rounded corners, minimum width=1.2cm, minimum height=0.8cm] (s) at (0,0) {Ich};
\node[draw=poporange, fill=poporangeL, rounded corners, minimum width=1.6cm, minimum height=0.8cm] (v) at (2.2,0) {\textbf{stehe}};
\node[draw=popmuted, fill=popcream, rounded corners, minimum width=2.4cm, minimum height=0.8cm] (m) at (5.2,0) {um 6 Uhr};
\node[draw=popgreen, fill=popgreenL, rounded corners, minimum width=1.2cm, minimum height=0.8cm] (p) at (8.2,0) {\textbf{auf}.};
\draw[-{Stealth[length=3mm]}, very thick, poporange] (s.east) -- (v.west);
\draw[-{Stealth[length=3mm]}, very thick, popgreen] (v.south) to[bend right=30] node[below, align=center, text=popdark]{particule rejetée\\en fin de phrase} (p.south);
\node[align=center, text=popdark] at (4.1,-2.2) {Le verbe conjugué (position 2) et sa particule forment une \textbf{pince} qui encadre la phrase.};
\end{tikzpicture}
\legende{La « pince » du verbe séparable : \all{Ich stehe um 6 Uhr auf.} (Je me lève à 6 heures.)}
\end{popfigure}

\courssub{La règle au présent}

\begin{propriete}[Verbe séparable au présent]
Au présent (et au Präteritum), dans une phrase affirmative :
\begin{itemize}
\item le \textbf{verbe conjugué} occupe la \textbf{position 2} (comme tout verbe conjugué allemand) ;
\item la \textbf{particule} est rejetée \textbf{en fin de phrase}.
\end{itemize}
\all{Die Parade fängt um 9 Uhr an.} « Le défilé commence à 9 h. » (\all{anfangen} = commencer)
\end{propriete}

Cette structure en « pince » est fondamentale en allemand : tout ce qui se trouve entre le verbe conjugué et la particule est « encadré » par les deux morceaux du verbe. Même dans une longue phrase, la particule attend patiemment à la toute fin :

\begin{exemplebox}[Exemples traduits au présent]
\all{Wir stehen am Feiertag spät auf.} « Nous nous levons tard le jour férié. »

\all{Der Gottesdienst fängt am Sonntag um 10 Uhr an.} « L'office religieux commence dimanche à 10 h. »

\all{Machst du bei der Parade mit?} « Tu participes au défilé ? »

\all{Die Familie kommt zu Weihnachten zusammen.} « La famille se réunit à Noël. » (\all{zusammenkommen} = se réunir)

\all{Ich lade meine Freunde zum Fest ein.} « J'invite mes amis à la fête. » (\all{einladen} = inviter)
\end{exemplebox}

Remarquez la phrase interrogative : \all{Machst du ... mit?} — même à l'interrogative, la particule reste en fin de phrase.

\begin{attention}[Erreur n° 1 des francophones]
Ne jamais coller la particule au verbe conjugué ! *\all{Ich aufstehe um 6 Uhr} est \textbf{faux}. On dit : \all{Ich stehe um 6 Uhr auf.} De même, *\all{Die Parade anfängt um 9 Uhr} est impossible : \all{Die Parade fängt um 9 Uhr an.}
\end{attention}

\courssub{Les préfixes séparables fréquents}

Voici les préfixes séparables les plus courants, avec des verbes utiles pour parler des fêtes :

\begin{center}
\begin{tabularx}{\linewidth}{|l|X|X|}
\hline
\rowcolor{popblueL} \textbf{Préfixe} & \textbf{Verbe} & \textbf{Français} \\
\hline
auf- & aufstehen & se lever \\
\hline
an- & anfangen & commencer \\
\hline
ein- & einladen & inviter \\
\hline
aus- & aussehen & avoir l'air, ressembler \\
\hline
mit- & mitmachen / mitkommen & participer / accompagner \\
\hline
statt- & stattfinden & avoir lieu \\
\hline
zurück- & zurückkommen & revenir \\
\hline
vor- & vorbereiten & préparer \\
\hline
\end{tabularx}
\end{center}

\begin{propriete}[L'accent tonique : le test infaillible]
Dans un verbe séparable, l'\textbf{accent tonique} porte sur le \textbf{préfixe} : on prononce « \textbf{AUF}stehen », « \textbf{AN}fangen », « \textbf{STATT}finden ». Dans un verbe inséparable, le préfixe n'est \textbf{pas accentué} : « be\textbf{SUCH}en », « ver\textbf{STEH}en ». L'oreille est donc le meilleur détecteur !
\end{propriete}

\begin{methode}[Comment reconnaître un verbe séparable ?]
\begin{enumerate}
\item \textbf{À l'oral} : écoute l'accent tonique. Préfixe accentué (\textbf{AUF}stehen) = séparable ; préfixe non accentué (be\textbf{SUCH}en) = inséparable.
\item \textbf{Au dictionnaire} : le verbe séparable est noté avec un point ou une barre : \all{auf·stehen}, \all{statt·finden}. L'inséparable n'a pas de point : \all{besuchen}.
\item \textbf{Par le sens} : si le préfixe garde son sens propre (\all{auf} = en haut → \all{aufstehen} = se lever « vers le haut »), le verbe est souvent séparable. Les préfixes \all{be-}, \all{ver-}, \all{er-}, \all{ent-}, \all{ge-}, \all{zer-} sont \textbf{toujours inséparables}.
\end{enumerate}
\end{methode}

\courssub{Le verbe séparable au parfait (Perfekt)}

Au parfait, le participe passé se forme en \textbf{insérant -ge- entre le préfixe et le radical} : \all{aufstehen} → \all{auf\textbf{ge}standen} ; \all{stattfinden} → \all{statt\textbf{ge}funden} ; \all{einladen} → \all{ein\textbf{ge}laden}. L'auxiliaire (\all{haben} ou \all{sein}) se choisit comme pour le verbe de base : \all{aufstehen} exprime un changement de position → auxiliaire \all{sein}.

\begin{exemplebox}[Exemples traduits au parfait]
\all{Er ist früh aufgestanden.} « Il s'est levé tôt. »

\all{Sie hat mich zum Fest eingeladen.} « Elle m'a invité à la fête. »

\all{Das Fest hat gestern stattgefunden.} « La fête a eu lieu hier. »

\all{Die Parade hat pünktlich angefangen.} « Le défilé a commencé à l'heure. »

\all{Wir haben bei dem Gottesdienst mitgemacht.} « Nous avons participé à l'office. »
\end{exemplebox}

Notez bien la place : au parfait, le participe (\all{aufgestanden}) va en fin de phrase \textbf{en un seul mot} — le -ge- recolle les deux morceaux !

\begin{attention}[Un cas particulier : vorbereiten]
Le verbe \all{vorbereiten} (préparer) est bien séparable : \all{Ich bereite das Fest vor.} Mais son participe est \all{vorbereitet}, \textbf{sans -ge-}. Pourquoi ? Parce que le radical \all{bereiten} commence déjà par un préfixe inséparable non accentué (\all{be-}), et le -ge- ne se place jamais devant une syllabe non accentuée. Même phénomène : \all{verkaufen} → \all{verkauft}. Retenez simplement : \all{hat vorbereitet}.
\end{attention}

\courssub{Avec un verbe modal ou à l'infinitif}

\begin{propriete}[Verbe séparable après un modal]
Quand le verbe séparable dépend d'un \textbf{verbe modal} (\all{müssen}, \all{wollen}, \all{können}...) ou reste à l'\textbf{infinitif}, il ne se sépare \textbf{pas} : il reste en \textbf{un seul mot}, à la fin de la phrase.

\all{Ich muss früh aufstehen.} « Je dois me lever tôt. »

\all{Willst du bei der Parade mitmachen?} « Veux-tu participer au défilé ? »

\all{Am Nationalfeiertag wollen wir das Fest vorbereiten.} « Le jour de la fête nationale, nous voulons préparer la fête. »
\end{propriete}

Logique : la séparation n'a lieu que quand le verbe est \textbf{conjugué}. À l'infinitif, il reste entier.

\courssub{Piège : les verbes inséparables}

Certains verbes ont aussi un préfixe, mais celui-ci \textbf{ne se détache jamais}. Ce sont les verbes \cle{inséparables} (\all{untrennbare Verben}), formés avec les préfixes \all{be-}, \all{ver-}, \all{er-}, \all{ent-}, \all{emp-}, \all{ge-}, \all{zer-} : \all{besuchen} (visiter), \all{verstehen} (comprendre), \all{gefallen} (plaire), \all{erklären} (expliquer), \all{beginnen} (commencer).

Deux conséquences capitales :
\begin{itemize}
\item \textbf{pas de séparation} : \all{Ich besuche die Kirche.} (jamais *\all{ich suche die Kirche be}) ;
\item \textbf{pas de -ge- au parfait} : \all{besucht}, \all{verstanden}, \all{gefallen} (jamais *\all{bege sucht}).
\end{itemize}

\begin{center}
\begin{tabularx}{\linewidth}{|X|X|X|}
\hline
\rowcolor{popblueL} \textbf{Critère} & \textbf{Séparable} & \textbf{Inséparable} \\
\hline
Accent tonique & sur le préfixe : \textbf{AUF}stehen & sur le radical : be\textbf{SUCH}en \\
\hline
Présent & Ich stehe früh \textbf{auf}. & Ich \textbf{besuche} die Kirche. \\
\hline
Participe passé & auf\textbf{ge}standen & besucht (sans -ge-) \\
\hline
Dictionnaire & auf·stehen (avec point) & besuchen (sans point) \\
\hline
Préfixes typiques & auf-, ein-, aus-, an-, mit-, statt-, vor-, zurück- & be-, ver-, er-, ent-, ge-, zer- \\
\hline
\end{tabularx}
\end{center}

\begin{attention}[Comparaison avec le français]
Le français ne possède \textbf{aucun} verbe séparable : on ne dit pas « je lève me tôt » ! Le phénomène se rapproche des verbes à particule de l'anglais (\emph{to get up}, \emph{to take part}). Conséquence : \textbf{traduire mot à mot est impossible}. Il faut apprendre chaque verbe séparable comme un vocabulaire entier : \all{aufstehen} = « se lever », \all{stattfinden} = « avoir lieu », et non pas additionner les sens des morceaux. Attention aussi au faux ami : \all{bekommen} ne signifie pas « devenir » mais « recevoir » !
\end{attention}

\courssub{Texte d'application : Unser Nationalfeiertag}

\begin{textebox}[Unser Nationalfeiertag in Yaoundé]
Am zwanzigsten Mai stehen wir in Yaoundé sehr früh auf. Der Nationalfeiertag fängt um 9 Uhr an, aber meine Familie bereitet alles schon am Abend vor. Meine Mutter kocht, mein Vater holt die Fahne aus dem Schrank, und ich lade meine Freunde ein. Alle wollen bei der Parade mitmachen.

Um 8 Uhr kommen die Schüler auf dem Boulevard zusammen. Sie tragen die Farben des Landes: Grün, Rot und Gelb. Die Parade fängt pünktlich an. Die Soldaten marschieren, die Schüler singen, und alle schauen stolz zu. Mein Bruder macht Fotos, denn er kommt jedes Jahr mit.

Am Nachmittag findet in unserer Kirche ein Gottesdienst statt. Der Pastor dankt Gott für unser Land. Danach gehen wir nach Hause zurück und essen zusammen. Am Abend lädt unser Onkel die ganze Familie zum Essen ein. Er hat ein großes Fest vorbereitet.

Letztes Jahr hat das Fest wegen des Regens nicht stattgefunden, aber dieses Jahr scheint die Sonne. „Unser Nationalfeiertag gefällt mir sehr“, sagt meine kleine Schwester. „Ich stehe gern früh auf, wenn es ein Fest gibt!“ Wir lachen alle. Der Feiertag ist für uns der schönste Tag des Jahres.
\end{textebox}

\textbf{Glossaire :} \all{der Schrank, die Schränke} (l'armoire) ; \all{die Fahne, die Fahnen} (le drapeau) ; \all{zusammenkommen} (se réunir) ; \all{marschieren} (marcher au pas) ; \all{zuschauen} (regarder, assister) ; \all{der Gottesdienst, die Gottesdienste} (l'office religieux) ; \all{zurückgehen} (rentrer) ; \all{wegen} (à cause de) ; \all{scheinen} (briller) ; \all{der Feiertag, die Feiertage} (le jour férié).

\textbf{Questions de compréhension :}
\begin{enumerate}
\item Wann fängt der Nationalfeiertag an?
\item Wer bereitet das Fest vor?
\item Was findet am Nachmittag statt?
\item Warum hat das Fest letztes Jahr nicht stattgefunden?
\end{enumerate}

\textbf{Chasse grammaticale :} soulignez dans le texte tous les verbes séparables conjugués et entourez leur particule en fin de phrase. Vous devez en trouver plus de dix : \all{stehen ... auf}, \all{fängt ... an} (deux fois), \all{bereitet ... vor}, \all{lade ... ein}, \all{mitmachen} (infinitif après \all{wollen} : non séparé !), \all{kommen ... zusammen}, \all{schauen ... zu}, \all{kommt ... mit}, \all{findet ... statt}, \all{gehen ... zurück}, \all{lädt ... ein}, \all{hat ... stattgefunden}, \all{stehe ... auf}. Repérez aussi les deux verbes inséparables : \all{gefällt} (de \all{gefallen}) et \all{hat ... vorbereitet} (participe sans -ge-, voir l'encadré ci-dessus).

\begin{exoresolu}[Transformez les phrases]
Énoncé — Mettez les phrases suivantes au présent en séparant correctement le verbe, puis traduisez : 1. \all{die Kinder / aufstehen / früh} ; 2. \all{das Fest / stattfinden / am Sonntag} ; 3. \all{du / mitmachen / bei der Parade ?} ; 4. Mettez au parfait : \all{er / einladen / seine Freunde}.
\tcblower
Solution détaillée :
\begin{enumerate}
\item \all{Die Kinder stehen früh auf.} « Les enfants se lèvent tôt. » — verbe conjugué \all{stehen} en position 2, particule \all{auf} en fin de phrase.
\item \all{Das Fest findet am Sonntag statt.} « La fête a lieu dimanche. » — \all{findet} en position 2, \all{statt} à la fin.
\item \all{Machst du bei der Parade mit?} « Tu participes au défilé ? » — à l'interrogative, le verbe en position 1, la particule toujours en fin.
\item \all{Er hat seine Freunde eingeladen.} « Il a invité ses amis. » — participe \all{ein\textbf{ge}laden} : le -ge- s'insère entre \all{ein-} et \all{laden} ; auxiliaire \all{haben} car \all{einladen} n'exprime ni mouvement ni changement.
\end{enumerate}
\end{exoresolu}

\begin{exercice}[À vous !]
1. Conjuguez au présent : \all{ich (aufstehen)}, \all{die Parade (anfangen)}, \all{wir (mitmachen)}, \all{ihr (zurückkommen)}.

2. Mettez au parfait : \all{das Fest (stattfinden)}, \all{sie (einladen) ihre Nachbarn}, \all{wir (aufstehen) um 6 Uhr}.

3. Séparable ou inséparable ? \all{besuchen}, \all{anfangen}, \all{verstehen}, \all{stattfinden}, \all{gefallen}, \all{einladen}. Justifiez avec l'accent tonique.

4. Traduisez : « Le défilé commence à 9 heures. » ; « Hier, la fête a eu lieu. » ; « Je dois me lever tôt. »
\end{exercice}

\begin{corrige}[Exercice : les verbes séparables]
1. \all{ich stehe auf} ; \all{die Parade fängt an} ; \all{wir machen mit} ; \all{ihr kommt zurück}.

2. \all{Das Fest hat stattgefunden.} ; \all{Sie hat ihre Nachbarn eingeladen.} ; \all{Wir sind um 6 Uhr aufgestanden.} (auxiliaire \all{sein} : changement de position).

3. Inséparables : \all{besuchen} (be\textbf{SUCH}en), \all{verstehen} (ver\textbf{STEH}en), \all{gefallen} (ge\textbf{FAL}len) — préfixes non accentués. Séparables : \all{anfangen}, \all{stattfinden}, \all{einladen} — préfixes accentués.

4. \all{Die Parade fängt um 9 Uhr an.} ; \all{Gestern hat das Fest stattgefunden.} ; \all{Ich muss früh aufstehen.}
\end{corrige}

\begin{aretenir}[L'essentiel sur les verbes séparables]
\begin{itemize}
\item Un verbe séparable = préfixe accentué + verbe de base (\all{auf·stehen}).
\item Au présent : verbe conjugué en \textbf{position 2}, particule en \textbf{fin de phrase} : \all{Ich stehe um 6 Uhr auf.}
\item Au parfait : -ge- entre préfixe et radical : \all{aufgestanden}, \all{stattgefunden} ; mais pas de -ge- si le radical commence par une syllabe non accentuée : \all{vorbereitet}.
\item Avec un modal : le verbe reste entier en fin de phrase : \all{Ich muss früh aufstehen.}
\item Les préfixes \all{be-}, \all{ver-}, \all{er-}, \all{ent-}, \all{ge-}, \all{zer-} sont toujours inséparables : pas de séparation, pas de -ge- : \all{besucht}, \all{verstanden}.
\end{itemize}
\end{aretenir}

\coursec{Commentaire et production guidée}

Après avoir maîtrisé les nombres ordinaux pour dater les événements et la structure des verbes à particule séparable pour décrire le déroulement des fêtes, nous allons maintenant mobiliser l'ensemble de ces acquis. Cette section finale vous guide pas à pas pour \cle{commenter le texte support} « Feste in Deutschland » (étudié en section 1) et pour \cle{produire votre propre paragraphe} décrivant une fête camerounaise à un correspondant germanophone. C'est l'épreuve type du Baccalauréat : compréhension, analyse linguistique et expression personnelle comparative.

\courssub{Méthode du commentaire de texte en 4 étapes}

\begin{methode}[Commenter un texte en allemand (niveau A2/B1)]
Suivez rigoureusement ces quatre étapes. Utilisez les phrases-modèles (en \textbf{gras}) pour structurer votre réponse.

\textbf{Étape 1 : Présenter le texte} (\textit{Vorstellen})
\begin{itemize}
    \item Identifiez la source (si connue), le thème principal et le type de texte.
    \item \textbf{Modèle :} \all{Der Text „Feste in Deutschland“ handelt von den wichtigsten Festen in Deutschland. Er informiert über Weihnachten, Ostern und den Nationalfeiertag.}
\end{itemize}

\textbf{Étape 2 : Résumer le contenu} (\textit{Zusammenfassen})
\begin{itemize}
    \item Résumez en 3 phrases maximum l'ordre chronologique ou thématique des idées.
    \item Utilisez les connecteurs logiques : \cle{zuerst} (d'abord), \cle{dann} (ensuite), \cle{schließlich} / \cle{zum Schluss} (enfin).
    \item \textbf{Modèle :} \all{Zuerst beschreibt der Autor Weihnachten am 25. Dezember. Dann erklärt er Ostern im Frühling. Schließlich nennt er den Nationalfeiertag am 3. Oktober.}
\end{itemize}

\textbf{Étape 3 : Analyser la langue} (\textit{Analysieren})
\begin{itemize}
    \item Relevez \textbf{3 à 4 mots du champ lexical des fêtes} (ex. \all{der Gottesdienst}, \all{die Parade}, \all{feiern}, \all{der Feiertag}).
    \item Identifiez \textbf{2 structures grammaticales} étudiées : des dates avec ordinaux (\all{am 3. Oktober}) et des verbes séparables (\all{stattfinden}, \all{aufstehen}).
    \item \textbf{Modèle :} \all{Im Text gibt es viele Wörter zum Thema Feste: „Gottesdienst“, „Parade“, „feiern“. Die Grammatik zeigt Datumsangaben („am 25. Dezember“) und trennbare Verben („findet statt“, „stehen auf“).}
\end{itemize}

\textbf{Étape 4 : Donner son avis personnel et comparer} (\textit{Beurteilen \& Vergleichen})
\begin{itemize}
    \item Exprimez votre opinion et faites le lien avec le Cameroun (contexte personnel).
    \item Structure : \cle{Ich finde ...} (je trouve que...) + \cle{weil} (parce que) + comparaison.
    \item \textbf{Modèle :} \all{Ich finde den Text interessant, weil Kamerun auch einen Nationalfeiertag hat. Der 20. Mai ist bei uns wie der 3. Oktober in Deutschland.}
\end{itemize}
\end{methode}

\begin{popfigure}
\begin{tikzpicture}[
    node distance=1.2cm and 1.5cm,
    >=Stealth,
    box/.style={draw=popdark, thick, rounded corners=6pt, align=center, minimum width=3.2cm, minimum height=1.8cm, font=\small},
    arrow/.style={->, thick, popdark}
]
\node[box, fill=popblueL, align=center] (p){Présentation\\ \textit{(Vorstellen)}\\ \all{Der Text handelt von...}};
\node[box, fill=popgreenL, right=of p, align=center] (r){Résumé\\ \textit{(Zusammenfassen)}\\ \all{Zuerst... dann... schließlich...}};
\node[box, fill=poporangeL, right=of r, align=center] (a){Analyse\\ \textit{(Analysieren)}\\ Lexique, Grammaire, Structure};
\node[box, fill=poppinkL, right=of a, align=center] (j){Avis personnel\\ \textit{(Beurteilen)}\\ \all{Ich finde..., weil... Kamerun...}};

\draw[arrow] (p) -- (r);
\draw[arrow] (r) -- (a);
\draw[arrow] (a) -- (j);
\end{tikzpicture}
\legende{Organigramme des 4 étapes du commentaire de texte}
\end{popfigure}

\courssub{Commentaire modèle rédigé : « Feste in Deutschland »}

\begin{textebox}[Commentaire modèle — Niveau A2+]
Der Text „Feste in Deutschland“ handelt von den drei wichtigsten Festen in Deutschland: Weihnachten, Ostern und der Nationalfeiertag.

Zuerst beschreibt der Autor Weihnachten. Es findet am 25. und 26. Dezember statt. Die Menschen gehen in den Gottesdienst und essen mit der Familie. Dann erklärt der Text Ostern im Frühling. Es ist ein religiöses Fest, aber auch ein Frühlingsfest mit Eiersuchen. Schließlich nennt der Autor den Nationalfeiertag am 3. Oktober. An diesem Tag gibt es Paraden und Reden. Die Fahne wird gehisst.

Ich finde den Text sehr informativ, weil er die Termine und Bräuche klar erklärt. Besonders interessant finde ich den Nationalfeiertag, weil Kamerun auch einen Nationalfeiertag hat. Unser Nationalfeiertag ist am 20. Mai. Da gibt es auch Paraden und die Schüler machen bei der Parade mit. Die Feste sind in beiden Ländern wichtig für die Identität.
\end{textebox}

\begin{definition}[Glossaire du commentaire modèle]
\begin{itemize}
    \item \all{der Brauch, die Bräuche} : la coutume, la tradition.
    \item \all{der Frühling} : le printemps.
    \item \all{das Eiersuchen} : la chasse aux œufs (de Pâques).
    \item \all{hissen} (regulier) : hisser (un drapeau). Participe passé : \all{gehisst}.
    \item \all{die Identität} : l'identité.
    \item \all{informativ} : informatif.
    \item \all{klären} : expliquer, clarifier.
\end{itemize}
\end{definition}

\courssub{Production écrite guidée : « Eine Feier in Kamerun »}

\begin{methode}[Rédiger un paragraphe descriptif (6--8 phrases)]
Voici la marche à suivre pour écrire votre texte à destination d'un correspondant allemand (ex. email, lettre, forum d'échange).

\textbf{1. Planifiez avec la grille d'aide (ci-dessous).} Notez les mots-clés en allemand.
\textbf{2. Rédigez la phrase d'accroche :} Nommez la fête et la date. \all{In Kamerun feiern wir den Nationalfeiertag am 20. Mai.}
\textbf{3. Décrivez le déroulement (2--3 phrases) :} Utilisez \textbf{au moins 2 verbes séparables} (ex. \all{stattfinden}, \all{aufstehen}, \all{mitmachen}, \all{anziehen}, \all{vorbereiten}). Respectez la place de la particule (fin de phrase).
\textbf{4. Ajoutez un détail culturel :} Nourriture, tenue, ambiance. \all{Wir essen Ndolé und Plantains.}
\textbf{5. Concluez par votre avis :} \all{Ich finde das Fest toll, weil...}
\textbf{6. Relisez-vous :} Vérifiez l'ordre des mots (V2), les ordinaux (\all{am zwanzigsten Mai}), les majuscules des noms.
\end{methode}

\begin{exemplebox}[Phrases-modèles pour la production (à adapter)]
\begin{itemize}
    \item \all{In Kamerun findet der Nationalfeiertag am zwanzigsten Mai statt.} (Au Cameroun, la fête nationale a lieu le 20 mai.)
    \item \all{Die Feierlichkeiten beginnen früh morgens.} (Les festivités commencent tôt le matin.)
    \item \all{Die Schüler stehen sehr früh auf und ziehen ihre Uniformen an.} (Les élèves se lèvent très tôt et enfilent leur uniforme.)
    \item \all{Bei der großen Parade machen alle Schulen mit.} (Au grand défilé, toutes les écoles participent.)
    \item \all{Nach der Parade essen die Familien zusammen: oft Ndolé oder Poulet DG.} (Après le défilé, les familles mangent ensemble : souvent Ndolé ou Poulet DG.)
    \item \all{Ich finde diesen Tag sehr wichtig, weil er an die Einheit des Landes erinnert.} (Je trouve ce jour très important car il rappelle l'unité du pays.)
    \item \all{Abends gibt es oft Musik und Tanz, zum Beispiel Bikutsi oder Makossa.} (Le soir, il y a souvent de la musique et de la danse, par exemple Bikutsi ou Makossa.)
\end{itemize}
\end{exemplebox}

\courssub{Grille d'aide à la production (Fiche élève)}

\begin{center}
\begin{tabularx}{\linewidth}{|X|X|X|}
\hline
\rowcolor{popblueL} \textbf{Étape / Funktion} & \textbf{Contenu attendu} & \textbf{Mots-clés / Structures obligatoires} \\ \hline
\textbf{1. Nommer} & Nom de la fête & \all{der Nationalfeiertag} (20. Mai), \all{der Jugendtag} (11. Februar), \all{Weihnachten}, \all{das Ende des Ramadan}, \all{Tabaski} \\ \hline
\textbf{2. Dater} & Date avec ordinal & \all{am} + ordinal (Datif) : \all{am zwanzigsten Mai}, \all{am elften Februar}, \all{am fünfundzwanzigsten Dezember} \\ \hline
\textbf{3. Décrire : L'événement principal} & Verbe séparable \textbf{stattfinden} & \all{Das Fest findet am ... statt.} (Particule \all{statt} à la fin) \\ \hline
\textbf{4. Décrire : Actions des gens} & \textbf{2 verbes séparables} au choix & \all{aufstehen} (se lever), \all{mitmachen} (participer), \all{anziehen} (s'habiller), \all{vorbereiten} (préparer), \all{aufräumen} (ranger) \\ \hline
\textbf{5. Détail culturel} & Nourriture, musique, tenue & \all{essen} (\all{das} Ndolé, \all{das} Koki, \all{die} Plantains), \all{tanzen} (\all{der} Bikutsi, \all{der} Makossa), \all{tragen} (\all{die} Uniform, \all{der} Pagne) \\ \hline
\textbf{6. Avis personnel} & Opinion + justification & \all{Ich finde das Fest ... (wichtig / schön / toll), weil ...} (verbe à la fin après \all{weil}) \\ \hline
\end{tabularx}
\end{center}

\begin{attention}[Pièges à éviter pour la production — Spécial Francophones]
\begin{enumerate}
    \item \textbf{Particule oubliée ou mal placée :} \all{*Das Fest stattfindet am 20. Mai.} $\rightarrow$ \textbf{Faux.} La particule \all{statt} DOIT aller à la fin : \all{Das Fest findet am 20. Mai \textbf{statt}.}
    \item \textbf{Ordinal mal accordé :} \all{*am zwanzigste Mai} $\rightarrow$ \textbf{Faux.} Après \all{am} (an dem), c'est le datif : \all{am \textbf{zwanzigsten} Mai}. (Terminaison \textbf{-en} pour tous les genres au datif singulier avec article défini).
    \item \textbf{Ordre des mots (V2) :} \all{*Am 20. Mai das Fest findet statt.} $\rightarrow$ \textbf{Faux.} Le verbe conjugué est en 2ème position : \all{Am 20. Mai \textbf{findet} das Fest statt.}
    \item \textbf{Subordonnée avec \all{weil} :} \all{*Ich finde es gut, weil es ist wichtig.} $\rightarrow$ \textbf{Faux.} \all{weil} renvoie le verbe à la fin : \all{Ich finde es gut, weil es wichtig \textbf{ist}.}
    \item \textbf{Majuscules :} \all{*der nationalfeiertag, die parade, die fahne} $\rightarrow$ \textbf{Faux.} \textbf{Tous les noms} prennent une majuscule : \all{der Nationalfeiertag, die Parade, die Fahne}.
\end{enumerate}
\end{attention}

\courssub{Expression orale guidée : Présentation et interaction}

\begin{experience}[Activité de classe : « Mein Fest — Dein Fest »]
\textbf{Consigne :}
\begin{enumerate}
    \item \textbf{Préparation (5 min) :} Relisez votre paragraphe écrit. Entraînez-vous à le dire sans lire (regardez vos mots-clés sur la grille).
    \item \textbf{Présentation (2 min / élève) :} Présentez votre fête à la classe (ou en binôme). Parlez distinctement. \textit{Ex: „Guten Tag. Ich erzähle euch vom Nationalfeiertag in Kamerun...“}
    \item \textbf{Questions / Réponses (3 min) :} Deux camarades posent une question chacun. Vous répondez par une phrase complète.
    \item \textbf{Exemples de questions types :}
    \begin{itemize}
        \item \all{Wann findet das Fest statt?} (Quand a lieu la fête ?)
        \item \all{Was machen die Schüler?} (Que font les élèves ?)
        \item \all{Was isst man an diesem Tag?} (Qu'est-ce qu'on mange ce jour-là ?)
        \item \all{Findest du das Fest wichtig? Warum?} (Trouves-tu cette fête importante ? Pourquoi ?)
    \end{itemize}
\end{enumerate}
\textbf{Critères d'évaluation orale :} Prononciation des Umlaute (\all{zwanzigsten}, \all{Mai}), fluidité (pas de lecture), correction de la particule finale, utilisation de \all{weil} + verbe final.
\end{experience}

\courssub{Auto-évaluation : Suis-je prêt pour le Bac ?}

Utilisez cette grille pour vérifier votre production écrite ou orale avant de la rendre. Total : \textbf{6 points}.

\begin{center}
\begin{tabular}{|p{5cm}|c|c|}
\hline
\rowcolor{popblueL} \textbf{Critère} & \textbf{Oui (2 pts)} & \textbf{Non / Partiel (0-1 pt)} \\ \hline
\textbf{Vocabulaire du thème} : Utilisation correcte d'au moins 5 mots du thème (fêtes, date, activités, nourriture, avis) avec articles et pluriels. & & \\ \hline
\textbf{Date correcte} : Utilisation de \all{am} + ordinal au datif (\all{am zwanzigsten Mai}) sans faute d'accord. & & \\ \hline
\textbf{Verbes séparables} : Au moins 2 verbes séparables employés correctement (particule en fin de phrase, verbe conjugué en 2ème position). & & \\ \hline
\textbf{TOTAL} & \multicolumn{2}{c|}{\textbf{/ 6}} \\ \hline
\end{tabular}
\end{center}

\begin{savaistu}[Le Bac et la comparaison interculturelle]
Au Baccalauréat (série A), l'épreuve d'expression écrite (souvent 10--12 points sur 20) propose fréquemment un sujet de type : \emph{„Vergleichen Sie ein Fest in Deutschland mit einem Fest in Kamerun.“} (Comparez une fête en Allemagne avec une fête au Cameroun) ou \emph{„Schreiben Sie einem deutschen Freund, wie man in Kamerun den Nationalfeiertag feiert.“}.

La production que vous venez de faire (\textbf{« Eine Feier in Kamerun »}) est \textbf{exactement} ce format d'examen. Maîtriser la structure « Présentation $\rightarrow$ Description (verbes séparables + dates) $\rightarrow$ Avis comparatif » vous assure les points de \textbf{contenu} et de \textbf{structure grammaticale}. Entraînez-vous à écrire ce paragraphe sur 3 fêtes différentes (Nationalfeiertag, Jugendtag, Fête religieuse) pour être prêt le jour J.
\end{savaistu}

\courssub{Exercice résolu : Correction d'une production élève}

\begin{exoresolu}[Corriger les erreurs classiques]
\textbf{Énoncé :} Un élève a écrit ce paragraphe. Repérez et corrigez les 5 erreurs (grammaire, syntaxe, orthographe).
\all{In Kamerun der Nationalfeiertag findet am 20. Mai statt. Die Schüler stehen früh auf und machen bei der Parade mit. Wir essen Ndolé und trinken Bière. Ich finde das Fest toll, weil es ist wichtig für die Einheit. Am Abend tanzen wir Makossa.}

\tcblower
\textbf{Solution détaillée :}

\begin{enumerate}
    \item \textbf{Ordre des mots (V2) :} \all{*In Kamerun der Nationalfeiertag findet...} $\rightarrow$ \all{In Kamerun \textbf{findet} der Nationalfeiertag...} (Le verbe \all{findet} doit être en 2ème position après le complément de lieu \all{In Kamerun}).
    \item \textbf{Ordinal (Datif) :} \all{*am 20. Mai} (chiffre) $\rightarrow$ En toutes lettres : \all{am \textbf{zwanzigsten} Mai}. (Rappel : \all{am} = \all{an dem} $\rightarrow$ Datif $\rightarrow$ adjectif en \textbf{-en}).
    \item \textbf{Verbe séparable \all{mitmachen} :} \all{machen ... mit} est correct ici (particule \all{mit} en fin de phrase). \all{bei der Parade} (Datif singulier féminin) est correct.
    \item \textbf{Majuscule / Genre :} \all{*Bière} $\rightarrow$ \all{\textbf{das} \textbf{Bier}} (Nom allemand, majuscule, genre neutre). On n'écrit pas le mot français dans une phrase allemande.
    \item \textbf{Subordonnée \all{weil} :} \all{*weil es ist wichtig} $\rightarrow$ \all{weil es wichtig \textbf{ist}}. (Verbe conjugué \textbf{à la fin} après \all{weil}).
\end{enumerate}

\textbf{Version corrigée :}
\all{In Kamerun findet der Nationalfeiertag am zwanzigsten Mai statt. Die Schüler stehen früh auf und machen bei der Parade mit. Wir essen Ndolé und trinken Bier. Ich finde das Fest toll, weil es wichtig für die Einheit ist. Am Abend tanzen wir Makossa.}
\end{exoresolu}

\courssub{Exercice d'entraînement : À vous de jouer !}

\begin{exercice}[Production écrite : « Der Jugendtag in Kamerun »]
Rédigez un paragraphe de \textbf{7 à 8 phrases} en allemand pour décrire la \textbf{Fête de la Jeunesse (11 février)} à un correspondant allemand. Utilisez la \textbf{Grille d'aide} (page précédente) et respectez la \textbf{Méthode en 6 étapes}.

\textbf{Contraintes obligatoires :}
\begin{itemize}
    \item Nommer la fête : \all{der Jugendtag} (ou \all{der Tag der Jugend}).
    \item Date correcte : \all{am elften Februar}.
    \item \textbf{3 verbes séparables} différents (ex. \all{stattfinden}, \all{aufstehen}, \all{mitmachen}, \all{anziehen}, \all{vorbereiten}, \all{aufräumen}).
    \item Un détail culturel camerounais (défilé jeunesse, chants, tenues, discours).
    \item Avis personnel avec \all{weil} + verbe final.
\end{itemize}

\textbf{Aide au démarrage :}
\all{In Kamerun feiern wir den Jugendtag am elften Februar. Dieses Fest findet jedes Jahr statt. ...}
\end{exercice}

\begin{corrige}[Exercice : Der Jugendtag in Kamerun (Proposition de correction)]
\all{In Kamerun feiern wir den Jugendtag am elften Februar. Dieses Fest findet jedes Jahr in allen Städten statt. Die Schüler stehen sehr früh auf und ziehen ihre schönen Uniformen an. Bei den Paraden machen alle Schulen und Jugendgruppen mit. Man hört viele Reden und singt die Nationalhymne. Nach dem Umzug essen die Familien oft zusammen Poulet DG oder Riz au Gras. Ich finde diesen Tag sehr wichtig, weil die Jugend die Zukunft des Landes ist. Abends tanzen wir oft Makossa und Bikutsi.}
\end{corrige}

\newpage

\coursec{Activité d'intégration}

\courssub{Objectifs de l'activité}

À l'issue de cette activité d'intégration, tu seras capable de :
\begin{itemize}
\item mobiliser le \cle{vocabulaire thématique des fêtes} (religieuses et nationales) dans une situation de communication réelle ;
\item écrire une \cle{date} correctement, avec un \cle{nombre ordinal} bien formé (\all{am 20. Mai}, \all{der 3. Oktober}) ;
\item employer au moins trois \cle{verbes à particule séparable} en respectant la place de la particule en fin de phrase ;
\item rédiger une \cle{lettre d'invitation} simple et correcte (8 à 10 phrases) adressée à un correspondant germanophone ;
\item lire ta production à voix haute avec une prononciation claire et \cle{comprendre} les lettres des camarades (repérage de la fête, de la date et des verbes séparables) ;
\item comparer une fête camerounaise avec une fête allemande, dans une démarche interculturelle.
\end{itemize}

\courssub{Situation}

Ton lycée de Yaoundé participe à un échange scolaire avec le Gymnasium de Bamberg (Bavière, Allemagne). Ta correspondante allemande, Lena, t'a envoyé une lettre dans laquelle elle décrit la fête nationale allemande, le \all{Tag der Deutschen Einheit}, et te pose une question : comment fête-t-on la fête nationale au Cameroun ? Lis d'abord sa lettre.

\begin{textebox}[Brief aus Bamberg — Lettre de Lena]
Liebe Freundin, lieber Freund,

wie geht es dir? Mir geht es gut. Heute schreibe ich über unsere Feste in Deutschland. Am 3. Oktober feiern wir den Tag der Deutschen Einheit. Das ist unser Nationalfeiertag. An diesem Tag stehen wir spät auf, denn die Schule und die Büros sind geschlossen. In Berlin findet ein großes Fest statt: Es gibt Musik, Essen und eine Parade. Viele Leute machen beim Fest mit und tragen die Farben der Fahne: Schwarz, Rot und Gold. Das Fest fängt am Morgen an und endet am Abend mit einem Feuerwerk.

Am 25. Dezember feiern wir Weihnachten mit der Familie, und an Ostern suchen die Kinder Eier im Garten. Am Sonntag gehen viele Familien zum Gottesdienst in die Kirche.

Und du? Wann ist der Nationalfeiertag in Kamerun? Was macht ihr an diesem Tag? Ich lade dich ein: Schreib mir einen Brief und lade mich zu eurem Fest ein!

Viele Grüße\\
Lena
\end{textebox}

\textbf{Glossaire :} \all{die Einheit} (l'unité) ; \all{der Tag der Deutschen Einheit} (le Jour de l'unité allemande, fête nationale du 3 octobre) ; \all{geschlossen} (fermé) ; \all{das Feuerwerk, die Feuerwerke} (le feu d'artifice) ; \all{suchen} (chercher) ; \all{das Ei, die Eier} (l'œuf) ; \all{der Garten, die Gärten} (le jardin) ; \all{tragen} (porter — ici : arborer) ; \all{Viele Grüße} (amitiés, formule de fin de lettre).

Tu vas maintenant répondre à Lena et l'inviter à une fête camerounaise de ton choix : la \cle{fête nationale du 20 mai}, \cle{Noël} ou \cle{Pâques} au village.

\courssub{Consignes}

Travail par binômes. Suis les étapes dans l'ordre.

\begin{enumerate}
\item \textbf{Comprendre le document.} Lis la lettre de Lena deux fois : une première fois pour le sens global, une deuxième fois en soulignant les informations demandées. Réponds par écrit, en allemand et en phrases complètes, aux questions suivantes :
\begin{enumerate}
\item \all{Wann ist der Nationalfeiertag in Deutschland?}
\item \all{Was findet in Berlin statt?}
\item \all{Welche Farben hat die deutsche Fahne?}
\item \all{Was suchen die Kinder an Ostern?}
\end{enumerate}
\item \textbf{Repérer la grammaire dans le texte.} Souligne dans la lettre de Lena tous les \cle{verbes séparables} conjugués (il y en a quatre) et entoure les deux \cle{dates avec ordinaux}. Recopie-les dans ton cahier en séparant le verbe et sa particule : \all{stehen ... auf}.
\item \textbf{Préparer le lexique.} Choisis ta fête (fête nationale du 20 mai, Noël ou Pâques au village). Avec ton camarade, liste au moins \textbf{cinq mots du lexique de la leçon} que tu vas utiliser, avec leur article : par exemple \all{der Nationalfeiertag, die Parade, die Fahne, der Gottesdienst, das Fest, die Familie, das Essen, die Musik}.
\item \textbf{Rédiger la lettre d'invitation (8 à 10 phrases).} Ta lettre doit obligatoirement contenir :
\begin{itemize}
\item l'en-tête de la lettre (lieu et date : \all{Yaoundé, den 15. Mai}) et la formule d'appel (\all{Liebe Lena,}) ;
\item la \textbf{date de la fête écrite avec un ordinal correct} (\all{am 20. Mai} — on lit « am zwanzigsten Mai ») ;
\item au moins \textbf{trois verbes séparables} parmi : \all{stattfinden, aufstehen, mitmachen, einladen, anfangen} (attention : la particule va à la fin de la phrase !) ;
\item au moins \textbf{cinq mots du lexique} de la leçon ;
\item \textbf{une phrase de comparaison} avec une fête allemande (structures utiles : \all{Wie in Deutschland ...} « comme en Allemagne » ou \all{Anders als in Deutschland ...} « contrairement à l'Allemagne ») ;
\item une formule de politesse finale (\all{Viele Grüße} / \all{Ich freue mich auf dich!}).
\end{itemize}
\item \textbf{Auto-correction en binôme.} Relis ta lettre avec la grille suivante : majuscule à tous les noms ? verbe conjugué en deuxième position ? particule séparable en fin de phrase ? ordinal avec un point (\all{am 20. Mai}) ? articles corrects ?
\item \textbf{Lecture orale et écoute active.} Chaque binôme lit sa lettre à la classe. Pendant l'écoute, note sur une fiche : (a) la fête choisie, (b) la date entendue, (c) les verbes séparables identifiés.
\item \textbf{Vote final.} La classe vote pour l'invitation la plus convaincante selon trois critères : correction grammaticale, richesse du vocabulaire, clarté de la prononciation.
\end{enumerate}

\courssub{Ressources à mobiliser}

\begin{aretenir}[Rappel 1 : les dates et les ordinaux]
L'ordinal se forme avec le nombre + \textbf{-te} (de 1 à 19) ou \textbf{-ste} (à partir de 20) : \all{der zweite} (2\textsuperscript{e}), \all{der fünfte} (5\textsuperscript{e}), \all{der zwanzigste} (20\textsuperscript{e}). \textbf{Exceptions} : \all{der \textbf{erste}} (1\textsuperscript{er}), \all{der \textbf{dritte}} (3\textsuperscript{e}), \all{der \textbf{siebte}} (7\textsuperscript{e}, sans -t-), \all{der \textbf{achte}} (8\textsuperscript{e}, sans -t-). Dans une date, on écrit le nombre suivi d'un \textbf{point} : \all{am 20. Mai} (le 20 mai), \all{am 3. Oktober} (le 3 octobre), \all{am 25. Dezember} (le 25 décembre). Préposition : \textbf{am} + date. À l'oral, l'ordinal se décline comme un adjectif : \all{am \textbf{zwanzigsten} Mai}, \all{am \textbf{dritten} Oktober}.
\end{aretenir}

\begin{center}
\begin{tabularx}{\linewidth}{|l|X|X|}
\hline
\rowcolor{popblueL} Nombre & Ordinal (écrit) & Date (exemple) \\
\hline
1 & der erste & am 1. Januar (le 1\textsuperscript{er} janvier) \\
\hline
2 & der zweite & am 2. Februar (le 2 février) \\
\hline
3 & der dritte & am 3. Oktober (le 3 octobre) \\
\hline
7 & der siebte & am 7. April (le 7 avril) \\
\hline
8 & der achte & am 8. Mai (le 8 mai) \\
\hline
20 & der zwanzigste & am 20. Mai (le 20 mai) \\
\hline
25 & der fünfundzwanzigste & am 25. Dezember (le 25 décembre) \\
\hline
\end{tabularx}
\end{center}

\begin{aretenir}[Rappel 2 : les verbes séparables]
Au Präsens, le verbe conjugué est en \textbf{deuxième position} et la particule va \textbf{à la fin} de la phrase : \all{Das Fest \textbf{fängt} am Morgen \textbf{an}.} (La fête commence le matin.) — \all{Ich \textbf{lade} dich zu unserem Fest \textbf{ein}.} (Je t'invite à notre fête.) — \all{Die Parade \textbf{findet} am 20. Mai \textbf{statt}.} (Le défilé a lieu le 20 mai.) Si un complément ouvre la phrase, le verbe reste deuxième et le sujet passe après lui : \all{\textbf{An diesem Tag} \textbf{stehen} \textbf{wir} spät \textbf{auf}.} (Ce jour-là, nous nous levons tard.)
\end{aretenir}

\begin{aretenir}[Rappel 3 : lexique des fêtes]
\all{das Fest, die Feste} (la fête) ; \all{der Feiertag, die Feiertage} (le jour férié) ; \all{der Nationalfeiertag, die Nationalfeiertage} (la fête nationale) ; \all{die Parade, die Paraden} (le défilé) ; \all{die Fahne, die Fahnen} (le drapeau) ; \all{der Gottesdienst, die Gottesdienste} (l'office religieux) ; \all{Weihnachten} (Noël) ; \all{Ostern} (Pâques) ; \all{die Familie, die Familien} (la famille) ; \all{das Essen, die Essen} (le repas) ; \all{die Musik} (la musique) ; \all{feiern} (fêter, célébrer).
\end{aretenir}

\begin{attention}[Pièges à éviter]
\begin{itemize}
\item Ne dis jamais \all{*Das Fest anfängt...} : la particule doit aller à la fin : \all{Das Fest \textbf{fängt} ... \textbf{an}.}
\item N'oublie pas le \textbf{point} après l'ordinal : \all{am 20. Mai}, pas \all{*am 20 Mai}.
\item \all{einladen} est un verbe fort : \all{ich lade ein, du \textbf{lädst} ein, er \textbf{lädt} ein}.
\item Majuscule à tous les noms : \all{die Parade, das Fest, die Fahne}.
\item Faux ami : \all{die Parade} signifie « le défilé » (militaire ou carnavalesque) ; ne l'emploie pas pour « une parade » au sens de riposte ou de manœuvre.
\item Attention à la préposition : on dit \all{\textbf{an} Ostern} et \all{\textbf{an} Weihnachten} (à Pâques, à Noël), sans article.
\end{itemize}
\end{attention}

\courssub{Corrigé des questions de compréhension}

\begin{corrige}[Consigne 1 : questions sur la lettre de Lena]
\begin{enumerate}
\item \all{Der Nationalfeiertag in Deutschland ist am 3. Oktober.} (La fête nationale en Allemagne est le 3 octobre.)
\item \all{In Berlin findet ein großes Fest statt: Es gibt Musik, Essen und eine Parade.} (À Berlin a lieu une grande fête : il y a de la musique, à manger et un défilé.)
\item \all{Die deutsche Fahne hat die Farben Schwarz, Rot und Gold.} (Le drapeau allemand a les couleurs noir, rouge et or.)
\item \all{An Ostern suchen die Kinder Eier im Garten.} (À Pâques, les enfants cherchent des œufs dans le jardin.)
\end{enumerate}
\textbf{Repérage grammatical (consigne 2)} — les quatre verbes séparables de la lettre : \all{wir \textbf{stehen} ... spät \textbf{auf}} ; \all{ein großes Fest \textbf{findet} ... \textbf{statt}} ; \all{viele Leute \textbf{machen} ... \textbf{mit}} ; \all{das Fest \textbf{fängt} ... \textbf{an}} (et sa variante \all{ich \textbf{lade} dich \textbf{ein}}). Les deux dates avec ordinaux : \all{am \textbf{3.} Oktober} et \all{am \textbf{25.} Dezember}.
\end{corrige}

\courssub{Proposition de production et corrigé commenté}

\begin{exoresolu}[Modèle de lettre d'invitation]
Rédige une lettre d'invitation (8 à 10 phrases) à ta correspondante Lena pour l'inviter à la fête nationale camerounaise, en respectant toutes les contraintes de la consigne 4.
\tcblower
\textbf{Production modèle :}

\all{Yaoundé, den 10. Mai}

\all{Liebe Lena,}

\all{vielen Dank für deinen Brief. Ich lade dich zu unserem Nationalfeiertag ein! Er findet am 20. Mai in Yaoundé statt. An diesem Tag stehen wir sehr früh auf, um sechs Uhr. Dann gehen wir zum Boulevard des 20. Mai. Dort gibt es eine große Parade: Die Schüler und die Soldaten marschieren, und alle tragen die Farben unserer Fahne: Grün, Rot und Gelb. Das Fest fängt um acht Uhr an und endet am Abend mit Musik und einem großen Essen mit der Familie. Viele Leute machen beim Fest mit und singen die Nationalhymne. Wie in Deutschland ist unser Nationalfeiertag ein Feiertag: Die Schule und die Büros sind geschlossen. Ich freue mich auf dich!}

\all{Viele Grüße}\\\all{Amina}

\textbf{Commentaire des choix de langue :}
\begin{itemize}
\item \textbf{Date avec ordinal} : \all{am 20. Mai} — ordinal correct avec point, préposition \all{am} ; l'en-tête suit le modèle allemand \all{Yaoundé, den 10. Mai} (accusatif avec \all{den}, sous-entendu « Yaoundé, le 10 mai »).
\item \textbf{Verbes séparables (cinq, deux de plus que le minimum)} : \all{ich \textbf{lade} ... \textbf{ein}} ; \all{er \textbf{findet} ... \textbf{statt}} ; \all{wir \textbf{stehen} ... \textbf{auf}} ; \all{das Fest \textbf{fängt} ... \textbf{an}} ; \all{viele Leute \textbf{machen} ... \textbf{mit}}. À chaque fois, la particule est bien en fin de phrase.
\item \textbf{Lexique de la leçon (plus de cinq mots)} : \all{der Nationalfeiertag, die Parade, die Fahne, das Fest, das Essen, die Familie, der Feiertag, die Musik}.
\item \textbf{Comparaison interculturelle} : \all{Wie in Deutschland ist unser Nationalfeiertag ein Feiertag} — la structure \all{wie in Deutschland} établit le parallèle avec la lettre de Lena (\all{die Schule und die Büros sind geschlossen}), ce qui montre la compréhension du texte support.
\item \textbf{Ordre des mots} : le verbe conjugué est toujours en deuxième position (\all{An diesem Tag \textbf{stehen} wir ... \textbf{auf}} : le complément de temps en tête déclenche l'inversion sujet-verbe).
\item \textbf{Formules de lettre} : appel \all{Liebe Lena,} avec virgule et minuscule à la ligne suivante ; clôture \all{Viele Grüße} et \all{Ich freue mich auf dich!} (je me réjouis de te voir).
\end{itemize}
\end{exoresolu}

\courssub{Bilan de l'activité}

Cette activité d'intégration t'a permis de réinvestir l'ensemble des savoir-faire de la leçon dans une tâche authentique de communication écrite et orale :

\begin{itemize}
\item \textbf{Compréhension de l'écrit} : tu as lu et exploité une lettre personnelle allemande sur les fêtes (repérage d'informations, de dates et de verbes séparables).
\item \textbf{Lexique} : tu as employé activement au moins cinq mots du champ lexical des fêtes, avec leurs articles.
\item \textbf{Grammaire} : tu as écrit une date avec un ordinal correct et placé correctement les particules séparables en fin de phrase.
\item \textbf{Expression écrite} : tu as produit une lettre d'invitation structurée (en-tête, appel, corps, clôture) de 8 à 10 phrases.
\item \textbf{Expression orale et écoute} : tu as lu ta lettre à voix haute et repéré, dans les productions des camarades, la fête, la date et les verbes séparables.
\item \textbf{Interculturalité} : tu as comparé une fête camerounaise (le 20 mai) avec une fête allemande (le 3 octobre), compétence essentielle du programme de Première.
\end{itemize}

\begin{methode}[Pour aller plus loin : auto-évaluation]
Coche ce que tu maîtrises désormais :
\begin{enumerate}
\item Je sais écrire une date en allemand avec un ordinal (\all{am 20. Mai}) et la lire à voix haute (\all{am zwanzigsten Mai}).
\item Je sais conjuguer et placer un verbe séparable dans une phrase, y compris après un complément en tête de phrase.
\item Je connais au moins huit mots du lexique des fêtes avec leur article.
\item Je sais rédiger une courte lettre d'invitation en allemand (en-tête, appel, corps, clôture).
\item Je sais comparer une fête camerounaise et une fête allemande en une phrase (\all{Wie in Deutschland ...} / \all{Anders als in Deutschland ...}).
\end{enumerate}
Si une case reste vide, reprends la partie correspondante de la leçon avant l'évaluation.
\end{methode}

\newpage

\coursec{Méthodes et exercices résolus}

\courssub{Méthode 1 : Comprendre un texte sur les fêtes (lecture globale puis détaillée)}

\begin{methode}[Lire et comprendre un texte sur les fêtes religieuses ou nationales]
\begin{enumerate}
\item \textbf{Lire le titre et la source} : ils annoncent le thème (fête religieuse : \all{Weihnachten}, \all{Ostern} ; fête nationale : \all{der Nationalfeiertag}, \all{der Tag der Deutschen Einheit}).
\item \textbf{Première lecture globale} : repérer \emph{qui} (\all{die Familie, die Bürger}), \emph{quoi} (\all{feiern, der Gottesdienst, die Parade}), \emph{quand} (les dates : \all{am 3. Oktober, im Dezember}), \emph{où} (\all{in der Kirche, auf dem Platz}).
\item \textbf{Repérer les mots de la famille thématique} : \all{feiern -- das Fest -- die Feier -- festlich} ; \all{die Fahne, die Kerze, das Geschenk, der Feiertag}.
\item \textbf{Deuxième lecture détaillée} : souligner les verbes séparables (\all{stattfinden, zusammenkommen, anfangen}) et vérifier la place de la particule (fin de phrase au Präsens).
\item \textbf{Répondre aux questions} : reformuler avec les mots du texte, phrase complète, verbe conjugué en deuxième position ; ne jamais recopier un paragraphe entier.
\item \textbf{Résumé} : 3 à 5 phrases, au Präsens, avec ses propres mots ; commencer par \all{Der Text handelt von...} ou \all{In dem Text geht es um...}.
\end{enumerate}
\end{methode}

\begin{exoresolu}[Compréhension de texte : type Bac]
\textbf{Texte.} \all{In Deutschland feiern die Menschen viele Feste. Am 25. und 26. Dezember feiern die Familien Weihnachten. Sie kommen zusammen, besuchen den Gottesdienst und schenken den Kindern Geschenke. Am 3. Oktober findet der Nationalfeiertag statt: die Deutschen feiern den Tag der Deutschen Einheit. In den Städten gibt es Paraden, und überall sieht man Fahnen.}

« En Allemagne, les gens fêtent de nombreuses fêtes. Les 25 et 26 décembre, les familles fêtent Noël. Elles se réunissent, assistent à l'office religieux et offrent des cadeaux aux enfants. Le 3 octobre, la fête nationale a lieu : les Allemands célèbrent le Jour de l'Unité allemande. Dans les villes, il y a des défilés, et partout on voit des drapeaux. »

\textbf{Questions.} 1. \all{Wann feiern die Familien Weihnachten?} 2. \all{Was besuchen die Familien?} 3. \all{Was feiern die Deutschen am 3. Oktober?} 4. \all{Was sieht man überall in den Städten?}
\tcblower
\textbf{Raisonnement.} Chaque question commence par un mot interrogatif (\all{wann} = quand, \all{was} = que/quoi) ; la réponse doit être une phrase complète avec le verbe conjugué en position 2.

1. La date est dans la phrase 2 : \all{Am 25. und 26. Dezember feiern die Familien Weihnachten.} « Les familles fêtent Noël les 25 et 26 décembre. » (Après le complément de temps placé en tête, le verbe reste en position 2 et le sujet suit : c'est l'inversion.)

2. \all{Sie besuchen den Gottesdienst.} « Elles assistent à l'office religieux. » (\all{der Gottesdienst} est masculin ; à l'accusatif, l'article devient \all{den}.)

3. \all{Am 3. Oktober feiern die Deutschen den Tag der Deutschen Einheit.} « Le 3 octobre, les Allemands fêtent le Jour de l'Unité allemande. »

4. \all{Überall in den Städten sieht man Fahnen.} « Partout dans les villes, on voit des drapeaux. »

\textbf{Conclusion.} Des réponses courtes, exactes, tirées du texte mais reformulées en phrases complètes : c'est ce que le correcteur attend.
\end{exoresolu}

\courssub{Méthode 2 : Employer le vocabulaire des fêtes}

\begin{methode}[Utiliser le lexique thématique des fêtes]
\begin{enumerate}
\item \textbf{Apprendre chaque nom avec son article et son pluriel} : \all{das Fest, die Feste} ; \all{der Feiertag, die Feiertage} ; \all{die Fahne, die Fahnen} ; \all{die Parade, die Paraden} ; \all{der Gottesdienst, die Gottesdienste} ; \all{die Kerze, die Kerzen} ; \all{das Geschenk, die Geschenke}.
\item \textbf{Mémoriser les noms de fêtes} (souvent sans article) : \all{Weihnachten} (Noël), \all{Ostern} (Pâques), \all{Silvester} (la Saint-Sylvestre), \all{der Nationalfeiertag} (la fête nationale).
\item \textbf{Utiliser les verbes associés} : \all{feiern} (fêter), \all{schenken} (offrir), \all{schmücken} (décorer), \all{einladen} (inviter), \all{zusammenkommen} (se réunir).
\item \textbf{Construire des familles de mots} : \all{feiern $\rightarrow$ die Feier, das Fest, festlich} ; \all{die Flagge / die Fahne} (synonymes : le drapeau).
\item \textbf{Attention au cas} : après \all{schenken}, la personne est au datif, la chose à l'accusatif : \all{Ich schenke meinem Bruder ein Geschenk.} « J'offre un cadeau à mon frère. »
\end{enumerate}
\end{methode}

\begin{exoresolu}[Lexique : compléter un texte à trous]
\textbf{Énoncé.} Complétez avec : \all{Fahnen, Gottesdienst, Geschenke, Parade, Feiertag}.

\all{Am 3. Oktober ist in Deutschland ein nationaler (1) \dots. Viele Familien gehen in die Kirche und besuchen den (2) \dots. In Berlin gibt es eine große (3) \dots\ mit Musik. Die Kinder tragen kleine (4) \dots. Am Abend bekommen sie manchmal kleine (5) \dots.}
\tcblower
\textbf{Raisonnement.} On analyse le contexte et l'article de chaque trou.

1. \all{ein nationaler Feiertag} : « un jour férié national » (\all{der Feiertag} est masculin ; article indéfini \all{ein} au nominatif).

2. \all{den Gottesdienst} : « l'office religieux » ; après \all{besuchen}, accusatif masculin $\rightarrow$ \all{den}.

3. \all{eine große Parade} : « un grand défilé » ; \all{die Parade} est féminin, d'où l'adjectif \all{große} après \all{eine}.

4. \all{kleine Fahnen} : « de petits drapeaux » ; pluriel de \all{die Fahne}, adjectif au pluriel \all{kleine}.

5. \all{kleine Geschenke} : « de petits cadeaux » ; pluriel de \all{das Geschenk} (le radical ne change pas au pluriel).

\textbf{Texte complété.} \all{Am 3. Oktober ist in Deutschland ein nationaler Feiertag. Viele Familien gehen in die Kirche und besuchen den Gottesdienst. In Berlin gibt es eine große Parade mit Musik. Die Kinder tragen kleine Fahnen. Am Abend bekommen sie manchmal kleine Geschenke.}

\textbf{Conclusion.} Avant de choisir le mot, vérifier toujours : article + cas + pluriel.
\end{exoresolu}

\courssub{Méthode 3 : Dire et écrire les dates (nombres ordinaux)}

\begin{methode}[Former les dates et les ordinaux]
\begin{enumerate}
\item \textbf{Formation des ordinaux} : de 1 à 19, cardinal + \all{-te} (\all{zwei $\rightarrow$ der zweite, vier $\rightarrow$ der vierte}) ; à partir de 20, cardinal + \all{-ste} (\all{zwanzig $\rightarrow$ der zwanzigste}).
\item \textbf{Irréguliers à connaître par cœur} : \all{eins $\rightarrow$ der erste} ; \all{drei $\rightarrow$ der dritte} ; \all{sieben $\rightarrow$ der siebte} (ou \all{siebente}) ; \all{acht $\rightarrow$ der achte} (un seul \emph{t}).
\item \textbf{Écriture} : on écrit un point après le chiffre : \all{der 3. Oktober} (= \all{der dritte Oktober}).
\item \textbf{Question} : \all{Der wievielte ist heute?} ou \all{Welches Datum haben wir heute?} « Quelle date sommes-nous ? »
\item \textbf{Réponse} : \all{Heute ist der 3. Oktober.} ; avec \all{am} pour « le » : \all{Am 3. Oktober feiern wir den Nationalfeiertag.} (\all{am} = an + dem, datif).
\item \textbf{Date complète} : \all{am Montag, dem 3. Oktober 2024} ; l'année se lit \all{zweitausendvierundzwanzig}.
\end{enumerate}
\end{methode}

\begin{exoresolu}[Thème : traduire des phrases avec des dates]
\textbf{Énoncé.} Traduisez en allemand : 1. Noël est le 25 décembre. 2. La fête nationale allemande est le 3 octobre. 3. Mon anniversaire est le 1\ier{} mai. 4. Nous fêtons Pâques au printemps.
\tcblower
\textbf{Raisonnement et solutions.}

1. « Le 25 décembre » : ordinal $\rightarrow$ \all{der 25. Dezember} (\all{der fünfundzwanzigste}). \all{Weihnachten ist am 25. Dezember.} « Noël est le 25 décembre. » (Avec le verbe \all{sein} + une date, on utilise \all{am}.)

2. \all{Der deutsche Nationalfeiertag ist am 3. Oktober.} « La fête nationale allemande est le 3 octobre. » (Adjectif \all{deutsche} : nominatif masculin après l'article défini \all{der}.)

3. « Premier » $\rightarrow$ ordinal irrégulier \all{der erste}. \all{Mein Geburtstag ist am 1. Mai.} « Mon anniversaire est le 1\ier{} mai. »

4. \all{Wir feiern Ostern im Frühling.} « Nous fêtons Pâques au printemps. » (\all{im} = in + dem ; \all{der Frühling} est masculin.)

\textbf{Conclusion.} Trois réflexes : ordinal suivi d'un point (\all{3.}), préposition \all{am} + datif pour « le + date », ordinaux irréguliers \all{erste, dritte}.
\end{exoresolu}

\courssub{Méthode 4 : Employer les verbes à particule séparable}

\begin{methode}[Conjuguer et placer les verbes séparables]
\begin{enumerate}
\item \textbf{Reconnaître} : la particule est accentuée et détachable : \all{anfangen, aufstehen, stattfinden, einladen, zusammenkommen, mitbringen}. Verbes inséparables : préfixes \all{be-, emp-, ent-, er-, ge-, ver-, zer-} (\all{besuchen, verstehen} : jamais séparés).
\item \textbf{Au Präsens, phrase déclarative} : verbe conjugué en position 2, particule renvoyée \textbf{à la fin} : \all{Das Fest fängt um 18 Uhr an.} « La fête commence à 18 h. »
\item \textbf{À l'infinitif avec \all{zu}} : le \all{zu} s'insère entre la particule et le verbe : \all{anzufangen, stattzufinden, einzuladen}.
\item \textbf{Au Perfekt} : le \all{ge-} s'insère aussi entre particule et verbe : \all{angefangen, stattgefunden, eingeladen} : \all{Das Fest hat um 18 Uhr angefangen.} « La fête a commencé à 18 h. »
\item \textbf{Phrase interrogative} : \all{Wann findet die Parade statt?} « Quand le défilé a-t-il lieu ? »
\end{enumerate}
\end{methode}

\begin{exoresolu}[Transformation : verbes séparables au Präsens et au Perfekt]
\textbf{Énoncé.} 1. Conjuguez au Präsens : \all{Die Familie (zusammenkommen) am 25. Dezember.} 2. Mettez au Perfekt : \all{Die Parade findet am Morgen statt.} 3. Traduisez : « Nous invitons nos amis à la fête. » (\all{einladen}) 4. Traduisez : « La fête commence à 20 heures. »
\tcblower
\textbf{Raisonnement et solutions.}

1. \all{zusammenkommen} : verbe fort construit sur \all{kommen} (\all{kommt}). Particule \all{zusammen} renvoyée à la fin : \all{Die Familie kommt am 25. Dezember zusammen.} « La famille se réunit le 25 décembre. »

2. Perfekt de \all{stattfinden} : auxiliaire \all{haben} + participe \all{stattgefunden} (le \all{ge-} entre particule et verbe ; \all{finden $\rightarrow$ gefunden}). \all{Die Parade hat am Morgen stattgefunden.} « Le défilé a eu lieu le matin. »

3. \all{einladen} : verbe à voyelle changeante au Präsens (\all{er lädt ein}, a $\rightarrow$ ä). « À la fête » = \all{zu dem Fest $\rightarrow$ zum Fest}. \all{Wir laden unsere Freunde zum Fest ein.} « Nous invitons nos amis à la fête. »

4. \all{anfangen} : \all{fängt an} (a $\rightarrow$ ä à la 2\ieme{} et 3\ieme{} personne du singulier). \all{Das Fest fängt um 20 Uhr an.} « La fête commence à 20 heures. »

\textbf{Conclusion.} Schéma à retenir : [sujet] + [verbe conjugué en position 2] + \dots\ + [particule en dernier]. Au Perfekt : \all{hat} + \dots\ + [particule-ge-verbe].
\end{exoresolu}

\courssub{Méthode 5 : Produire un court paragraphe sur une fête}

\begin{methode}[Rédiger un court texte sur une fête (expression écrite guidée)]
\begin{enumerate}
\item \textbf{Plan en 4 étapes} : (a) nom de la fête et date ; (b) qui la fête et où ; (c) les activités (2 verbes séparables si possible) ; (d) une opinion (\all{Ich finde das Fest schön / interessant.}).
\item \textbf{Structures modèles} : \all{Am \dots\ feiern wir \dots} ; \all{Die Familie kommt zusammen.} ; \all{Wir besuchen den Gottesdienst.} ; \all{Das Fest fängt um \dots\ Uhr an.} ; \all{Es gibt eine Parade und viele Fahnen.}
\item \textbf{Varier les débuts de phrase} : commencer une phrase sur deux par un complément de temps (\all{Am Morgen\dots, Am Abend\dots}) avec inversion du sujet après le verbe.
\item \textbf{Relire} : majuscule à tous les noms, verbe en position 2, particule séparable en fin de phrase, Umlaute corrects.
\end{enumerate}
\end{methode}

\begin{exoresolu}[Rédaction guidée : une fête au Cameroun]
\textbf{Énoncé.} Rédigez un court paragraphe (5--6 phrases) en allemand : « La fête nationale au Cameroun » (le 20 mai : défilé, drapeaux, écoliers, famille réunie).
\tcblower
\textbf{Raisonnement.} On suit le plan : date $\rightarrow$ participants $\rightarrow$ activités $\rightarrow$ opinion. On place au moins deux verbes séparables (\all{stattfinden, zusammenkommen}) et une date avec \all{am} + ordinal.

\textbf{Production corrigée.}

\all{Am 20. Mai feiern wir in Kamerun den Nationalfeiertag. An diesem Tag findet in Yaoundé eine große Parade statt. Die Schüler und die Soldaten marschieren, und überall sieht man Fahnen. Am Abend kommt die ganze Familie zusammen. Wir essen zusammen und sehen die Parade im Fernsehen. Ich finde dieses Fest sehr schön, denn alle Menschen sind froh.}

« Le 20 mai, nous fêtons au Cameroun la fête nationale. Ce jour-là, un grand défilé a lieu à Yaoundé. Les écoliers et les soldats défilent, et partout on voit des drapeaux. Le soir, toute la famille se réunit. Nous mangeons ensemble et regardons le défilé à la télévision. Je trouve cette fête très belle, car tout le monde est joyeux. »

\textbf{Vérifications grammaticales.} \all{Am 20. Mai} : ordinal + \all{am} ; \all{findet \dots\ statt} et \all{kommt \dots\ zusammen} : particules en fin de phrase ; \all{im Fernsehen} : expression figée au datif ; \all{denn} : conjonction de coordination, le verbe reste en position 2 après elle (pas d'inversion).

\textbf{Conclusion.} Un paragraphe simple, correct et personnel vaut mieux qu'un texte long et fautif.
\end{exoresolu}

\begin{attention}[Les 5 erreurs qui coûtent des points au Bac]
\begin{enumerate}
\item \textbf{Particule séparable oubliée ou mal placée.} Écrire \all{*Das Fest anfängt um 18 Uhr} est faux : la particule va en fin de phrase : \all{Das Fest fängt um 18 Uhr an.} Au Perfekt : \all{hat angefangen} (et non \all{*hat gefangen}).
\item \textbf{Date sans \all{am} ni ordinal correct.} « Le 3 octobre » se dit \all{am 3. Oktober} (avec le point !), et non \all{*auf 3 Oktober}. « Premier » = \all{der erste}, « troisième » = \all{der dritte} (irréguliers).
\item \textbf{Majuscules oubliées sur les noms.} \all{*die fahne, *das geschenk, *der gottesdienst} : en allemand, \textbf{tous} les noms prennent la majuscule : \all{die Fahne, das Geschenk, der Gottesdienst}. Faute lourdement sanctionnée dans une copie.
\item \textbf{Faux amis et calques du français.} \all{die Parade} = le défilé ; « assister à l'office » = \all{den Gottesdienst besuchen}, pas \all{*assistieren} ; « fêter » = \all{feiern}, pas \all{*festen}.
\item \textbf{Place du verbe après un complément en tête.} \all{*Am 3. Oktober die Deutschen feiern\dots} est faux : le verbe reste en position 2 : \all{Am 3. Oktober feiern die Deutschen den Nationalfeiertag.} (inversion sujet--verbe obligatoire).
\end{enumerate}
\end{attention}

\newpage

\coursec{Exercices}

\courssub{Je vérifie mes connaissances}

\begin{exercice}[QCM : les fêtes]
Entourez la bonne réponse.
\begin{enumerate}
\item \all{Weihnachten} est :
\begin{enumerate}
\item une fête nationale \item une fête religieuse \item une fête sportive
\end{enumerate}
\item \all{Der Tag der Deutschen Einheit} est célébré :
\begin{enumerate}
\item am 1. Mai \item am 3. Oktober \item am 25. Dezember
\end{enumerate}
\item \all{Die Fahne} signifie :
\begin{enumerate}
\item le défilé \item le drapeau \item l'église
\end{enumerate}
\item \all{Der Gottesdienst} a lieu :
\begin{enumerate}
\item auf dem Marktplatz \item in der Kirche \item im Stadion
\end{enumerate}
\end{enumerate}
\end{exercice}

\begin{exercice}[Vrai ou faux ?]
Dites si les phrases suivantes sont vraies ou fausses et justifiez votre réponse par une phrase complète en allemand.
\begin{enumerate}
\item \all{Ostern ist ein religiöses Fest.}
\item \all{Der Nationalfeiertag in Kamerun ist am 3. Oktober.}
\item \all{Bei einer Parade tragen die Soldaten Fahnen.}
\item \all{Weihnachten feiert man am 31. Dezember.}
\end{enumerate}
\end{exercice}

\begin{exercice}[Vocabulaire : articles et pluriels]
Complétez le tableau avec l'article et le pluriel de chaque nom.
\begin{center}
\begin{tabular}{|l|l|l|}
\hline
\rowcolor{popblueL} Nom & Article & Pluriel \\
\hline
Fest & \ldots & \ldots \\
\hline
Fahne & \ldots & \ldots \\
\hline
Parade & \ldots & \ldots \\
\hline
Gottesdienst & \ldots & \ldots \\
\hline
Feiertag & \ldots & \ldots \\
\hline
\end{tabular}
\end{center}
\end{exercice}

\begin{exercice}[Conjugaison à compléter]
Complétez les phrases avec le verbe entre parenthèses conjugué au présent. Attention à la place de la particule !
\begin{enumerate}
\item Die Parade \ldots\ um 9 Uhr \ldots\ (anfangen).
\item Die Schüler \ldots\ am Nationalfeiertag sehr früh \ldots\ (aufstehen).
\item Der Bürgermeister \ldots\ alle Familien \ldots\ (einladen).
\item Das Fest \ldots\ auf dem Marktplatz \ldots\ (stattfinden).
\item Wir \ldots\ bei der Parade \ldots\ (mitmachen).
\end{enumerate}
\end{exercice}

\courssub{Je m'entraîne}

\begin{exercice}[Les dates en toutes lettres]
Écrivez les dates suivantes en toutes lettres, puis employez chacune d'elles dans une phrase complète avec \all{am}.

Exemple : \all{25. Dezember} $\rightarrow$ \all{der fünfundzwanzigste Dezember} — \all{Am fünfundzwanzigsten Dezember feiern wir Weihnachten.}
\begin{enumerate}
\item 1. Mai
\item 3. Oktober
\item 20. Mai
\item 25. Dezember
\item 31. Dezember
\end{enumerate}
\end{exercice}

\begin{exercice}[Séparable ou inséparable ?]
Classez les dix verbes suivants en deux colonnes : « verbes séparables » et « verbes inséparables ». Puis conjuguez deux verbes de chaque colonne à la 3\up{e} personne du singulier au présent.

\all{aufstehen — besuchen — mitmachen — verstehen — einladen — gefallen — stattfinden — erklären — anfangen — wiederholen}
\end{exercice}

\begin{exercice}[Texte à trous]
Complétez le texte suivant avec les mots de la liste : \all{Fahnen — Gottesdienst — stattfinden — Parade — Nationalfeiertag}.

\all{Am 20. Mai ist in Kamerun \ldots . Am Morgen besuchen viele Familien den \ldots\ in der Kirche. Um 10 Uhr beginnt die \ldots\ auf dem Boulevard. Die Soldaten und die Schüler tragen \ldots\ in den Nationalfarben. Die Feierlichkeiten \ldots\ in allen Städten des Landes \ldots .}
\end{exercice}

\begin{exercice}[Appariements]
Reliez chaque début de phrase à la bonne fin.
\begin{center}
\begin{tabularx}{\linewidth}{|X|X|}
\hline
\rowcolor{popblueL} Début & Fin \\
\hline
1. \all{Weihnachten feiert man} & a. \all{in der Kirche statt.} \\
\hline
2. \all{Der Gottesdienst findet} & b. \all{am 25. Dezember.} \\
\hline
3. \all{Die Schüler stehen} & c. \all{die deutsche Einheit.} \\
\hline
4. \all{Am 3. Oktober feiert Deutschland} & d. \all{am Nationalfeiertag früh auf.} \\
\hline
5. \all{Zu Ostern suchen die Kinder} & e. \all{bunte Eier im Garten.} \\
\hline
\end{tabularx}
\end{center}
\end{exercice}

\courssub{Je me prépare au Bac}

\begin{exercice}[Compréhension écrite : type Bac]
Lisez le texte suivant, puis répondez aux questions en allemand, par des phrases complètes.

\begin{textebox}[Der Tag der Deutschen Einheit]
Jedes Jahr am 3. Oktober feiert Deutschland den Tag der Deutschen Einheit. Dieser Feiertag erinnert an das Jahr 1990, als Ost- und Westdeutschland wieder ein Land wurden. Am Morgen findet in vielen Städten ein Gottesdienst statt, danach beginnt das große Fest auf den Straßen und Plätzen. Viele Familien stehen früh auf und machen bei den Feierlichkeiten mit. Es gibt Musik, Essen und Reden von Politikern. Kinder tragen kleine Fahnen in den Farben Schwarz, Rot und Gold. In einer Stadt findet jedes Jahr eine besonders große offizielle Feier statt. Am Abend laden viele Deutsche ihre Freunde und Nachbarn ein, und man spricht über die Geschichte des Landes. Für die meisten Deutschen ist dieser Tag ein Symbol der Freiheit und der Einheit.
\end{textebox}

\textbf{Questions globales :}
\begin{enumerate}
\item Worum geht es in diesem Text ? Fassen Sie den Text in zwei Sätzen zusammen.
\item Was bedeutet der 3. Oktober für die Deutschen ?
\end{enumerate}

\textbf{Questions détaillées :}
\begin{enumerate}
\setcounter{enumi}{2}
\item Wann stehen die Familien am Feiertag auf ?
\item Was machen die Kinder während des Festes ?
\item Wen laden die Deutschen am Abend ein ?
\end{enumerate}

\textbf{Vrai ou faux ?} Justifiez votre réponse par une citation du texte.
\begin{enumerate}
\setcounter{enumi}{5}
\item \all{Der Gottesdienst findet am Abend statt.}
\item \all{Die offizielle Feier findet jedes Jahr in derselben Stadt statt.}
\end{enumerate}

\textbf{Questions de langue :}
\begin{enumerate}
\setcounter{enumi}{7}
\item Trouvez dans le texte trois verbes séparables et conjuguez-les au présent à la 3\up{e} personne du singulier.
\item Écrivez la date \all{3. Oktober} en toutes lettres et employez-la dans une phrase avec \all{am}.
\end{enumerate}
\end{exercice}

\begin{exercice}[Thème et version : type Bac]
\textbf{Thème.} Traduisez en allemand :
\begin{enumerate}
\item « Le défilé a lieu le 20 mai. »
\item « Les élèves se lèvent tôt et participent à la fête. »
\item « Le 25 décembre, ma famille va à l'église. »
\end{enumerate}

\textbf{Version.} Traduisez en français le paragraphe suivant :

\begin{textebox}[Unser Nationalfeiertag]
Am 20. Mai stehen alle Schüler in Yaoundé sehr früh auf. Die Parade fängt um 9 Uhr auf dem Boulevard an. Viele Schüler machen mit und tragen die Farben des Landes. Am Nachmittag lädt unsere Schulleiterin die besten Schüler zu einer kleinen Feier ein. Am 31. Dezember feiert meine Familie dann Silvester zu Hause.
\end{textebox}
\end{exercice}

\begin{exercice}[Expression écrite : type Bac]
\textbf{Sujet :} \all{Beschreiben Sie ein Fest in Ihrem Land.}

Rédigez un court texte en allemand (environ \textbf{60 mots}) dans lequel vous présentez une fête religieuse ou nationale de votre pays. Vous indiquerez :
\begin{itemize}
\item le nom de la fête et sa date (avec \all{am} + ordinal) ;
\item les activités de la journée (utilisez au moins \textbf{deux verbes séparables} au présent) ;
\item votre opinion personnelle (\all{Ich finde dieses Fest \ldots}, \all{Mir gefällt \ldots}).
\end{itemize}

\textbf{Méthode :} faites d'abord un plan en trois parties (introduction — activités — opinion), vérifiez la place du verbe dans chaque phrase, la majuscule aux noms et la forme des ordinaux.
\end{exercice}

\newpage

\coursec{Corrigés des exercices}

\begin{corrige}[Exercice 1 — QCM : les fêtes]
\begin{enumerate}
\item \textbf{b)} \all{Weihnachten} est une fête religieuse (Noël).
\item \textbf{b)} \all{Der Tag der Deutschen Einheit} est célébré \all{am 3. Oktober}.
\item \textbf{b)} \all{Die Fahne} signifie « le drapeau ».
\item \textbf{b)} \all{Der Gottesdienst} a lieu \all{in der Kirche}.
\end{enumerate}
\end{corrige}

\begin{corrige}[Exercice 2 — Vrai ou faux ?]
\begin{enumerate}
\item \textbf{Vrai.} \all{Ostern ist ein religiöses Fest.} « Pâques est une fête religieuse. »
\item \textbf{Faux.} \all{Der Nationalfeiertag in Kamerun ist am 20. Mai.} « La fête nationale au Cameroun est le 20 mai. » (Le 3 octobre est la fête nationale allemande.)
\item \textbf{Vrai.} \all{Bei einer Parade tragen die Soldaten Fahnen.} « Lors d'un défilé, les soldats portent des drapeaux. »
\item \textbf{Faux.} \all{Weihnachten feiert man am 25. und 26. Dezember.} « On fête Noël les 25 et 26 décembre. » (Le 31 décembre, c'est \all{Silvester}, la Saint-Sylvestre.)
\end{enumerate}
\end{corrige}

\begin{corrige}[Exercice 3 — Vocabulaire : articles et pluriels]
\begin{center}
\begin{tabular}{|l|l|l|}
\hline
\rowcolor{popblueL} Nom & Article & Pluriel \\
\hline
Fest & \all{das} & \all{die Feste} \\
\hline
Fahne & \all{die} & \all{die Fahnen} \\
\hline
Parade & \all{die} & \all{die Paraden} \\
\hline
Gottesdienst & \all{der} & \all{die Gottesdienste} \\
\hline
Feiertag & \all{der} & \all{die Feiertage} \\
\hline
\end{tabular}
\end{center}
\emph{Note :} \all{das Fest} fait son pluriel en \all{-e} : \all{die Feste}. \all{der Gottesdienst} et \all{der Feiertag} prennent également \all{-e} au pluriel : \all{die Gottesdienste}, \all{die Feiertage}. Les noms féminins en \all{-e} comme \all{die Fahne} et \all{die Parade} ajoutent simplement \all{-n} : \all{die Fahnen}, \all{die Paraden}.
\end{corrige}

\begin{corrige}[Exercice 4 — Conjugaison à compléter]
\begin{enumerate}
\item Die Parade \all{fängt} um 9 Uhr \all{an}. (\all{anfangen}, séparable ; verbe fort : \all{a} $\rightarrow$ \all{ä})
\item Die Schüler \all{stehen} am Nationalfeiertag sehr früh \all{auf}. (\all{aufstehen}, séparable)
\item Der Bürgermeister \all{lädt} alle Familien \all{ein}. (\all{einladen}, séparable ; verbe fort : \all{a} $\rightarrow$ \all{ä})
\item Das Fest \all{findet} auf dem Marktplatz \all{statt}. (\all{stattfinden}, séparable)
\item Wir \all{machen} bei der Parade \all{mit}. (\all{mitmachen}, séparable)
\end{enumerate}
\emph{Rappel :} Au présent, dans une proposition principale, le verbe conjugué est en \textbf{2\textup{e} position} et la particule séparable va en \textbf{position finale} : \all{Die Parade fängt um 9 Uhr an.}
\end{corrige}

\begin{corrige}[Exercice 5 — Les dates en toutes lettres]
\begin{enumerate}
\item \all{der erste Mai} — \all{Am ersten Mai ist Tag der Arbeit.} « Le 1\textup{er} mai, c'est la fête du Travail. »
\item \all{der dritte Oktober} — \all{Am dritten Oktober feiert Deutschland die Einheit.} « Le 3 octobre, l'Allemagne fête l'Unité. »
\item \all{der zwanzigste Mai} — \all{Am zwanzigsten Mai ist Nationalfeiertag in Kamerun.} « Le 20 mai, c'est la fête nationale au Cameroun. »
\item \all{der fünfundzwanzigste Dezember} — \all{Am fünfundzwanzigsten Dezember feiern wir Weihnachten.} « Le 25 décembre, nous fêtons Noël. »
\item \all{der einunddreißigste Dezember} — \all{Am einunddreißigsten Dezember feiern wir Silvester.} « Le 31 décembre, nous fêtons la Saint-Sylvestre. »
\end{enumerate}
\emph{Attention :} Après \all{am} (= \all{an dem}, datif), l'ordinal prend la terminaison \textbf{-n} : \all{erste} $\rightarrow$ \all{ersten}, \all{dritte} $\rightarrow$ \all{dritten}, \all{zwanzigste} $\rightarrow$ \all{zwanzigsten}. Ne pas oublier les ordinaux irréguliers : \all{1. = erste}, \all{3. = dritte}, \all{7. = siebte}.
\end{corrige}

\begin{corrige}[Exercice 6 — Séparable ou inséparable ?]
\begin{center}
\begin{tabular}{|l|l|}
\hline
\rowcolor{popblueL} Verbes séparables & Verbes inséparables \\
\hline
\all{anfangen} & \all{besuchen} \\
\all{aufstehen} & \all{verstehen} \\
\all{einladen} & \all{gefallen} \\
\all{mitmachen} & \all{erklären} \\
\all{stattfinden} & \all{wiederholen} \\
\hline
\end{tabular}
\end{center}

\textbf{Conjugaison à la 3\textup{e} personne du singulier (Präsens) :}
\begin{itemize}
\item \textbf{Séparables :} \all{er fängt an}, \all{er steht auf}, \all{er lädt ein}, \all{er macht mit}, \all{es findet statt}.
\item \textbf{Inséparables :} \all{er besucht}, \all{er versteht}, \all{es gefällt}, \all{er erklärt}, \all{er wiederholt}.
\end{itemize}
\emph{Règle :} Les préfixes \all{an-}, \all{auf-}, \all{ein-}, \all{mit-}, \all{statt-} sont \textbf{séparables} (ils portent l'accent tonique : \emph{ANfangen}). Les préfixes \all{be-}, \all{ver-}, \all{ge-}, \all{er-} sont \textbf{inséparables} (non accentués). Le verbe \all{wiederholen} (« répéter ») est inséparable car l'accent tombe sur \all{-ho-} : \emph{wiederHOlen}.
\end{corrige}

\begin{corrige}[Exercice 7 — Texte à trous]
\all{Am 20. Mai ist in Kamerun \textbf{Nationalfeiertag}. Am Morgen besuchen viele Familien den \textbf{Gottesdienst} in der Kirche. Um 10 Uhr beginnt die \textbf{Parade} auf dem Boulevard. Die Soldaten und die Schüler tragen \textbf{Fahnen} in den Nationalfarben. Die Feierlichkeiten \textbf{finden} in allen Städten des Landes \textbf{statt}.}

« Le 20 mai, c'est la fête nationale au Cameroun. Le matin, beaucoup de familles assistent à l'office religieux à l'église. À 10 heures, le défilé commence sur le boulevard. Les soldats et les élèves portent des drapeaux aux couleurs nationales. Les festivités ont lieu dans toutes les villes du pays. »

\emph{Note :} Pour le dernier trou, le verbe \all{stattfinden} est séparable : la particule \all{statt} se place en fin de phrase. Le sujet \all{die Feierlichkeiten} est au pluriel, d'où \all{finden} : \all{Die Feierlichkeiten finden ... statt.}
\end{corrige}

\begin{corrige}[Exercice 8 — Appariements]
\begin{enumerate}
\item \all{1 — b} : \all{Weihnachten feiert man am 25. Dezember.} « On fête Noël le 25 décembre. »
\item \all{2 — a} : \all{Der Gottesdienst findet in der Kirche statt.} « L'office religieux a lieu à l'église. »
\item \all{3 — d} : \all{Die Schüler stehen am Nationalfeiertag früh auf.} « Les élèves se lèvent tôt le jour de la fête nationale. »
\item \all{4 — c} : \all{Am 3. Oktober feiert Deutschland die deutsche Einheit.} « Le 3 octobre, l'Allemagne fête l'unité allemande. »
\item \all{5 — e} : \all{Zu Ostern suchen die Kinder bunte Eier im Garten.} « À Pâques, les enfants cherchent des œufs colorés dans le jardin. »
\end{enumerate}
\end{corrige}

\begin{corrige}[Exercice 9 — Compréhension écrite : type Bac]
\textbf{Barème indicatif :} 10 points au total (questions globales 3 pts, questions détaillées 3 pts, vrai/faux 2 pts, questions de langue 2 pts).

\textbf{Questions globales :}
\begin{enumerate}
\item \all{Der Text handelt vom Tag der Deutschen Einheit am 3. Oktober. Er beschreibt die Feierlichkeiten: den Gottesdienst, das Fest auf den Straßen, die Parade, die Reden und die Einladungen am Abend.} « Le texte parle du Jour de l'Unité allemande, le 3 octobre. Il décrit les festivités : l'office religieux, la fête dans les rues, le défilé, les discours et les invitations du soir. »
\item \all{Der 3. Oktober ist ein Symbol der Freiheit und der Einheit für die Deutschen. Er erinnert an die Wiedervereinigung von 1990.} « Le 3 octobre est un symbole de liberté et d'unité pour les Allemands. Il rappelle la réunification de 1990. »
\end{enumerate}

\textbf{Questions détaillées :}
\begin{enumerate}
\setcounter{enumi}{2}
\item \all{Die Familien stehen am Morgen früh auf.} « Les familles se lèvent tôt le matin. » (Texte : \all{„Viele Familien stehen früh auf.“})
\item \all{Die Kinder tragen kleine Fahnen in den Farben Schwarz, Rot und Gold.} « Les enfants portent de petits drapeaux aux couleurs noir, rouge et or. »
\item \all{Am Abend laden die Deutschen ihre Freunde und Nachbarn ein.} « Le soir, les Allemands invitent leurs amis et leurs voisins. »
\end{enumerate}

\textbf{Vrai ou faux ? (avec citation justificative) :}
\begin{enumerate}
\setcounter{enumi}{5}
\item \textbf{Faux.} Citation : \all{„Am Morgen findet in vielen Städten ein Gottesdienst statt.“} L'office a lieu le \emph{matin}, et non le soir.
\item \textbf{Faux.} Citation : \all{„In einer Stadt findet jedes Jahr eine besonders große offizielle Feier statt.“} Le texte dit \all{„in einer Stadt“} (une ville, indéfinie) : la grande fête officielle n'a donc pas lieu chaque année dans la même ville.
\end{enumerate}

\textbf{Questions de langue :}
\begin{enumerate}
\setcounter{enumi}{7}
\item Verbes séparables relevés dans le texte : \all{aufstehen} (se lever), \all{stattfinden} (avoir lieu), \all{mitmachen} (participer), \all{einladen} (inviter). Conjugaison à la 3\textup{e} personne du singulier :
\begin{itemize}
\item \all{er steht auf}
\item \all{es findet statt}
\item \all{er macht mit}
\item \all{er lädt ein}
\end{itemize}
\emph{Remarque :} \all{beginnen} (« commencer ») n'est \textbf{pas} un verbe séparable ; on ne le compte donc pas.
\item \all{der dritte Oktober} — \all{Am dritten Oktober feiert Deutschland den Tag der Deutschen Einheit.} « Le 3 octobre, l'Allemagne fête le Jour de l'Unité allemande. »
\end{enumerate}
\end{corrige}

\begin{corrige}[Exercice 10 — Thème et version : type Bac]
\textbf{Thème (propositions de traduction) :}
\begin{enumerate}
\item \all{Die Parade findet am 20. Mai statt.} (ou : \all{Am 20. Mai findet die Parade statt.})
\item \all{Die Schüler stehen früh auf und machen bei der Feier mit.}
\item \all{Am 25. Dezember geht meine Familie in die Kirche.} (ou : \all{Am 25. Dezember besucht meine Familie den Gottesdienst.})
\end{enumerate}
\emph{Points de vigilance :}
\begin{itemize}
\item \textbf{Ordinaux :} \all{20. = zwanzigste}, \all{25. = fünfundzwanzigste} ; après \all{am}, terminaison \textbf{-n} au datif : \all{am zwanzigsten Mai}, \all{am fünfundzwanzigsten Dezember}.
\item \textbf{Verbes séparables :} \all{stattfinden} $\rightarrow$ \all{findet ... statt} ; \all{aufstehen} $\rightarrow$ \all{stehen ... auf} ; \all{mitmachen} $\rightarrow$ \all{machen ... mit}. La particule se place en fin de proposition.
\item \textbf{Place du verbe :} quand la phrase commence par un complément de temps (\all{Am 20. Mai ...}), le verbe reste en 2\textup{e} position et le sujet suit : \all{Am 20. Mai findet die Parade statt.}
\end{itemize}

\textbf{Version (traduction du texte allemand) :}

« Le 20 mai, tous les élèves de Yaoundé se lèvent très tôt. Le défilé commence à 9 heures sur le boulevard. Beaucoup d'élèves participent et portent les couleurs du pays. L'après-midi, notre directrice invite les meilleurs élèves à une petite fête. Le 31 décembre, ma famille fête ensuite le réveillon à la maison. »

\emph{Rappel de vocabulaire :} \all{der Umzug} = le défilé ; \all{teilnehmen / mitmachen} = participer ; \all{die Farben des Landes} = les couleurs du pays ; \all{Silvester} = la Saint-Sylvestre, le réveillon.
\end{corrige}

\begin{corrige}[Exercice 11 — Expression écrite : type Bac]
\textbf{Barème indicatif (6 pts) :} respect des consignes (nom de la fête, date, 2 verbes séparables, opinion) 3 pts ; correction grammaticale (ordre des mots, majuscules, ordinaux, conjugaison) 2 pts ; richesse lexicale et fluidité 1 pt.

\textbf{Proposition de rédaction modèle (environ 60 mots) :}

\all{In Kamerun feiern wir am 20. Mai den Nationalfeiertag. Am Morgen stehen alle Schüler früh auf, und die Parade fängt um 9 Uhr an. Viele Jugendliche machen bei dem Umzug mit und tragen die Nationalfarben Grün, Rot und Gelb. Am Nachmittag laden die Familien Freunde zum Essen ein. Ich finde dieses Fest sehr schön, weil es die Einheit des Landes zeigt. Mir gefällt die fröhliche Atmosphäre.}

« Au Cameroun, nous fêtons la fête nationale le 20 mai. Le matin, tous les élèves se lèvent tôt, et le défilé commence à 9 heures. Beaucoup de jeunes participent au cortège et portent les couleurs nationales : vert, rouge et jaune. L'après-midi, les familles invitent des amis à manger. Je trouve cette fête très belle, parce qu'elle montre l'unité du pays. L'atmosphère joyeuse me plaît. »

\textbf{Analyse du modèle :}
\begin{itemize}
\item \textbf{Introduction :} nom de la fête (\all{Nationalfeiertag}) + date (\all{am 20. Mai}).
\item \textbf{Activités (au moins 2 verbes séparables) :} \all{stehen ... auf}, \all{fängt ... an}, \all{machen ... mit}, \all{laden ... ein} (4 verbes séparables utilisés).
\item \textbf{Opinion personnelle :} \all{Ich finde dieses Fest sehr schön}, \all{Mir gefällt die fröhliche Atmosphäre} (attention : \all{gefallen} se construit avec le datif : \all{mir gefällt}).
\item \textbf{Grammaire :} ordinal correct (\all{am zwanzigsten Mai}), majuscules à tous les noms, verbe en 2\textup{e} position, particule séparable en fin de proposition.
\end{itemize}

\textbf{Points de vigilance pour l'élève :}
\begin{itemize}
\item Vérifier \textbf{l'ordre des mots} : Sujet — Verbe (2\textup{e} place) — Compléments — Particule (fin de phrase).
\item \textbf{Ordinaux :} \all{1. = erste}, \all{3. = dritte}, \all{20. = zwanzigste}, \all{25. = fünfundzwanzigste}, \all{31. = einunddreißigste}. Après \all{am} : terminaison \textbf{-n}.
\item \textbf{Majuscules :} tous les noms prennent une majuscule (\all{Feiertag}, \all{Parade}, \all{Farben}, \all{Einheit}, \all{Atmosphäre}).
\item \textbf{Verbes séparables :} ne jamais écrire \all{aufsteht} en position 2 ; on écrit \all{steht ... auf} avec la particule à la fin.
\end{itemize}
\end{corrige}

\newpage

\coursec{Fiche bilan}

\begin{popfigure}
\begin{tikzpicture}[mindmap, grow cyclic, every node/.style={concept, font=\small},
root concept/.append style={concept color=popblue, font=\bfseries},
level 1 concept/.append style={level distance=3.6cm, sibling angle=51.5, font=\scriptsize},
level 2 concept/.append style={level distance=2.1cm, font=\tiny}]
\node[root concept]{Religiöse und nationale Feste}
  child[concept color=poporange]{ node {Wortschatz : Weihnachten, Ostern, der Feiertag}
    child { node {die Fahne, die Parade} }
    child { node {der Gottesdienst} }
  }
  child[concept color=popgreen]{ node {Ordnungszahlen}
    child { node {der erste, der zweite} }
    child { node {der dritte Mai} }
  }
  child[concept color=poppurple]{ node {Datumsangaben}
    child { node {am 3. Oktober} }
    child { node {im Dezember} }
  }
  child[concept color=poppink]{ node {Trennbare Verben}
    child { node {anfangen : er fängt an} }
    child { node {stattfinden} }
  }
  child[concept color=popgold]{ node {Textverständnis}
    child { node {global, dann detailliert} }
  }
  child[concept color=popblue!60]{ node {Produktion}
    child { node {Zusammenfassung, kurzer Text} }
  }
  child[concept color=popgreen!60]{ node {Kultur}
    child { node {Tag der Deutschen Einheit} }
  };
\end{tikzpicture}
\legende{Carte mentale de la leçon : les fêtes religieuses ou nationales}
\end{popfigure}

\begin{bilan}
\textbf{1. Lexique clé (à connaître avec article et pluriel)}
\begin{itemize}
\item \all{das Fest, -e} : la fête ; \all{der Feiertag, -e} : le jour férié ; \all{der Nationalfeiertag, -e} : la fête nationale
\item \all{Weihnachten} (pl.) : Noël ; \all{Ostern} (pl.) : Pâques ; \all{der Gottesdienst, -e} : l'office religieux
\item \all{die Fahne, -n} : le drapeau ; \all{die Parade, -n} : le défilé ; \all{das Feuerwerk, -e} : le feu d'artifice
\item \all{feiern} : fêter, célébrer ; \all{stattfinden} (findet statt, fand statt, hat stattgefunden) : avoir lieu ; \all{sich versammeln} : se rassembler
\end{itemize}

\textbf{2. Grammaire à connaître par cœur}
\begin{center}
\begin{tabularx}{\linewidth}{|X|X|}
\hline
\rowcolor{popblueL} Règle & Exemple \\
\hline
Ordinaux de 1 à 19 : cardinal + \all{-te} ; exceptions : \all{erste, dritte, siebte, achte} & \all{der erste Januar, der dritte Oktober, der achte Mai} \\
\hline
À partir de 20 : cardinal + \all{-ste} & \all{der zwanzigste Mai} \\
\hline
Date au jour : \all{am} + ordinal au datif (\all{-ten / -sten}) & \all{am ersten Mai, am zwanzigsten Juni} \\
\hline
Mois : \all{im} + mois (sans article) & \all{im Dezember} \\
\hline
Verbe séparable au présent : particule en fin de phrase & \all{Die Parade fängt um 10 Uhr an.} \\
\hline
Avec un modal ou à l'infinitif : le verbe se réunit en un seul mot, à la fin & \all{Das Fest muss bald anfangen.} \\
\hline
\end{tabularx}
\end{center}

\textbf{3. Méthodes express}
\begin{itemize}
\item Compréhension : lire d'abord le titre et la première phrase, repérer \emph{qui, quand, où, quoi}, puis répondre par des phrases complètes avec le verbe en deuxième position.
\item Date en allemand : \all{der 3. Oktober} se lit \all{der dritte Oktober} ; avec préposition : \all{am dritten Oktober} (adjectif ordinal au datif).
\item Verbe séparable : repérer la particule (\all{an-, auf-, ein-, statt-}...), conjuguer le radical en deuxième position et placer la particule à la fin : \all{Das Konzert findet am Abend statt.} « Le concert a lieu le soir. »
\end{itemize}
\end{bilan}

\begin{attention}[Pièges fréquents des francophones]
\begin{itemize}
\item Ne jamais écrire \all{*der einte} : 1er se dit \all{der erste} ; 3e se dit \all{der dritte} (et non \all{*der dreite}).
\item Après \all{am}, l'ordinal prend la terminaison datif : \all{am ersten Mai}, \all{am dritten Oktober} (et non \all{*am erste Mai}).
\item Majuscule obligatoire à tous les noms : \all{die Fahne, der Gottesdienst, das Feuerwerk}.
\item Ne pas oublier la particule en fin de phrase : \all{Die Feier fängt um 18 Uhr an.} « La fête commence à 18 heures. »
\item Faux ami : \all{die Parade} signifie « le défilé », pas « la parade » au sens d'esquive.
\end{itemize}
\end{attention}

\begin{aretenir}[Checklist avant le Bac]
\begin{itemize}
\item[$\square$] Je connais le vocabulaire des fêtes avec article et pluriel.
\item[$\square$] Je sais former les ordinaux : \all{erste, zweite, dritte, vierte ... zwanzigste}.
\item[$\square$] Je sais dire une date : \all{am ersten Mai}, \all{im Juli}.
\item[$\square$] Je place correctement la particule des verbes séparables en fin de phrase.
\item[$\square$] Je réunis le verbe séparable avec un modal : \all{Wir müssen früh aufstehen.} « Nous devons nous lever tôt. »
\item[$\square$] Je mets la majuscule à tous les noms : \all{die Fahne, der Gottesdienst}.
\item[$\square$] Je sais répondre à une question de compréhension par une phrase complète.
\item[$\square$] Je sais résumer un texte en 3 à 5 phrases avec mes propres mots.
\item[$\square$] Je connais la fête nationale allemande : \all{der Tag der Deutschen Einheit, am 3. Oktober}.
\item[$\square$] Je vérifie la place du verbe conjugué en deuxième position dans la phrase.
\end{itemize}
\end{aretenir}

\findecours

\end{document}
